% The Circular flow of income — Economics Upper Sixth — Cours signé OnBuch+
% Official MINESEC syllabus. Compiler avec : tectonic ECO02-the-circular-flow-of-income.tex
\newcommand{\DOCMATIERE}{Economics}
\newcommand{\DOCNIVEAU}{Upper Sixth}
\newcommand{\DOCMODULE}{Upper Sixth — Economics}
\newcommand{\DOCLECON}{2}
\newcommand{\DOCTITRE}{The Circular flow of income}
% OnBuch+ — préambule « cours signés » ENRICHI (compilé avec tectonic / XeLaTeX)
% Système visuel « Pop » d'OnBuch+ : fond crème (jamais blanc), bordures encre
% épaisses, ombres « dures » décalées, titres Archivo Black en capitales.
% Version MATHÉMATIQUES : graphes (pgfplots/tikz), boîtes pédagogiques
% (définition, propriété, méthode, attention, le savais-tu, exercice résolu,
% exercice, TP, bilan), page de garde et signature OnBuch+ ancrée en bas.
%
% Macros attendues dans main.tex AVANT \input{preamble} :
%   \DOCMATIERE  \DOCNIVEAU  \DOCMODULE  \DOCLECON (numéro)  \DOCTITRE
\documentclass[11pt]{article}

\usepackage{fontspec}
\usepackage[english]{babel}
\usepackage[a4paper,margin=2cm,top=3.4cm,bottom=2.8cm,headheight=1.4cm,footskip=1.3cm]{geometry}
\usepackage{amsmath,amssymb,amsfonts}
\usepackage{mathtools} % \xmapsto, \coloneqq... souvent utilisés par les modèles
\usepackage{cancel} % simplifications barrées (bilans de forces)
% \abs{z} (module, valeur absolue) et \norm{u} (norme), utilisés par les modèles
\providecommand{\abs}{}\let\abs\relax\DeclarePairedDelimiter{\abs}{\lvert}{\rvert}
\providecommand{\norm}{}\let\norm\relax\DeclarePairedDelimiter{\norm}{\lVert}{\rVert}
\usepackage[table]{xcolor} % [table] : fournit aussi \rowcolors (alternance de lignes)
\usepackage{pagecolor}
\usepackage{tikz}
\usetikzlibrary{shapes.symbols,circuits.logic.US,circuits.logic.IEC,arrows.meta,positioning,calc,shapes.geometric,shapes.misc,shapes.arrows,decorations.pathmorphing,decorations.pathreplacing,decorations.markings,patterns,shadows,mindmap,trees,fit,backgrounds,matrix,intersections,angles,quotes}
\usepackage{pgfplots}
\usepackage[version=4]{mhchem} % équations nucléaires \ce{^{235}_{92}U}
\usepackage[european,siunitx]{circuitikz} % schémas électriques
\pgfplotsset{compat=1.18}
\usepgfplotslibrary{fillbetween,groupplots} % aires entre courbes (calcul intégral)
\usepackage{enumitem}
\usepackage{fancyhdr}
% [most] charge toutes les bibliothèques tcolorbox (listings, minted, raster,
% documentation...) et gonfle inutilement la mémoire TeX sur un document long
% (~300 boîtes) : on ne charge que ce qui est réellement utilisé.
\usepackage[skins,breakable]{tcolorbox}
\usepackage{graphicx}
\usepackage{colortbl}
\usepackage{booktabs}
\usepackage{tabularx}
\usepackage{diagbox} % case d'en-tête diagonale des tableaux à double entrée
% Colonnes extensibles centrée / gauche / droite (C, L, R), souvent utilisées par les modèles
\newcolumntype{C}{>{\centering\arraybackslash}X}
\newcolumntype{L}{>{\raggedright\arraybackslash}X}
\newcolumntype{R}{>{\raggedleft\arraybackslash}X}
\usepackage{array}
\usepackage{multicol}
\usepackage{siunitx}
% parse-numbers=false : accepte aussi \SI{2,5 \times 10^{-3}}{...}
\sisetup{locale=UK,output-decimal-marker={.},group-minimum-digits=4,parse-numbers=false}
\DeclareSIUnit\ohm{\ensuremath{\symup{\Omega}}} % U+2126 absent de la police
\DeclareSIUnit\division{div} % réglages d'oscilloscope (ms/div, V/div)
\DeclareSIUnit\year{yr}\DeclareSIUnit\annum{yr}\DeclareSIUnit\Ma{Ma}\DeclareSIUnit\tr{tr} % tours (vitesse de rotation, ex. \SI{1500}{\tr\per\minute})
\usepackage{qrcode}
% \degree : commande fréquemment utilisée par les modèles pour un angle ou une
% température en dehors de \SI{}{} (non fournie par les paquets chargés).
\providecommand{\degree}{^{\circ}}
% \qed (fin de démonstration), utilisé par les modèles sans amsthm
\providecommand{\qed}{\hfill$\square$}
% \barre : double barre des valeurs interdites dans un tableau de variations (tkz-tab)
\providecommand{\barre}{\ensuremath{\|}}
% environnement proof (amsthm non chargé) utilisé par les modèles
\ifdefined\proof\else\newenvironment{proof}[1][Proof]{\par\noindent\textit{#1.}\ }{\qed\par}\fi
\usepackage{lastpage}
\usepackage{hyperref}
\hypersetup{hidelinks}

% ============================================================
% Palette « Pop » OnBuch+ — mêmes hex que la classe P (plus_ui.dart)
% ============================================================
\definecolor{poporange}{HTML}{F59321}
\definecolor{popdark}{HTML}{E07A0C}
\definecolor{popcream}{HTML}{FFF6E8}
\definecolor{popink}{HTML}{1C1714}
\definecolor{poppurple}{HTML}{7A5AE0}
\definecolor{popgreen}{HTML}{2E9E6A}
\definecolor{poppink}{HTML}{E0457B}
\definecolor{popblue}{HTML}{2F7DE1}
\definecolor{popgold}{HTML}{C98A12}
\definecolor{popmuted}{HTML}{8A7F76}
\definecolor{popline}{HTML}{EADFD2}
\definecolor{popgray}{HTML}{8A7F76}
\colorlet{popgrey}{popgray}
% Alias tolérants (le modèle nomme parfois les couleurs différemment)
\colorlet{popwhite}{white}
\colorlet{popblack}{popink}
\colorlet{popred}{poppink}
\colorlet{popyellow}{popgold}
% Teintes claires pour fonds de tableaux / zones
\colorlet{popblueL}{popblue!10!white}
\colorlet{poporangeL}{poporange!14!white}
\colorlet{popgreenL}{popgreen!12!white}
\colorlet{poppurpleL}{poppurple!12!white}
\colorlet{poppinkL}{poppink!12!white}
\colorlet{popgoldL}{popgold!14!white}

\pagecolor{popcream}

% ============================================================
% Typographie
% ============================================================
\defaultfontfeatures{Ligatures=TeX}
\setmainfont{PlusJakartaSans-Medium.ttf}[Path=../../../fonts/,BoldFont=PlusJakartaSans-Bold.ttf]
% Maths en sans-serif (Fira Math) avec chiffres et lettres droites de Plus
% Jakarta, pour que les formules se fondent dans le texte.
\usepackage[math-style=ISO,bold-style=ISO]{unicode-math}
\setmathfont{FiraMath-Regular.otf}[Path=../../../fonts/]
\setmathfont{PlusJakartaSans-Medium.ttf}[Path=../../../fonts/,range={up/num,up/{Latin,latin},\mathup}]
% (icomma retiré : décimales avec point en anglais)
% Caractères mathématiques Unicode absents des polices de texte (Plus Jakarta,
% Archivo Black) : rendus par leur équivalent mathématique, en texte comme en
% formule — sinon ils s'affichent en carré vide (ex. « Arithmétique dans ℤ »).
\usepackage{newunicodechar}
\newunicodechar{ℝ}{\ensuremath{\mathbb{R}}}
\newunicodechar{ℤ}{\ensuremath{\mathbb{Z}}}
\newunicodechar{ℂ}{\ensuremath{\mathbb{C}}}
\newunicodechar{ℕ}{\ensuremath{\mathbb{N}}}
\newunicodechar{ℚ}{\ensuremath{\mathbb{Q}}}
\newunicodechar{𝔻}{\ensuremath{\mathbb{D}}}
\newunicodechar{α}{\ensuremath{\alpha}}
\newunicodechar{β}{\ensuremath{\beta}}
\newunicodechar{γ}{\ensuremath{\gamma}}
\newunicodechar{δ}{\ensuremath{\delta}}
\newunicodechar{ε}{\ensuremath{\varepsilon}}
\newunicodechar{θ}{\ensuremath{\theta}}
\newunicodechar{λ}{\ensuremath{\lambda}}
\newunicodechar{μ}{\ensuremath{\mu}}
\newunicodechar{σ}{\ensuremath{\sigma}}
\newunicodechar{τ}{\ensuremath{\tau}}
\newunicodechar{φ}{\ensuremath{\varphi}}
\newunicodechar{ψ}{\ensuremath{\psi}}
\newunicodechar{ω}{\ensuremath{\omega}}
\newunicodechar{Σ}{\ensuremath{\Sigma}}
% majuscules produites par \MakeUppercase dans les titres (α-aminés -> Α)
\newunicodechar{Α}{\ensuremath{\alpha}}
\newunicodechar{Β}{\ensuremath{\beta}}
\newunicodechar{Γ}{\ensuremath{\Gamma}}
\newunicodechar{Δ}{\ensuremath{\Delta}}
\newunicodechar{Ε}{\ensuremath{\varepsilon}}
\newunicodechar{Θ}{\ensuremath{\Theta}}
\newunicodechar{Λ}{\ensuremath{\Lambda}}
\newunicodechar{Μ}{\ensuremath{\mu}}
\newunicodechar{Τ}{\ensuremath{\tau}}
\newunicodechar{Φ}{\ensuremath{\Phi}}
\newunicodechar{Ψ}{\ensuremath{\Psi}}
\newunicodechar{Ω}{\ensuremath{\Omega}}
\newunicodechar{↦}{\ensuremath{\mapsto}}
\newunicodechar{∈}{\ensuremath{\in}}
\newunicodechar{∉}{\ensuremath{\notin}}
\newunicodechar{∧}{\ensuremath{\wedge}}
\newunicodechar{⇔}{\ensuremath{\Leftrightarrow}}
\newunicodechar{⟺}{\ensuremath{\Longleftrightarrow}}
\newunicodechar{⇒}{\ensuremath{\Rightarrow}}
\newunicodechar{⟹}{\ensuremath{\Longrightarrow}}
\newunicodechar{∘}{\ensuremath{\circ}}
\newunicodechar{∪}{\ensuremath{\cup}}
\newunicodechar{∩}{\ensuremath{\cap}}
\newunicodechar{≡}{\ensuremath{\equiv}}
\newunicodechar{⊂}{\ensuremath{\subset}}
\newunicodechar{⊥}{\ensuremath{\perp}}
\newunicodechar{∃}{\ensuremath{\exists}}
\newunicodechar{∀}{\ensuremath{\forall}}
\newunicodechar{∛}{\ensuremath{\sqrt[3]{\,}}}
\newunicodechar{∜}{\ensuremath{\sqrt[4]{\,}}}
\newunicodechar{⃗}{\ensuremath{{}^{\to}}}
\newunicodechar{ⁿ}{\ifmmode^{n}\else\textsuperscript{\ensuremath{n}}\fi}
\newunicodechar{ˣ}{\ifmmode^{x}\else\textsuperscript{\ensuremath{x}}\fi}
\newunicodechar{ᵇ}{\ifmmode^{b}\else\textsuperscript{\ensuremath{b}}\fi}
\newunicodechar{ʳ}{\ifmmode^{r}\else\textsuperscript{\ensuremath{r}}\fi}
\newunicodechar{ᵘ}{\ifmmode^{u}\else\textsuperscript{\ensuremath{u}}\fi}
\newunicodechar{ʸ}{\ifmmode^{y}\else\textsuperscript{\ensuremath{y}}\fi}
\newunicodechar{ᵃ}{\ifmmode^{a}\else\textsuperscript{\ensuremath{a}}\fi}
\newunicodechar{ᵉ}{\ifmmode^{e}\else\textsuperscript{\ensuremath{e}}\fi}
\newunicodechar{ᵏ}{\ifmmode^{k}\else\textsuperscript{\ensuremath{k}}\fi}
\newunicodechar{ᵖ}{\ifmmode^{p}\else\textsuperscript{\ensuremath{p}}\fi}
\newunicodechar{ᴺ}{\ifmmode^{N}\else\textsuperscript{\ensuremath{N}}\fi}
\newunicodechar{ᶜ}{\ifmmode^{c}\else\textsuperscript{\ensuremath{c}}\fi}
\newunicodechar{ʰ}{\ifmmode^{h}\else\textsuperscript{\ensuremath{h}}\fi}
\newunicodechar{ᵠ}{\ifmmode^{q}\else\textsuperscript{\ensuremath{q}}\fi}
\newunicodechar{ₙ}{\ifmmode_{n}\else\textsubscript{\ensuremath{n}}\fi}
\newunicodechar{ₐ}{\ifmmode_{a}\else\textsubscript{\ensuremath{a}}\fi}
\newunicodechar{ᵢ}{\ifmmode_{i}\else\textsubscript{\ensuremath{i}}\fi}
\newunicodechar{ᵣ}{\ifmmode_{r}\else\textsubscript{\ensuremath{r}}\fi}
\newunicodechar{ₖ}{\ifmmode_{k}\else\textsubscript{\ensuremath{k}}\fi}
\newunicodechar{ᵦ}{\ifmmode_{\beta}\else\textsubscript{\ensuremath{\beta}}\fi}
\newunicodechar{ₓ}{\ifmmode_{x}\else\textsubscript{\ensuremath{x}}\fi}
\newfontfamily\popdisplay{ArchivoBlack-Regular.ttf}[Path=../../../fonts/]
\newfontfamily\poplogo{SpaceGrotesk-Bold.ttf}[Path=../../../fonts/]
\newfontfamily\popsemi{PlusJakartaSans-SemiBold.ttf}[Path=../../../fonts/]
\newfontfamily\popxbold{PlusJakartaSans-ExtraBold.ttf}[Path=../../../fonts/]
\newfontfamily\popxboldls{PlusJakartaSans-ExtraBold.ttf}[Path=../../../fonts/,LetterSpace=3.0]

\setlist[enumerate]{leftmargin=1.7em,itemsep=5pt,topsep=4pt}
\setlist[itemize]{leftmargin=1.7em,itemsep=4pt,topsep=4pt,label={\color{poporange}\textbullet}}
\setlength{\parindent}{0pt}
\setlength{\parskip}{5pt}
\renewcommand{\arraystretch}{1.25}

% pgfplots : style Pop par défaut
\pgfplotsset{
  every axis/.append style={axis line style={popink,line width=1pt},tick style={popink},
    grid style={popline,line width=0.5pt},label style={font=\small},tick label style={font=\footnotesize}},
  popcurve/.style={poporange,line width=1.8pt,no markers,smooth},
}
% Même style disponible dans un \draw TikZ simple (hors pgfplots)
\tikzset{popcurve/.style={poporange,line width=1.8pt}}
\tikzset{popgrid/.style={popline,very thin}}
\colorlet{popcurve}{poporange} % le nom du style est parfois utilisé comme couleur

% ============================================================
% Pill / squiggle
% ============================================================
\newcommand{\poppill}[3]{%
  \tikz[baseline=-0.55ex]\node[draw=popink,line width=1.2pt,rounded corners=2.6pt,%
    fill=#2,inner xsep=6.5pt,inner ysep=3pt]{%
    \popxboldls\fontsize{7}{7}\selectfont\color{#3}\MakeUppercase{#1}};%
}
\newcommand{\popsquiggle}[1]{%
  \tikz[baseline=-0.4ex]{
    \draw[#1,line width=2.6pt,line cap=round]
      plot [smooth] coordinates {(0,0.09) (0.34,0.01) (0.68,0.21) (1.02,0.01) (1.36,0.10)};
  }%
}
% Mot-clé surligné (marqueur orange sous le texte)
\newcommand{\cle}[1]{{\popxbold\color{popdark}#1}}

% ============================================================
% En-tête / pied
% ============================================================
\pagestyle{fancy}
\fancyhf{}
\renewcommand{\headrulewidth}{0pt}
\renewcommand{\footrulewidth}{0pt}
\lhead{\raisebox{-0.15ex}{\poplogo\fontsize{13}{13}\selectfont\color{popink}OnBuch\color{poporange}+}}
\rhead{\poppill{\DOCMATIERE\ \textbullet\ \DOCNIVEAU\ \textbullet\ Chapter \DOCLECON}{white}{popink}}
\lfoot{\popxbold\fontsize{7.6}{7.6}\selectfont\color{popmuted}\MakeUppercase{Signed course OnBuch+}\ \textbullet\ {\popsemi\DOCTITRE}}
\rfoot{\tikz[baseline=-0.6ex]\node[rounded corners=8.5pt,fill=popink,minimum height=17pt,minimum width=17pt,inner xsep=5pt,inner sep=0]{\popxbold\fontsize{8}{8}\selectfont\color{popcream}\thepage\,/\,\pageref*{LastPage}};}

% ============================================================
% Titres
% ============================================================
\newcommand{\chaptitre}[2]{%
  \begin{tcolorbox}[enhanced,colback=poporange,colframe=popink,boxrule=2.6pt,arc=11pt,
      drop shadow=popink,left=15pt,right=15pt,top=12pt,bottom=11pt,before skip=3pt,after skip=18pt]
    \poppill{#2}{popink}{popcream}\par\vspace{9pt}
    {\raggedright\hyphenpenalty=10000\exhyphenpenalty=10000\popdisplay\fontsize{22}{24}\selectfont\color{popcream}\MakeUppercase{#1}\par}
  \end{tcolorbox}
}
% Section numérotée : pastille orange avec le numéro + titre Archivo
\newcounter{popsec}
\newcommand{\coursec}[1]{%
  \stepcounter{popsec}\setcounter{popsub}{0}%
  \par\vspace{16pt}\noindent
  \begin{minipage}{\linewidth}
  \tikz[baseline=-0.7ex]\node[draw=popink,line width=1.6pt,fill=poporange,rounded corners=4pt,minimum size=22pt,inner sep=0]{\popdisplay\fontsize{12}{12}\selectfont\color{popcream}\thepopsec};%
  \hspace{8pt}{\popdisplay\fontsize{15}{17}\selectfont\color{popink}\MakeUppercase{#1}}\par
  \vspace{4pt}\noindent{\color{poporange}\rule{36pt}{3.6pt}}
  \end{minipage}\par\nopagebreak\vspace{8pt}
}
\newcounter{popsub}
\newcommand{\courssub}[1]{%
  \stepcounter{popsub}%
  \par\vspace{9pt}\noindent{\popxbold\fontsize{11.5}{13}\selectfont\color{popink}{\color{poporange}\thepopsec.\thepopsub}\hspace{6pt}#1}\par\nopagebreak\vspace{4pt}
}

% ============================================================
% Boîtes pédagogiques — même recette : fond blanc, bordure épaisse colorée,
% ombre dure encre, badge plein. Titre optionnel [#1].
% ============================================================
\tcbset{popbase/.style={breakable,enhanced,colback=white,boxrule=2.3pt,arc=9pt,
  shadow={3pt}{-3pt}{0pt}{popink},left=14pt,right=14pt,top=11pt,bottom=11pt,before skip=11pt,after skip=15pt,
  fontupper=\popsemi\fontsize{10.6}{14.5}\selectfont\color{popink},
  fontlower=\popsemi\fontsize{10.6}{14.5}\selectfont\color{popink}}}
% Coche dessinée (le glyphe ✓ n'existe pas dans les polices du document)
\newcommand{\popcheck}[1]{\tikz[baseline=-0.5ex]\draw[#1,line width=1.6pt,line cap=round,line join=round](-0.13,0.0)--(-0.04,-0.09)--(0.13,0.1);}
\providecommand{\coche}{\popcheck{popgreen}}
\providecommand{\resultat}[1]{\par\smallskip\noindent\fcolorbox{popgreen}{popgreenL}{\parbox{\dimexpr\linewidth-2\fboxsep-2\fboxrule\relax}{\textbf{Result.} #1}}\par\smallskip}
\newcommand{\popboxhead}[3]{\poppill{#1}{#2}{white}\ifx\relax#3\relax\else\hspace{6pt}{\popxbold\fontsize{10.6}{12}\selectfont\color{popink}#3}\fi\par\vspace{7pt}}

\newtcolorbox{aretenir}[1][]{popbase,colframe=popgreen,before upper={\popboxhead{Key points}{popgreen}{#1}}}
\newtcolorbox{exemplebox}[1][]{popbase,colframe=poporange,before upper={\popboxhead{Exemple}{poporange}{#1}}}
\newtcolorbox{definition}[1][]{popbase,colframe=popblue,before upper={\popboxhead{Definition}{popblue}{#1}}}
\newtcolorbox{propriete}[1][]{popbase,colframe=poppurple,before upper={\popboxhead{Property}{poppurple}{#1}}}
\newtcolorbox{methode}[1][]{popbase,colframe=popgold,before upper={\popboxhead{Method}{popgold}{#1}}}
\newtcolorbox{attention}[1][]{popbase,colframe=poppink,before upper={\popboxhead{Warning}{poppink}{#1}}}
\newtcolorbox{experience}[1][]{popbase,colframe=popblue,colback=popblueL,before upper={\popboxhead{Experiment}{popblue}{#1}}}
% « Le savais-tu ? » : carte violette pleine (contexte, culture, Cameroun)
\newtcolorbox{savaistu}[1][]{popbase,colback=poppurple,colframe=popink,
  fontupper=\popsemi\fontsize{10.4}{14.2}\selectfont\color{white},
  before upper={\poppill{Did you know?}{white}{poppurple}\ifx\relax#1\relax\else\hspace{6pt}{\popxbold\color{white}#1}\fi\par\vspace{7pt}}}
% Exercice résolu : énoncé en haut, \tcblower puis la solution sur fond crème
\newcounter{popexo}\newcounter{popexr}
\newtcolorbox{exoresolu}[1][]{popbase,colframe=poporange,colbacklower=popcream,
  segmentation style={popink,line width=1.2pt,dashed},
  before upper={\stepcounter{popexr}\popboxhead{Worked example \thepopexr}{poporange}{#1}},
  before lower={\poppill{Solution}{popink}{popcream}\par\vspace{6pt}}}
% Exercice (non résolu en place) — la correction est dans la partie Corrigés
\newtcolorbox{exercice}[1][]{popbase,colframe=popink,
  before upper={\stepcounter{popexo}\popboxhead{Exercise \thepopexo}{popink}{#1}}}
% Corrigé
\newtcolorbox{corrige}[1][]{popbase,colframe=popgreen,colback=popgreenL,
  before upper={\popboxhead{Solution}{popgreen}{#1}}}
% Objectifs / prérequis / bilan
\newtcolorbox{objectifs}{popbase,colframe=popink,colback=poporangeL,
  before upper={\popboxhead{Chapter objectives}{popink}{}}}
\newtcolorbox{prerequis}{popbase,colframe=popink,colback=popblueL,
  before upper={\popboxhead{Prerequisites}{popblue}{}}}
\newtcolorbox{bilan}{popbase,colframe=popink,colback=white,boxrule=2.8pt,
  before upper={\popboxhead{Fiche bilan}{poporange}{L'essentiel à connaître pour l'examen}}}
% Figure encadrée avec légende
\newtcolorbox{popfigure}[1][]{enhanced,breakable=false,colback=white,colframe=popline,boxrule=1.4pt,arc=8pt,
  left=10pt,right=10pt,top=10pt,bottom=8pt,before skip=10pt,after skip=12pt,halign=center,
  lower separated=false,halign lower=center,
  fontlower=\popsemi\fontsize{9}{11.5}\selectfont\color{popmuted},
  #1}
\newcommand{\legende}[1]{\par\vspace{6pt}{\popsemi\fontsize{9}{11.5}\selectfont\color{popmuted}#1}}

% ============================================================
% Signature OnBuch+
% ============================================================
\newcommand{\poplogoinline}[1]{{\poplogo\fontsize{#1}{#1}\selectfont\color{popink}OnBuch\color{poporange}+}}
\newcommand{\signatureonbuch}{%
  \begin{tcolorbox}[enhanced,colback=popink,colframe=popink,boxrule=2.2pt,arc=10pt,
      drop shadow=poporange,left=14pt,right=14pt,top=10pt,bottom=10pt,before skip=0pt,after skip=0pt]
    \tikz[baseline=-0.7ex]\node[circle,fill=popgreen,draw=popcream,line width=1.3pt,minimum size=22pt,inner sep=0]{\popcheck{white}};%
    \hspace{9pt}\begin{minipage}[c]{0.62\linewidth}
      {\popxbold\fontsize{12}{13}\selectfont\color{popcream}Signed\ \textbullet\ {\poplogo OnBuch\color{poporange}+}}\par\vspace{2pt}
      {\raggedright\popsemi\fontsize{8}{10.5}\selectfont\color{popline}\MakeUppercase{OnBuch+ teaching team}\\\MakeUppercase{Follows the official MINESEC syllabus}\par}
    \end{minipage}\hfill
    {\poplogo\fontsize{22}{22}\selectfont\color{popcream}OnBuch\color{poporange}+}
  \end{tcolorbox}%
}

% ============================================================
% Page de garde
% #1 titre · #2 matière · #3 niveau · #4 module · #5 n° leçon · #6 durée ·
% #7 description / objectifs (texte ou itemize)
% ============================================================
\newcommand{\pagegarde}[7]{%
  \thispagestyle{empty}
  \newgeometry{a4paper,margin=1.7cm,top=1.6cm,bottom=1.4cm}%
  \begin{tikzpicture}[remember picture,overlay]
    \fill[poporange!18] ([xshift=-3.2cm,yshift=-3.6cm]current page.north east) circle (4.6cm);
    \fill[poppurple!14] ([xshift=2.4cm,yshift=7.6cm]current page.south west) circle (3.8cm);
  \end{tikzpicture}%
  {\poplogo\fontsize{30}{30}\selectfont\color{popink}OnBuch\color{poporange}+}\hfill
  \raisebox{0.6ex}{\poppill{Signed course \textbullet\ Premium}{popink}{popcream}}\par
  \vspace{1pt}\popsquiggle{poppurple}\par
  \vspace{8pt}
  \begin{tcolorbox}[enhanced,colback=poporange,colframe=popink,boxrule=2.8pt,arc=13pt,
      drop shadow=popink,left=18pt,right=18pt,top=12pt,bottom=13pt,after skip=10pt]
    \poppill{#2 \textbullet\ #3}{popink}{popcream}\hspace{5pt}\poppill{#4}{white}{popink}\par\vspace{8pt}
    {\popdisplay\fontsize{14}{14}\selectfont\color{popink}CHAPTER #5}\par\vspace{4pt}
    {\raggedright\hyphenpenalty=10000\exhyphenpenalty=10000\popdisplay\fontsize{25}{27}\selectfont\color{popcream}\MakeUppercase{#1}\par}
    \vspace{8pt}
    {\popxbold\fontsize{9.5}{9.5}\selectfont\color{popink}\MakeUppercase{Suggested time: #6}}
  \end{tcolorbox}
  \begin{tcolorbox}[enhanced,colback=white,colframe=popink,boxrule=2.2pt,arc=10pt,
      drop shadow=popink,left=16pt,right=16pt,top=10pt,bottom=10pt,after skip=8pt]
    \poppill{In this chapter}{poppurple}{white}\par\vspace{5pt}
    \popsemi\fontsize{9.3}{12.6}\selectfont\color{popink}#7
  \end{tcolorbox}
  \noindent\begin{minipage}[t]{0.487\linewidth}
    \begin{tcolorbox}[enhanced,colback=popgreen,colframe=popink,boxrule=2.2pt,arc=10pt,
        drop shadow=popink,left=11pt,right=11pt,top=10pt,bottom=10pt,height=3.1cm,valign=center]
      \tcbox[colback=white,colframe=popink,boxrule=1.3pt,arc=5pt,boxsep=0pt,nobeforeafter,
        left=3pt,right=3pt,top=3pt,bottom=3pt,valign=center]{\qrcode[height=1.9cm]{https://play.google.com/store/apps/details?id=com.luvvix.onbuch&hl=en}}%
      \hspace{8pt}\begin{minipage}[c]{0.5\linewidth}
        {\popdisplay\fontsize{11}{12.5}\selectfont\color{white}\MakeUppercase{Download OnBuch}}\par\vspace{4pt}
        {\raggedright\popsemi\fontsize{8.4}{11}\selectfont\color{white}Free \textbullet\ Google Play\\AI tutor, past papers, results\par}
      \end{minipage}
    \end{tcolorbox}
  \end{minipage}\hfill
  \begin{minipage}[t]{0.487\linewidth}
    \begin{tcolorbox}[enhanced,colback=poppurple,colframe=popink,boxrule=2.2pt,arc=10pt,
        drop shadow=popink,left=11pt,right=11pt,top=10pt,bottom=10pt,height=3.1cm,valign=center]
      \tikz[baseline=-0.7ex]\node[circle,fill=white,draw=popink,line width=1.3pt,minimum size=1.6cm,inner sep=0]{\popdisplay\fontsize{24}{24}\selectfont\color{poppurple}+};%
      \hspace{8pt}\begin{minipage}[c]{0.6\linewidth}
        {\popdisplay\fontsize{11}{12.5}\selectfont\color{white}\MakeUppercase{Go OnBuch+}}\par\vspace{4pt}
        {\raggedright\popsemi\fontsize{8.4}{11}\selectfont\color{white}All signed courses, unlimited AI tutor, PDF sheets — OnBuch+ tab in the app\par}
      \end{minipage}
    \end{tcolorbox}
  \end{minipage}
  \vspace{8pt}
  \popcontenu
  \vfill
  \signatureonbuch
  \restoregeometry
  \newpage
}

% Rangée « Ce cours contient » (page de garde)
\newcommand{\poptile}[3]{% #1 couleur #2 gros texte #3 légende
  \begin{tcolorbox}[enhanced,colback=#1,colframe=popink,boxrule=2pt,arc=9pt,shadow={2.5pt}{-2.5pt}{0pt}{popink},
      left=6pt,right=6pt,top=9pt,bottom=9pt,halign=center,valign=center,height=1.75cm,width=\linewidth,nobeforeafter]
    {\popdisplay\fontsize{12.5}{13}\selectfont\color{white}\MakeUppercase{#2}}\par\vspace{3pt}
    {\centering\popsemi\fontsize{7.8}{9.5}\selectfont\color{white}#3\par}
  \end{tcolorbox}}
\newcommand{\popcontenu}{%
  \par\noindent{\popxboldls\fontsize{7.5}{7.5}\selectfont\color{popmuted}THIS SIGNED COURSE CONTAINS}\par\vspace{7pt}
  \noindent\begin{minipage}[t]{0.235\linewidth}\poptile{popblue}{Course}{complete, illustrated, step by step}\end{minipage}\hfill
  \begin{minipage}[t]{0.235\linewidth}\poptile{poporange}{Methods}{and worked examples}\end{minipage}\hfill
  \begin{minipage}[t]{0.235\linewidth}\poptile{poppink}{Exercises}{exam style + solutions}\end{minipage}\hfill
  \begin{minipage}[t]{0.235\linewidth}\poptile{popgreen}{Summary sheet}{the essentials to remember}\end{minipage}\par}

% Sommaire
\newtcolorbox{sommaire}{enhanced,colback=white,colframe=popink,boxrule=2.2pt,arc=10pt,
  drop shadow=popink,left=18pt,right=18pt,top=13pt,bottom=13pt,before skip=6pt,after skip=20pt,
  before upper={\poppill{Contents}{poppurple}{white}\par\vspace{9pt}\popsemi\fontsize{10.6}{14.5}\selectfont\color{popink}}}

% Signature de fin de cours — poussée tout en bas de la dernière page
\newcommand{\findecours}{%
  \par\vfill\nopagebreak
  \begin{center}\popsquiggle{poporange}\end{center}\vspace{4pt}
  \signatureonbuch
}
\AtBeginDocument{\def\llbracket{\mathopen{[\mkern-3.5mu[}}\def\rrbracket{\mathclose{]\mkern-3.5mu]}}} % intervalles d'entiers (glyphe ⟦ absent de la police)
% Glyphes absents de Fira Math : redéfinis à partir de glyphes disponibles
\AtBeginDocument{%
  \def\setminus{\mathbin{\backslash}}\let\smallsetminus\setminus
  \def\vdots{\mathord{\vbox{\baselineskip3.2pt\lineskiplimit0pt\kern4pt\hbox{.}\hbox{.}\hbox{.}}}}%
  \def\ddots{\mathinner{\mkern1mu\raise6.5pt\hbox{.}\mkern2mu\raise3.6pt\hbox{.}\mkern2mu\raise0.7pt\hbox{.}\mkern1mu}}%
  \def\triangle{\mathord{\scalebox{0.9}{$\symup{\Delta}$}}}%
  \def\nabla{\mathord{\raisebox{\depth}{\scalebox{1}[-1]{$\symup{\Delta}$}}}}%
  \def\checkmark{\popcheck{popdark}}%
  \def\star{\mathbin{\ast}}\def\bigstar{\mathord{\ast}}%
  \def\diamond{\mathbin{\scalebox{0.6}{$\Diamond$}}}%
  \def\bigcup{\mathop{\mathchoice{\raisebox{-0.3em}{\scalebox{1.7}{$\cup$}}}{\scalebox{1.25}{$\cup$}}{\cup}{\cup}}\displaylimits}%
  \def\bigcap{\mathop{\mathchoice{\raisebox{-0.3em}{\scalebox{1.7}{$\cap$}}}{\scalebox{1.25}{$\cap$}}{\cap}{\cap}}\displaylimits}%
  \def\lesseqqgtr{\mathrel{\lessgtr}}\def\bigoplus{\mathop{\oplus}\displaylimits}%
}
\providecommand{\ssi}{if and only if} % abréviation fréquente
\AtBeginDocument{\providecommand{\wideparen}{\overparen}} % arc de cercle

% Fonctions usuelles que les modèles utilisent sans que amsmath les fournisse
\providecommand{\arsinh}{\operatorname{arsinh}}\providecommand{\arcosh}{\operatorname{arcosh}}
\providecommand{\artanh}{\operatorname{artanh}}\providecommand{\arsech}{\operatorname{arsech}}
\providecommand{\arcsch}{\operatorname{arcsch}}\providecommand{\arcoth}{\operatorname{arcoth}}
\providecommand{\arcsinh}{\operatorname{arsinh}}\providecommand{\arccosh}{\operatorname{arcosh}}
\providecommand{\arctanh}{\operatorname{artanh}}\providecommand{\sech}{\operatorname{sech}}
\providecommand{\csch}{\operatorname{csch}}\providecommand{\arcsec}{\operatorname{arcsec}}
\providecommand{\arccsc}{\operatorname{arccsc}}\providecommand{\arccot}{\operatorname{arccot}}
\providecommand{\sgn}{\operatorname{sgn}}\providecommand{\lcm}{\operatorname{lcm}}
\providecommand{\Var}{\operatorname{Var}}\providecommand{\Cov}{\operatorname{Cov}}
\providecommand{\permil}{\ensuremath{{}^{0}\!/_{00}}}\providecommand{\textperthousand}{\ensuremath{{}^{0}\!/_{00}}}

%% FIGFIX-BEGIN v3 — correctifs globaux des figures (docs/onbuch-generation/tools/figfix)
\usepackage{adjustbox}\usepackage{etoolbox}
% toute figure TikZ/pgfplots plus large que la ligne est réduite à la largeur disponible
\BeforeBeginEnvironment{tikzpicture}{\begin{adjustbox}{max width=\linewidth}}
\AfterEndEnvironment{tikzpicture}{\end{adjustbox}}
% légendes pgfplots lisibles par-dessus les courbes
\pgfplotsset{every axis/.append style={legend style={fill=white,fill opacity=0.92,text opacity=1,draw=black!15,inner sep=3pt}}}
% étiquettes d'arêtes (syntaxe edge["..."]) avec fond blanc
\tikzset{every edge quotes/.append style={fill=white,inner sep=1.2pt}}
% bulles de cartes mentales : texte à la ligne dans la bulle, bulle un peu plus grande
\tikzset{every concept/.append style={text width=2.0cm,align=center,minimum size=2.3cm,inner sep=2pt}}
%% FIGFIX-END
\begin{document}

\pagegarde{The Circular flow of income}{Economics}{Upper Sixth}{Upper Sixth — Economics}{2}{24 periods}{This chapter explains how income circulates between households and firms, and how the key macroeconomic variables — consumption, saving and investment — interact to determine the equilibrium level of national income. Students master the multiplier and accelerator mechanisms and learn to compute APC, MPC, APS, MPS, NPV and equilibrium income, core skills for the GCE Advanced Level Economics examination.
\vspace{4pt}
\begin{multicols}{2}\small
\begin{itemize}
\item By the end of this chapter I can identify and define the main macroeconomic variables (income, consumption, saving, investment, aggregate demand).
\item By the end of this chapter I can draw and explain the circular flow of income in a two-sector and a multi-sector economy, identifying injections and withdrawals.
\item By the end of this chapter I can explain the consumption function, the Keynesian psychological law of consumption, and other consumption theories.
\item By the end of this chapter I can calculate and interpret APC, MPC, APS and MPS from data and verify that APC+APS=1 and MPC+MPS=1.
\item By the end of this chapter I can explain the determinants of saving and investment and appraise investment projects using payback, ARR and net present value.
\item By the end of this chapter I can determine the equilibrium level of income using both the aggregate demand approach and the injections/withdrawals approach.
\end{itemize}\end{multicols}}

\begin{sommaire}
\begin{multicols}{2}
\begin{enumerate}
\item Macroeconomic variables and the circular flow of income
\item Consumption and the consumption function
\item Saving and the saving function
\item Practical work: calculating APC, MPC, APS and MPS
\item Investment and investment appraisal
\item Practical work: net present value (NPV)
\item Equilibrium level of national income
\item Practical work: equilibrium and full-employment income
\item The multiplier
\item The accelerator and the multiplier–accelerator interaction
\item Integration activity
\item Methods and worked examples
\item Exercises
\item Solutions to the exercises
\item Summary sheet
\end{enumerate}
\end{multicols}
\end{sommaire}

\begin{objectifs}
\begin{itemize}
\item By the end of this chapter I can identify and define the main macroeconomic variables (income, consumption, saving, investment, aggregate demand).
\item By the end of this chapter I can draw and explain the circular flow of income in a two-sector and a multi-sector economy, identifying injections and withdrawals.
\item By the end of this chapter I can explain the consumption function, the Keynesian psychological law of consumption, and other consumption theories.
\item By the end of this chapter I can calculate and interpret APC, MPC, APS and MPS from data and verify that APC+APS=1 and MPC+MPS=1.
\item By the end of this chapter I can explain the determinants of saving and investment and appraise investment projects using payback, ARR and net present value.
\item By the end of this chapter I can determine the equilibrium level of income using both the aggregate demand approach and the injections/withdrawals approach.
\item By the end of this chapter I can distinguish equilibrium income from full-employment income and illustrate inflationary and deflationary gaps.
\item By the end of this chapter I can calculate the multiplier and explain how an initial injection leads to a multiplied change in national income.
\item By the end of this chapter I can explain and compute the accelerator and analyse the interaction between the multiplier and the accelerator.
\end{itemize}
\end{objectifs}

\begin{prerequis}
\begin{itemize}
\item Basic economic agents and their roles (households, firms, government) from Lower Sixth
\item National income concepts: GDP, GNP and their measurement (Lower Sixth)
\item Elementary algebra: solving simple linear equations and working with percentages and ratios
\item Demand and supply analysis and the idea of market equilibrium (Lower Sixth)
\item Factors of production and the rewards of factors (wages, rent, interest, profit)
\end{itemize}
\end{prerequis}

\newpage

\coursec{Macroeconomic variables and the circular flow of income}

Every morning in Douala, a trader at Marché Mboppi sells plantains to a civil servant, who was paid by the State, which collected taxes from the trader's own business. Money seems to travel in circles: one person's spending becomes another person's income, which becomes someone else's spending again. But why does this circular flow sometimes speed up, creating jobs and prosperity, and sometimes slow down, as during the 2016–2018 economic slowdown linked to falling oil prices? Why does the government inject billions into road projects hoping the whole economy will grow by even more? This chapter unveils the engine of the macroeconomy: the circular flow of income and the forces — consumption, saving, investment, the multiplier and the accelerator — that determine a nation's level of income and employment.

\courssub{Macroeconomics versus microeconomics}

Economics is traditionally split into two complementary branches. \textbf{Microeconomics} studies the behaviour of individual decision‑making units — households, firms, markets — and how they interact to determine prices and quantities of specific goods (e.g. the price of cocoa per kilogramme in Kumba). \textbf{Macroeconomics}, by contrast, looks at the economy as a whole: it aggregates millions of individual decisions into economy‑wide totals such as national income, total employment, the general price level, and the balance of payments.

\begin{definition}[Macroeconomic variable]
A \cle{macroeconomic variable} is an aggregate magnitude that describes the overall performance of an economy. It is measured at the national level (or for a monetary union like CEMAC) over a given period, usually a year or a quarter. The main macroeconomic variables are:
\begin{itemize}
    \item \textbf{National Income (Y)} – the total value of final goods and services produced in the domestic territory during a period (GDP at market prices or factor cost).
    \item \textbf{Consumption (C)} – total spending by households on goods and services.
    \item \textbf{Saving (S)} – the part of disposable income not spent on consumption.
    \item \textbf{Investment (I)} – spending on capital goods (machinery, buildings, infrastructure) and changes in inventories.
    \item \textbf{Government Spending (G)} – current and capital expenditure by all levels of government.
    \item \textbf{Taxes (T)} – compulsory transfers from households and firms to the government.
    \item \textbf{Exports (X)} – sales of domestically produced goods and services to the rest of the world.
    \item \textbf{Imports (M)} – purchases of foreign‑produced goods and services by domestic residents.
\end{itemize}
\end{definition}

\begin{propriete}[Stock vs. flow variables]
Macroeconomic variables are either \textbf{stocks} or \textbf{flows}:
\begin{itemize}
    \item A \cle{stock variable} is measured at a specific point in time (e.g. national wealth, money supply, public debt on 31 December).
    \item A \cle{flow variable} is measured per unit of time (e.g. national income per year, investment per quarter, exports per month).
\end{itemize}
National income (Y), consumption (C), saving (S), investment (I), government spending (G), taxes (T), exports (X) and imports (M) are all \textbf{flow variables}.
\end{propriete}

\begin{propriete}[Nominal vs. real variables]
\begin{itemize}
    \item \textbf{Nominal variables} are measured in current money values (e.g. nominal GDP in CFA francs of the current year).
    \item \textbf{Real variables} are adjusted for inflation, measured in constant prices of a base year (e.g. real GDP in 2015 CFA francs).
\end{itemize}
In the circular flow model, we usually work with real variables to analyse physical output and welfare, but the monetary flows are nominal. The GCE exam often asks you to distinguish between them.
\end{propriete}

\begin{attention}[Micro vs. macro: do not mix the levels]
A common mistake is to apply microeconomic reasoning to macroeconomic problems. For example, what is rational for one household (saving more during a recession) can be harmful for the whole economy if everyone does it simultaneously (\emph{paradox of thrift}). In the GCE exam, always check whether the question asks for a micro or a macro analysis.
\end{attention}

\courssub{The two‑sector circular flow model}

The simplest representation of the macroeconomy involves only two sectors: \textbf{households} and \textbf{firms}. We make the following \textbf{assumptions}:
\begin{enumerate}
    \item There is no government (no taxes, no public spending).
    \item The economy is closed (no exports, no imports).
    \item All output is either consumed or invested; there is no inventory accumulation beyond planned investment.
    \item Households own all factors of production (land, labour, capital) and supply them to firms.
    \item Firms produce all goods and services and sell them to households.
\end{enumerate}

In this model, two distinct but simultaneous flows circulate in opposite directions:

\begin{itemize}
    \item \textbf{Real flow (physical flow):} Factors of production flow from households to firms; finished goods and services flow from firms to households.
    \item \textbf{Money flow (monetary flow):} Factor payments (wages, rent, interest, profit) flow from firms to households; consumption expenditure flows from households to firms.
\end{itemize}

\begin{popfigure}
\begin{tikzpicture}[>=Stealth, node distance=3.5cm, every node/.style={font=\small}]
    % Sectors
    \node[draw, rounded corners, fill=popblueL, thick, minimum width=2.8cm, minimum height=1.8cm, align=center] (HH) at (0,0){\textbf{Households}\\\textit{Own factors of production}\\\textit{Demand goods \& services}};
    \node[draw, rounded corners, fill=popgreenL, thick, minimum width=2.8cm, minimum height=1.8cm, align=center] (F) at (7,0){\textbf{Firms}\\\textit{Hire factors}\\\textit{Supply goods \& services}};

    % Real flows (solid arrows, blue)
    \draw[->, thick, popblue] (HH.east) -- node[above, midway] {\textbf{Factor services}} (F.west);
    \draw[->, thick, popblue] (F.west) -- node[below, midway] {\textbf{Goods \& services}} (HH.east);

    % Money flows (dashed arrows, orange)
    \draw[<-, dashed, thick, poporange] ([yshift=0.5cm]HH.east) -- node[above, midway] {\textbf{Factor payments (Y)}} ([yshift=0.5cm]F.west);
    \draw[<-, dashed, thick, poporange] ([yshift=-0.5cm]F.west) -- node[below, midway] {\textbf{Consumption spending (C)}} ([yshift=-0.5cm]HH.east);

    % Labels for flow types
    \node[align=center, popblue] at (3.5, 1.5) {\textbf{Real Flow}};
    \node[align=center, poporange] at (3.5, -1.5) {\textbf{Money Flow}};
\end{tikzpicture}
\legende{Two‑sector circular flow: real flows (solid blue) and money flows (dashed orange) move in opposite directions.}
\end{popfigure}

\begin{attention}[Real flow vs. money flow — exam trap]
Never confuse the two. The \textbf{real flow} consists of physical goods and services and factor services; it is measured in quantities (tonnes of cocoa, hours of labour). The \textbf{money flow} consists of the payments that finance those real flows; it is measured in CFA francs. In the diagram, draw them as separate arrows (often solid for real, dashed for monetary) and label each arrow clearly: \textit{Factor services}, \textit{Goods \& services}, \textit{Factor payments (Y)}, \textit{Consumption (C)}.
\end{attention}

\courssub{Withdrawals (leakages) and injections}

In the two‑sector model, the circular flow is closed: every CFA franc earned flows back as spending. In reality, households do not spend all their income, and firms, the government, and foreigners also spend. This creates \textbf{withdrawals (leakages)} — money leaving the circular flow — and \textbf{injections} — money entering the flow from outside the household‑firm circuit.

\begin{definition}[Withdrawal (Leakage) and Injection]
A \cle{withdrawal (leakage)} is any use of income that does not return directly as spending on domestic output. The three withdrawals are:
\begin{itemize}
    \item \textbf{Saving (S)} – income not spent on consumption.
    \item \textbf{Taxation (T)} – income taken by the government.
    \item \textbf{Imports (M)} – spending on foreign‑produced goods.
\end{itemize}
An \cle{injection} is any spending on domestic output that does not originate from current household income. The three injections are:
\begin{itemize}
    \item \textbf{Investment (I)} – firms' spending on capital goods.
    \item \textbf{Government spending (G)} – public expenditure on goods and services.
    \item \textbf{Exports (X)} – foreigners' spending on domestic output.
\end{itemize}
\end{definition}

When we add the financial sector (banks, capital markets), the government sector, and the foreign sector, we obtain the \textbf{five‑sector circular flow model}.

\begin{popfigure}
\begin{tikzpicture}[>=Stealth, node distance=2.8cm, every node/.style={font=\small, align=center}]
    % Sectors
    \node[draw, rounded corners, fill=popblueL, thick, minimum width=2.5cm, minimum height=1.6cm] (HH) at (0,0) {\textbf{Households}};
    \node[draw, rounded corners, fill=popgreenL, thick, minimum width=2.5cm, minimum height=1.6cm] (F) at (5,0) {\textbf{Firms}};
    \node[draw, rounded corners, fill=poppink, thick, minimum width=2.5cm, minimum height=1.6cm] (FIN) at (2.5,2.5) {\textbf{Financial Sector}};
    \node[draw, rounded corners, fill=popgold, thick, minimum width=2.5cm, minimum height=1.6cm] (GOV) at (-2.5,0) {\textbf{Government}};
    \node[draw, rounded corners, fill=poppurple, thick, minimum width=2.5cm, minimum height=1.6cm] (FOR) at (7.5,0) {\textbf{Foreign Sector}};

    % Core flows between HH and Firms
    \draw[->, thick, popblue] (HH.east) -- node[above] {Factor services} (F.west);
    \draw[->, thick, popblue] (F.west) -- node[below] {Goods \& services} (HH.east);
    \draw[<-, dashed, thick, poporange] ([yshift=0.4cm]HH.east) -- node[above] {Factor payments (Y)} ([yshift=0.4cm]F.west);
    \draw[<-, dashed, thick, poporange] ([yshift=-0.4cm]F.west) -- node[below] {Consumption (C)} ([yshift=-0.4cm]HH.east);

    % Financial sector: Saving and Investment
    \draw[->, thick, popdark] (HH) -- node[left, pos=0.4] {Saving (S)} (FIN);
    \draw[->, thick, popdark] (FIN) -- node[right, pos=0.6] {Investment (I)} (F);

    % Government: Taxes and Govt spending
    \draw[->, thick, popdark] (HH) -- node[above, pos=0.4] {Taxes (T)} (GOV);
    \draw[->, thick, popdark] (F) -- node[above, pos=0.6] {Taxes (T)} (GOV);
    \draw[->, thick, popdark] (GOV) -- node[above] {Govt spending (G)} (F);
    \draw[->, thick, popdark] (GOV) -- node[below] {Transfers/Subsidies} (HH);

    % Foreign sector: Imports and Exports
    \draw[->, thick, popdark] (HH) -- node[above] {Imports (M)} (FOR);
    \draw[->, thick, popdark] (F) -- node[below] {Imports (M)} (FOR);
    \draw[->, thick, popdark] (FOR) -- node[above] {Exports (X)} (F);

    % Legend
    \node[align=left, anchor=north west] at (-4.5,-2.2) {
        \textbf{Leakages (Withdrawals):} S, T, M \\
        \textbf{Injections:} I, G, X
    };
\end{tikzpicture}
\legende{Five‑sector circular flow showing leakages (S, T, M) and injections (I, G, X). Solid arrows = real flows; dashed arrows = money flows (only core shown for clarity).}
\end{popfigure}

\courssub{Equilibrium in the circular flow}

The circular flow is in \textbf{equilibrium} when the total amount of money flowing out (withdrawals) equals the total amount flowing in (injections). If withdrawals exceed injections, the flow of income shrinks (recession); if injections exceed withdrawals, the flow expands (boom).

\begin{aretenir}[Equilibrium condition]
In the five‑sector model, the equilibrium condition for the level of national income is:
\[
\boxed{I + G + X = S + T + M}
\]
\textbf{Total Injections = Total Withdrawals}
This identity holds \emph{ex post} (after the fact) in national accounts. \emph{Ex ante} (planned), the economy moves toward this equality through changes in national income (Y).
\end{aretenir}

\begin{exemplebox}[Identifying leakages and injections in Cameroon]
Consider the following transactions in a given year:
\begin{enumerate}
    \item A household in Yaoundé deposits 500\,000 FCFA in a savings account at Afriland First Bank.
    \item The government awards a contract to build a new bridge over the Wouri River.
    \item A cocoa cooperative in the South‑West exports 10\,000 tonnes of cocoa beans to the Netherlands.
    \item A firm in Douala imports machinery from China.
    \item A civil servant pays income tax (IRPP).
\end{enumerate}
Classify each as a leakage or an injection.
\begin{center}
\begin{tabularx}{\linewidth}{|l|X|X|}\hline
\textbf{Transaction} & \textbf{Classification} & \textbf{Reason} \\ \hline
1. Saving deposit & Leakage (S) & Income not spent on domestic output \\ \hline
2. Bridge construction & Injection (G) & Government spending on domestic output \\ \hline
3. Cocoa exports & Injection (X) & Foreign spending on domestic output \\ \hline
4. Machinery imports & Leakage (M) & Domestic spending on foreign output \\ \hline
5. Income tax payment & Leakage (T) & Income transferred to government \\ \hline
\end{tabularx}
\end{center}
\end{exemplebox}

\begin{methode}[Drawing the circular flow diagram for the GCE exam]
\begin{enumerate}
    \item Draw the two core sectors (Households, Firms) as boxes or circles.
    \item Add the Financial Sector, Government, Foreign Sector around them.
    \item Draw \textbf{solid arrows} for real flows (factor services, goods/services, exports/imports).
    \item Draw \textbf{dashed arrows} for money flows (factor payments, consumption, S, I, T, G, X, M).
    \item Label \textbf{every arrow} with its standard symbol (S, I, T, G, X, M, C, Y).
    \item Indicate the direction: leakages flow \textbf{out} of the core circle; injections flow \textbf{in}.
    \item Write the equilibrium condition $I+G+X = S+T+M$ below the diagram.
\end{enumerate}
\end{methode}

\begin{savaistu}[The CFA franc zone and the circular flow]
Because Cameroon uses the CFA franc (pegged to the euro), the \textbf{foreign sector} has a special feature: the \emph{operations account} at the French Treasury acts as a buffer. When export earnings (X) exceed import payments (M), the surplus accumulates in the operations account — an injection that does not immediately expand domestic money supply unless the central bank (BEAC) decides to inject liquidity. This institutional detail matters for monetary policy but does not change the fundamental identity $I+G+X = S+T+M$.
\end{savaistu}

\begin{exoresolu}[Equilibrium income in a simple model]
Suppose a closed economy with government (no foreign sector) has the following behavioural equations (all values in billions of FCFA):
\begin{align*}
C &= 200 + 0.75 Y_d \quad &\text{(Consumption function)}\\
Y_d &= Y - T \quad &\text{(Disposable income)}\\
T &= 100 \quad &\text{(Lump‑sum tax)}\\
I &= 150 \quad &\text{(Autonomous investment)}\\
G &= 200 \quad &\text{(Government spending)}
\end{align*}
Find the equilibrium level of national income $Y^*$ using the injection–withdrawal approach.

\textbf{Solution}
\tcblower
\textbf{Step 1: Identify injections and withdrawals.}
Injections: $J = I + G = 150 + 200 = 350$.
Withdrawals: $W = S + T$. We need the saving function.
$S = Y_d - C = (Y - T) - (200 + 0.75(Y - T)) = 0.25(Y - T) - 200$.
With $T = 100$: $S = 0.25(Y - 100) - 200 = 0.25Y - 25 - 200 = 0.25Y - 225$.
Thus $W = S + T = (0.25Y - 225) + 100 = 0.25Y - 125$.

\textbf{Step 2: Set $J = W$.}
$350 = 0.25Y - 125 \implies 0.25Y = 475 \implies Y^* = 1900$ billion FCFA.

\textbf{Step 3: Verify with $Y = C + I + G$.}
$Y_d = 1900 - 100 = 1800$.
$C = 200 + 0.75(1800) = 1550$.
$C + I + G = 1550 + 150 + 200 = 1900 = Y^*$. \quad $\checkmark$
\end{exoresolu}

\begin{exercice}[Practice: Circular flow and equilibrium]
An open economy with government has the following data (billions FCFA):
$C = 500 + 0.8(Y - T)$, $T = 200$, $I = 300$, $G = 400$, $X = 250$, $M = 100 + 0.1Y$.
\begin{enumerate}
    \item Write the expressions for total injections $J$ and total withdrawals $W$.
    \item Calculate the equilibrium national income $Y^*$.
    \item At $Y^*$, compute the values of $C$, $S$, $M$ and verify $J = W$.
\end{enumerate}
\end{exercice}

\begin{corrige}[Exercise 1]
\begin{enumerate}
    \item $J = I + G + X = 300 + 400 + 250 = 950$.
    $S = Y_d - C = (Y - 200) - [500 + 0.8(Y - 200)] = 0.2(Y - 200) - 500 = 0.2Y - 40 - 500 = 0.2Y - 540$.
    $W = S + T + M = (0.2Y - 540) + 200 + (100 + 0.1Y) = 0.3Y - 240$.
    \item Set $J = W$: $950 = 0.3Y - 240 \implies 0.3Y = 1190 \implies Y^* = 3966.\overline{6} \approx 3966.7$ billion FCFA.
    \item $Y_d = 3966.7 - 200 = 3766.7$.
    $C = 500 + 0.8(3766.7) = 3513.3$.
    $S = 3766.7 - 3513.3 = 253.4$ (or $0.2(3966.7) - 540 = 253.3$).
    $M = 100 + 0.1(3966.7) = 496.7$.
    $W = 253.3 + 200 + 496.7 = 950 = J$. \quad $\checkmark$
\end{enumerate}
\end{corrige}

\coursec{Consumption and the consumption function}

In the previous section we identified the main macroeconomic variables and saw how they circulate between households and firms in a simple two-sector economy. We now zoom in on the most important component of aggregate demand: \textbf{consumption expenditure}. Understanding how households decide to spend their income is the first step to explaining how the equilibrium level of national income is determined.

\courssub{What is consumption?}

\begin{definition}[Consumption expenditure]
Consumption expenditure (\cle{C}) is the total spending by households on goods and services that directly satisfy their current wants and needs. In national income accounting, it is the value of all final goods and services purchased by households during a given period.
\end{definition}

Economists classify consumption goods according to their durability:

\begin{itemize}
    \item \textbf{Durable goods:} Goods that last for a relatively long time (typically more than three years) and yield utility over several periods. Examples: cars, refrigerators, furniture, smartphones. In Cameroon, a second-hand vehicle imported from Europe or a locally manufactured bed frame are durable goods.
    \item \textbf{Non-durable goods:} Goods that are used up in a single act of consumption or last for a short period (less than three years). Examples: food (plantains, rice, fish), fuel, clothing, newspapers.
    \item \textbf{Services:} Intangible acts of consumption that are produced and consumed simultaneously. Examples: haircuts, transport (taxi, bus fares), education fees, medical consultations, mobile phone airtime.
\end{itemize}

\begin{savaistu}[Consumption in Cameroon]
In Cameroon, as in many developing economies, a large share of household consumption goes to food and non-alcoholic beverages (often above 40\% of total household expenditure). The informal sector plays a huge role: buying roasted plantain (\emph{bobolo}) from a street vendor or taking a \emph{bensikin} (motorcycle taxi) are both consumption expenditures that are part of the circular flow, even if they are not always captured perfectly in official statistics.
\end{savaistu}

\courssub{The consumption function}

The relationship between consumption and its main determinant — disposable income — is called the \textbf{consumption function}. In its simplest linear form, which Keynes used in his \emph{General Theory}, it is written as:

\begin{propriete}[The linear consumption function]
\[
C = a + bY_d
\]
where:
\begin{itemize}
    \item $C$ = planned consumption expenditure,
    \item $Y_d$ = disposable income (income after direct taxes and transfer payments),
    \item $a$ = \textbf{autonomous consumption} (the intercept),
    \item $b$ = \textbf{marginal propensity to consume} (MPC), the slope of the function, with $0 < b < 1$.
\end{itemize}
\end{propriete}

\begin{definition}[Autonomous consumption ($a$)]
Autonomous consumption is the level of consumption that occurs even when disposable income is zero. It represents the minimum subsistence spending that households must undertake (food, shelter) and is financed by drawing down past savings (dissaving), borrowing, or receiving transfers.
\end{definition}

\begin{definition}[Induced consumption ($bY_d$)]
Induced consumption is the part of consumption that varies directly with disposable income. As income rises, households spend a fraction $b$ of each additional franc on consumption.
\end{definition}

\begin{aretenir}[Key features of the linear consumption function]
\begin{itemize}
    \item The function is a straight line with a positive intercept $a > 0$ on the vertical axis.
    \item The slope $b$ (MPC) is positive but less than 1: $0 < b < 1$.
    \item As income rises, consumption rises but \emph{by less than the rise in income}. This is the core of Keynes' psychological law.
\end{itemize}
\end{aretenir}

\courssub{Keynes' fundamental psychological law of consumption}

\begin{propriete}[Keynes' psychological law of consumption]
\begin{quote}
\textit{``Men are disposed, as a rule and on the average, to increase their consumption as their income increases, but not by as much as the increase in their income.''}
\end{quote}
— John Maynard Keynes, \emph{The General Theory of Employment, Interest and Money} (1936), Chapter 8.
\end{propriete}

This law has three propositions:
\begin{enumerate}
    \item When aggregate income increases, consumption expenditure also increases, but by a smaller amount.
    \item The increment of income is divided between consumption and saving in a certain ratio (MPC and MPS).
    \item Both consumption and saving increase when income increases (since $0 < MPC < 1$, then $0 < MPS < 1$).
\end{enumerate}

\begin{attention}[Common misinterpretation]
The law does \emph{not} say that consumption falls when income rises. It says the \emph{proportion} of income consumed (APC) falls as income rises, while the \emph{amount} consumed still rises.
\end{attention}

\courssub{Average and marginal propensities to consume}

Two key ratios measure the relationship between consumption and income:

\begin{propriete}[Average Propensity to Consume (APC)]
\[
APC = \frac{C}{Y_d}
\]
It is the fraction of total disposable income that is spent on consumption. In the linear function $C = a + bY_d$,
\[
APC = \frac{a}{Y_d} + b
\]
Since $a > 0$, APC declines as $Y_d$ rises (the term $a/Y_d$ gets smaller).
\end{propriete}

\begin{propriete}[Marginal Propensity to Consume (MPC)]
\[
MPC = \frac{\Delta C}{\Delta Y_d}
\]
It is the fraction of an \emph{additional} unit of disposable income that is spent on consumption. For the linear function, MPC is constant and equal to the slope $b$.
\end{propriete}

\begin{attention}[APC vs MPC — a classic exam trap]
\begin{itemize}
    \item \textbf{APC} is a ratio of \textbf{levels} ($C/Y$). It changes along the consumption function.
    \item \textbf{MPC} is a ratio of \textbf{changes} ($\Delta C/\Delta Y$). For a linear function it is constant (the slope).
    \item At low incomes, $APC > MPC$ (because of autonomous consumption). At very high incomes, $APC \to MPC$.
    \item Never confuse $C/Y$ with $\Delta C/\Delta Y$.
\end{itemize}
\end{attention}

\courssub{Graphical representation: the consumption function and the 45° line}

The consumption function is usually drawn with disposable income ($Y_d$) on the horizontal axis and consumption ($C$) on the vertical axis. The 45° line (where $C = Y_d$) is a crucial reference: it shows all points where consumption exactly equals income.

\begin{popfigure}
\begin{tikzpicture}
\begin{axis}[
    width=0.9\linewidth,
    height=0.6\linewidth,
    axis lines=middle,
    xlabel={Disposable Income ($Y_d$)},
    ylabel={Consumption ($C$)},
    xmin=0, xmax=1000,
    ymin=0, ymax=1000,
    xtick={200,400,600,800,1000},
    ytick={200,400,600,800,1000},
    xticklabels={200,400,600,800,1000},
    yticklabels={200,400,600,800,1000},
    domain=0:1000,
    samples=2,
    clip=false,
    legend style={at={(0.95,0.05)}, anchor=south east, font=\small},
]
% 45 degree line
\addplot [popline, thick, dashed] coordinates {(0,0) (1000,1000)};
\addlegendentry{$45^\circ$ line ($C=Y_d$)}

% Consumption function C = 200 + 0.6Y
\addplot [popblue, very thick, domain=0:1000] {200 + 0.6*x};
\addlegendentry{$C = 200 + 0.6Y_d$}

% Autonomous consumption label
\node[popblue, anchor=west] at (0,200) {$a=200$};

% Dissaving and saving zones
\node[poporange, align=center, font=\small] at (300,500) {Dissaving\\$(C > Y_d)$};
\node[popgreen, align=center, font=\small] at (800,300) {Saving\\$(C < Y_d)$};

% Break-even point
\fill[popdark] (500,500) circle (2pt);
\node[above left, font=\small] at (500,500) {Break-even ($Y_0$)};

% Vertical arrow for MPC illustration
\draw[poppurple, <->, thick] (400,440) -- (600,440) node[midway, below, font=\small] {$\Delta Y_d = 200$};
\draw[poppurple, <->, thick] (600,440) -- (600,560) node[midway, right, font=\small] {$\Delta C = 120$};
\node[poppurple, font=\small] at (650,500) {MPC $= 120/200 = 0.6$};
\end{axis}
\end{tikzpicture}
\legende{The linear consumption function $C = a + bY_d$ and the 45° line. Below the break-even point $Y_0$, consumption exceeds income (dissaving). Above $Y_0$, income exceeds consumption (saving). The slope of the consumption function is the MPC.}
\end{popfigure}

\begin{exemplebox}[Reading the graph]
Consider the function $C = 200 + 0.6Y_d$ plotted above.
\begin{itemize}
    \item At $Y_d = 0$, $C = 200$ (autonomous consumption). The household dissaves 200.
    \item At $Y_d = 500$, $C = 200 + 0.6(500) = 500$. This is the \textbf{break-even income} ($Y_0$) where $C = Y_d$ and saving is zero.
    \item At $Y_d = 1000$, $C = 800$. Saving $S = Y_d - C = 200$.
    \item The slope (MPC) is 0.6: for every 1000 CFA francs extra income, consumption rises by 600 CFA francs.
\end{itemize}
\end{exemplebox}

\courssub{Determinants of consumption (beyond current income)}

While the simple Keynesian function makes consumption depend solely on current disposable income, in reality several other factors shift the consumption function (change $a$ or $b$):

\begin{center}
\begin{tabularx}{\linewidth}{|l|X|}\hline
\textbf{Determinant} & \textbf{Effect on Consumption} \\ \hline
\textbf{Wealth (net worth)} & Higher wealth (houses, land, shares, savings) increases autonomous consumption ($a \uparrow$). A household in Douala owning a plot of land feels richer and spends more. \\ \hline
\textbf{Interest rates} & Higher rates encourage saving (substitution effect) and raise the cost of borrowing for durables, reducing $a$ and possibly $b$. Lower rates have the opposite effect. \\ \hline
\textbf{Expectations} & Optimism about future income or prices raises current consumption (shift up). Fear of job loss or inflation reduces it. \\ \hline
\textbf{Distribution of income} & Redistribution from rich (low MPC) to poor (high MPC) raises aggregate consumption because the poor spend a larger fraction of each franc. \\ \hline
\textbf{Availability of credit} & Easier access to loans (e.g., microfinance in Cameroon) allows households to smooth consumption and finance durables, raising $a$. \\ \hline
\textbf{Fiscal policy} & Direct taxes reduce $Y_d$, shifting the function down along the axis. Transfer payments (pensions, subsidies) shift it up. \\ \hline
\end{tabularx}
\end{center}

\courssub{Further studies: other consumption hypotheses}

The simple Keynesian function ($C = a + bY_d$) is the \textbf{absolute income hypothesis}: current consumption depends on current absolute income. Three major alternatives were developed to explain empirical puzzles (e.g., why the long-run APC is constant while the short-run APC falls).

\begin{definition}[Relative Income Hypothesis (Duesenberry, 1949)]
\begin{itemize}
    \item Consumption depends on \textbf{relative income} — an individual's income compared to others (demonstration effect / ``keeping up with the Joneses'').
    \item \textbf{Ratchet effect:} Once a household reaches a certain consumption level, it is hard to cut back even if income falls (past peak income matters).
    \item Implies the long-run consumption function is proportional (passes through origin), while short-run functions are flatter.
\end{itemize}
\end{definition}

\begin{definition}[Permanent Income Hypothesis (Friedman, 1957)]
\begin{itemize}
    \item Income is split into \textbf{permanent income} ($Y_p$ — expected long-run average) and \textbf{transitory income} ($Y_t$ — windfalls, temporary shocks).
    \item Consumption depends mainly on permanent income: $C = k Y_p$ (where $k$ depends on interest rate, wealth, tastes).
    \item Transitory income is mostly saved. This explains why short-run APC $<$ long-run APC.
\end{itemize}
\end{definition}

\begin{definition}[Life-Cycle Hypothesis (Modigliani \& Brumberg, 1954)]
\begin{itemize}
    \item Individuals plan consumption over their entire lifetime to smooth it.
    \item Young: borrow/dissave (income low, consumption high for education, housing).
    \item Middle age: high income, save for retirement (high saving).
    \item Old: dissave (run down assets).
    \item Aggregate consumption depends on wealth and age structure of population, not just current income.
\end{itemize}
\end{definition}

\begin{aretenir}[Comparison of hypotheses]
\begin{center}
\begin{tabularx}{\linewidth}{|l|X|X|X|}\hline
\textbf{Hypothesis} & \textbf{Key Variable} & \textbf{Short-run APC vs MPC} & \textbf{Long-run APC} \\ \hline
Absolute Income (Keynes) & Current $Y_d$ & $APC > MPC$ (APC falls as $Y$ rises) & Falls as economy grows \\ \hline
Relative Income (Duesenberry) & Relative $Y$, Past peak $Y$ & $APC > MPC$ & Constant (proportional) \\ \hline
Permanent Income (Friedman) & Permanent $Y_p$ & $APC > MPC$ (transitory income saved) & Constant ($= k$) \\ \hline
Life-Cycle (Modigliani) & Lifetime resources, Wealth, Age & Varies by age group & Constant for given wealth/age structure \\ \hline
\end{tabularx}
\end{center}
\end{aretenir}

\courssub{Worked example: calculating APC and MPC from data}

\begin{exoresolu}[APC and MPC from a hypothetical Cameroonian household survey]
\textbf{Statement:} The table below shows hypothetical annual disposable income and consumption data for a household in Yaoundé over five years (in millions of CFA francs).

\begin{center}
\begin{tabular}{|c|c|c|c|c|c|}\hline
Year & 1 & 2 & 3 & 4 & 5 \\ \hline
$Y_d$ (million CFAF) & 2.0 & 3.0 & 4.0 & 5.0 & 6.0 \\ \hline
$C$ (million CFAF) & 2.2 & 2.9 & 3.6 & 4.3 & 5.0 \\ \hline
\end{tabular}
\end{center}

\textbf{Tasks:}
\begin{enumerate}
    \item Calculate APC and MPC for each year (where possible).
    \item Determine the linear consumption function $C = a + bY_d$ that best fits this data.
    \item Identify the break-even income level.
    \item Comment on the trend of APC.
\end{enumerate}

\tcblower
\textbf{Detailed solution:}

\textbf{1. Calculations:}

\begin{center}
\begin{tabular}{|c|c|c|c|c|c|c|}\hline
Year & $Y_d$ & $C$ & $APC = C/Y_d$ & $\Delta Y_d$ & $\Delta C$ & $MPC = \Delta C/\Delta Y_d$ \\ \hline
1 & 2.0 & 2.2 & 1.10 & — & — & — \\ \hline
2 & 3.0 & 2.9 & 0.967 & 1.0 & 0.7 & 0.70 \\ \hline
3 & 4.0 & 3.6 & 0.900 & 1.0 & 0.7 & 0.70 \\ \hline
4 & 5.0 & 4.3 & 0.860 & 1.0 & 0.7 & 0.70 \\ \hline
5 & 6.0 & 5.0 & 0.833 & 1.0 & 0.7 & 0.70 \\ \hline
\end{tabular}
\end{center}

\textbf{2. Consumption function:}
The MPC is constant at $b = 0.7$.
Using Year 1: $2.2 = a + 0.7(2.0) \Rightarrow a = 2.2 - 1.4 = 0.8$.
Check with Year 5: $0.8 + 0.7(6.0) = 0.8 + 4.2 = 5.0$. Correct.
\[
\boxed{C = 0.8 + 0.7Y_d}
\]

\textbf{3. Break-even income ($Y_0$):}
Set $C = Y_d$:
\[
Y_0 = 0.8 + 0.7Y_0 \Rightarrow 0.3Y_0 = 0.8 \Rightarrow Y_0 = \frac{0.8}{0.3} \approx 2.67 \text{ million CFAF}.
\]
Below 2.67 million, the household dissaves; above it, the household saves.

\textbf{4. Comment on APC:}
APC falls from 1.10 to 0.83 as income rises. This is exactly what the Keynesian psychological law predicts: the proportion of income consumed declines because autonomous consumption (0.8) is spread over a larger income base. Note that $APC > MPC$ at all income levels shown, and APC approaches MPC (0.7) as $Y_d$ grows very large.
\end{exoresolu}

\courssub{Practical application: a Cameroonian household budget}

\begin{experience}[Case study: The Njoya family in Bafoussam]
The Njoya family has a monthly disposable income of 450,000 CFAF. Their consumption breakdown is:
\begin{itemize}
    \item Food (non-durable): 180,000 CFAF
    \item Rent \& utilities (services): 100,000 CFAF
    \item Transport (services): 30,000 CFAF
    \item Education/health (services): 40,000 CFAF
    \item Clothing (non-durable): 20,000 CFAF
    \item Phone airtime/internet (services): 15,000 CFAF
    \item Maintenance of motorcycle (durable repair): 10,000 CFAF
\end{itemize}
Total consumption $C = 395,000$ CFAF.
Saving $S = Y_d - C = 55,000$ CFAF.
$APC = 395/450 = 0.878$.
$MPC$ cannot be calculated from a single observation; we would need data from another month with different income.

\textbf{Question:} If the family receives a 50,000 CFAF bonus (transitory income) and their MPC out of transitory income is 0.3 (permanent income hypothesis), by how much does consumption rise?
\textbf{Answer:} $\Delta C = 0.3 \times 50,000 = 15,000$ CFAF. Most of the bonus is saved.
\end{experience}

\courssub{Summary of key points}

\begin{aretenir}[Consumption and the consumption function — what you must know]
\begin{itemize}
    \item \textbf{Consumption ($C$)} is household spending on final goods (durable, non-durable) and services.
    \item \textbf{Consumption function:} $C = a + bY_d$.
    \item \textbf{Autonomous consumption ($a$)}: spending at zero income (dissaving/borrowing). $a > 0$.
    \item \textbf{Induced consumption ($bY_d$)}: varies with income. $b = MPC$, $0 < b < 1$.
    \item \textbf{Keynes' psychological law:} As income rises, consumption rises but by a smaller proportion $\Rightarrow$ APC falls, $APC > MPC$.
    \item \textbf{APC} $= C/Y_d$ (average). \textbf{MPC} $= \Delta C/\Delta Y_d$ (marginal, slope).
    \item \textbf{Graph:} Consumption function cuts vertical axis at $a$, crosses 45° line at break-even income $Y_0 = a/(1-b)$. Below $Y_0$: dissaving. Above $Y_0$: saving.
    \item \textbf{Determinants shifting the function:} wealth, interest rates, expectations, income distribution, credit availability, fiscal policy.
    \item \textbf{Alternative hypotheses:} Relative income (Duesenberry), Permanent income (Friedman), Life-cycle (Modigliani) — all explain why long-run APC is roughly constant while short-run APC falls.
\end{itemize}
\end{aretenir}

\begin{exercice}[Practice questions]
\begin{enumerate}
    \item A household has the consumption function $C = 150 + 0.75Y_d$. Calculate the break-even level of income.
    \item If disposable income rises from 1,000 to 1,200 (in billions CFAF) and consumption rises from 850 to 1,000, calculate MPC and the new APC.
    \item Explain why the Life-Cycle Hypothesis predicts that an aging population (like in many Western countries) will lower the national saving rate.
    \item Distinguish between the \emph{demonstration effect} (Duesenberry) and the \emph{ratchet effect}.
    \item \textbf{Data response:} Using the table in the worked example above, calculate the Average Propensity to Save (APS) and Marginal Propensity to Save (MPS) for Year 3. Verify that $APC + APS = 1$ and $MPC + MPS = 1$.
\end{enumerate}
\end{exercice}

\begin{corrige}[Exercise 1]
Break-even: $Y = C \Rightarrow Y = 150 + 0.75Y \Rightarrow 0.25Y = 150 \Rightarrow Y = 600$ (billion CFAF).
\end{corrige}

\begin{corrige}[Exercise 2]
$\Delta Y = 200$, $\Delta C = 150 \Rightarrow MPC = 150/200 = 0.75$.
New $Y = 1200$, new $C = 1000 \Rightarrow APC = 1000/1200 = 0.833$.
\end{corrige}

\begin{corrige}[Exercise 5]
Year 3: $Y_d = 4.0$, $C = 3.6 \Rightarrow S = 0.4$.
$APS = S/Y = 0.4/4.0 = 0.1$. $APC = 0.9$. $APC + APS = 1.0$ ✓.
$\Delta Y = 1.0$, $\Delta C = 0.7 \Rightarrow \Delta S = 0.3$.
$MPS = \Delta S/\Delta Y = 0.3/1.0 = 0.3$. $MPC = 0.7$. $MPC + MPS = 1.0$ ✓.
\end{corrige}

\coursec{Saving and the saving function}

In the previous section we studied consumption: we saw that households divide their disposable income between consumption and ``what is left over''. That leftover is precisely the subject of this section. Every franc of disposable income that is \emph{not} spent on goods and services is \cle{saved}. Just as we built a consumption function $C = a + bY$, we shall now build a \cle{saving function}, define the average and marginal propensities to save, show how the saving curve is derived graphically from the consumption curve, and finish with one of the most famous results in macroeconomics: the \cle{paradox of thrift}.

\courssub{Definition of saving and dissaving}

\begin{definition}[Saving and dissaving]
\cle{Saving} ($S$) is the part of \cle{disposable income} that is not consumed:
\[ S = Y - C \]
where $Y$ is disposable national (or household) income and $C$ is consumption expenditure.
\begin{itemize}
\item When $Y > C$, saving is \textbf{positive}: the household (or nation) puts income aside.
\item When $Y = C$, saving is \textbf{zero}: this is the \cle{break-even point}.
\item When $Y < C$, saving is \textbf{negative}: this is called \cle{dissaving}. The household finances its consumption by drawing on past savings or by borrowing.
\end{itemize}
\end{definition}

\begin{exemplebox}[A household in Bamenda]
A civil servant in Bamenda earns a disposable income of $200\,000$ FCFA per month and spends $160\,000$ FCFA on food, transport, school fees and airtime. Her saving is
\[ S = Y - C = 200\,000 - 160\,000 = 40\,000 \text{ FCFA}. \]
If in a difficult month she spends $220\,000$ FCFA (for example because of a funeral contribution), then $S = -20\,000$ FCFA: she is \emph{dissaving}, financing the gap from her previous savings or from a loan.
\end{exemplebox}

\begin{attention}[Saving is a withdrawal, not an investment]
In everyday language, ``saving'' and ``investing'' are often confused. In macroeconomics they are different acts, often performed by different people. \cle{Saving} is a \textbf{withdrawal (leakage)} from the circular flow of income: it is income that does not return immediately to firms as spending. \cle{Investment} is an \textbf{injection}: spending by firms on capital goods. Saving is \emph{not} automatically invested. Money kept at home under a mattress, or even deposited in a bank that does not lend it out, leaks out of the flow without financing any purchase of machines or buildings. Equality between saving and investment is a \emph{condition of equilibrium} (studied in the section on the equilibrium level of income), not an automatic identity of behaviour.
\end{attention}

\courssub{Deriving the saving function from the consumption function}

Since disposable income is either consumed or saved, $Y = C + S$. Substituting the Keynesian consumption function $C = a + bY$ (where $a > 0$ is autonomous consumption and $b$ is the marginal propensity to consume, $0 < b < 1$):

\begin{align*}
S &= Y - C \\
  &= Y - (a + bY) \\
  &= -a + (1-b)Y.
\end{align*}

\begin{propriete}[The saving function]
If the consumption function is $C = a + bY$, then the saving function is
\[ S = -a + (1-b)Y = -a + sY, \]
where $s = 1 - b$ is the \cle{marginal propensity to save}. The intercept is $-a$ (negative): at zero income, households still consume $a$, so they must dissave by $a$. The break-even income, where $S = 0$, is
\[ Y^{*} = \frac{a}{1-b} = \frac{a}{s}. \]
\end{propriete}

The economic logic is simple: whatever is not consumed out of an extra franc of income is saved. If the MPC is $b = 0.8$, then out of each additional franc of income, $0.80$ FCFA is consumed and $0.20$ FCFA is saved, so $s = 0.2$.

\begin{exemplebox}[From consumption to saving]
Let $C = 100 + 0.8Y$ (in billions of FCFA). Then
\[ S = -100 + (1-0.8)Y = -100 + 0.2Y. \]
Break-even income: $Y^{*} = \dfrac{100}{0.2} = 500$ billion FCFA. Check: at $Y = 500$, $C = 100 + 0.8(500) = 500$, so $S = 0$. At $Y = 800$, $S = -100 + 0.2(800) = 60$ billion FCFA.
\end{exemplebox}

\courssub{Average and marginal propensity to save}

\begin{definition}[APS and MPS]
The \cle{average propensity to save} (APS) is the fraction of total income that is saved:
\[ APS = \frac{S}{Y}. \]
The \cle{marginal propensity to save} (MPS) is the fraction of an \emph{additional} unit of income that is saved:
\[ MPS = \frac{\Delta S}{\Delta Y} = 1 - b = s. \]
For the linear saving function $S = -a + sY$, the MPS is constant and equal to the slope $s$, while the APS $= -a/Y + s$ rises as income rises (it is negative below the break-even point and approaches $s$ as $Y$ becomes very large).
\end{definition}

\begin{propriete}[Fundamental relationships]
Since $Y = C + S$, dividing through by $Y$ gives
\[ APC + APS = 1 \]
and since any change in income is either consumed or saved, $\Delta Y = \Delta C + \Delta S$, so dividing by $\Delta Y$ gives
\[ MPC + MPS = 1. \]
These identities hold \emph{always}, for any consumption/saving function. They are frequently tested at the GCE and are the key to quick calculations.
\end{propriete}

\begin{exemplebox}[Numerical verification]
Take $C = 100 + 0.8Y$ at $Y = 1\,000$ billion FCFA. Then $C = 900$ and $S = 100$.
\[ APC = \frac{900}{1000} = 0.9, \qquad APS = \frac{100}{1000} = 0.1, \qquad APC + APS = 0.9 + 0.1 = 1. \checkmark \]
Now let income rise to $Y = 1\,100$. Then $C = 980$, $S = 120$, so $\Delta C = 80$, $\Delta S = 20$, $\Delta Y = 100$:
\[ MPC = \frac{80}{100} = 0.8, \qquad MPS = \frac{20}{100} = 0.2, \qquad MPC + MPS = 0.8 + 0.2 = 1. \checkmark \]
Note also that $APC = 0.9 \neq MPC = 0.8$: the average and marginal propensities are equal \emph{only} when autonomous consumption is zero.
\end{exemplebox}

\begin{attention}[Common mistakes]
\begin{itemize}
\item Confusing $APS$ with $MPS$: the APS uses \emph{levels} ($S/Y$), the MPS uses \emph{changes} ($\Delta S / \Delta Y$).
\item Writing the saving function as $S = a + (1-b)Y$. The intercept is $-a$, because at $Y = 0$ consumption $a$ is financed by dissaving.
\item Believing that saving can never be negative. Below the break-even income, dissaving is normal and observed (students, unemployed persons, retirees).
\item Forgetting that $APC + APS = 1$ holds at \emph{every} level of income, not only at equilibrium.
\end{itemize}
\end{attention}

\courssub{Graphical derivation of the saving curve}

The saving curve is obtained directly from the consumption curve. On a diagram with income on the horizontal axis, draw the $45^{\circ}$ line (on which $C = Y$) and the consumption line $C = a + bY$. The vertical gap between the $45^{\circ}$ line and the consumption line \emph{is} saving: $S = Y - C$. Where the consumption line cuts the $45^{\circ}$ line, the gap is zero: this is the break-even point, and the saving curve crosses the income axis exactly below it. To the left of that point, consumption exceeds income (dissaving, $S < 0$); to the right, saving is positive and grows at the rate $s = 1-b$.

\begin{popfigure}
\begin{center}
\begin{tikzpicture}[scale=0.9]
% Top panel: consumption (C = 1 + 0.5Y, break-even at Y = 2)
\begin{scope}
\draw[->, thick] (0,0) -- (8.6,0) node[below left, font=\small] {$Y$};
\draw[->, thick] (0,0) -- (0,5.6) node[left, font=\small] {$C,\ Y$};
\draw[popmuted, dashed] (0,0) -- (5.2,5.2) node[above right=-2pt, font=\small] {$45^{\circ}$};
\draw[very thick, popblue] (0,1) -- (8.2,5.1) node[above left=-2pt, font=\small] {$C = a + bY$};
\fill[poporange] (2,2) circle (2pt);
\draw[poporange, dashed] (2,0) -- (2,2);
\node[below, font=\small, poporange] at (2,0) {$Y^{*}$};
\draw[<->, popgreen, thick] (4.5,3.25) -- (4.5,4.5);
\node[right, font=\small, popgreen] at (4.55,3.9) {$S>0$};
\draw[<->, poppink, thick] (1,1) -- (1,1.5);
\node[right, font=\small, poppink] at (1.05,1.25) {$S<0$};
\node[font=\small\bfseries, popink] at (6.8,5.5) {Consumption};
\end{scope}
% Bottom panel: saving (S = -1 + 0.5Y)
\begin{scope}[yshift=-4.8cm]
\draw[->, thick] (0,0) -- (8.6,0) node[below left, font=\small] {$Y$};
\draw[->, thick] (0,-1.6) -- (0,2.8) node[left, font=\small] {$S$};
\draw[very thick, poppurple] (0,-1) -- (6.6,2.3) node[above left=-2pt, font=\small] {$S = -a + sY$};
\fill[poporange] (2,0) circle (2pt);
\draw[poporange, dashed] (2,0) -- (2,1.2);
\node[below right=-2pt, font=\small, poporange] at (2,0) {$Y^{*}$};
\node[left, font=\small, poppurple] at (0,-1) {$-a$};
\node[font=\small, popgreen] at (5.6,1.1) {saving};
\node[font=\small, poppink] at (1,-1.35) {dissaving};
\node[font=\small\bfseries, popink] at (6.8,2.6) {Saving};
\end{scope}
\end{tikzpicture}
\end{center}
\legende{Derivation of the saving curve from the consumption curve. The saving curve crosses the income axis at the break-even income $Y^{*}$, exactly below the intersection of the consumption line with the $45^{\circ}$ line. To the left of $Y^{*}$ the consumption line lies above the $45^{\circ}$ line (dissaving); to the right it lies below (positive saving).}
\end{popfigure}

The same saving function can be drawn on its own:

\begin{popfigure}
\begin{center}
\begin{tikzpicture}
\begin{axis}[
  width=11cm, height=7cm,
  axis lines=left,
  xlabel={Income $Y$ (billions FCFA)}, ylabel={Saving $S$},
  xmin=0, xmax=1000, ymin=-150, ymax=150,
  xtick={0,500,1000}, ytick={-100,0,100},
  grid=both, grid style={popline, very thin},
  legend style={at={(0.03,0.97)}, anchor=north west, font=\small},
  clip=false]
\addplot[very thick, poppurple, domain=0:1000] {-100 + 0.2*x};
\addlegendentry{$S = -100 + 0.2Y$}
\addplot[popmuted, dashed, domain=0:1000] {0};
\addplot[only marks, mark=*, poporange, mark size=2pt] coordinates {(500,0)};
\node[font=\small, poporange, anchor=south west] at (axis cs:510,5) {break-even $Y^{*}=500$};
\node[font=\small, poppink] at (axis cs:180,-70) {dissaving zone};
\node[font=\small, popgreen] at (axis cs:760,90) {saving zone};
\end{axis}
\end{tikzpicture}
\end{center}
\legende{The saving function $S = -100 + 0.2Y$: negative below the break-even income $Y^{*} = 500$, zero at $Y^{*}$, and positive above it, with slope $MPS = 0.2$.}
\end{popfigure}

\courssub{Determinants of saving}

Why do households and nations save more or less? The main determinants are:

\begin{itemize}
\item \textbf{The level of income.} This is the central Keynesian determinant: saving rises with disposable income ($MPS > 0$). Richer households save a higher \emph{proportion} of their income; the poorest households often dissave.
\item \textbf{The rate of interest.} A higher interest rate raises the reward for saving and may encourage it, but the effect is ambiguous: for people saving towards a fixed target (for example, a plot of land in Douala), a higher return means they need to save \emph{less} each month. Empirically, income matters far more than the interest rate for most households.
\item \textbf{Inflation.} When prices rise quickly, the real value of money savings falls. Households may save less in money form and prefer real assets (land, livestock, building materials). Conversely, uncertainty created by inflation can push some households to accumulate precautionary balances.
\item \textbf{Precautionary and life-cycle motives.} People save to face emergencies (illness, funerals, unemployment), to finance education, to prepare for retirement, or to leave wealth to their children. Where social protection is weak, as in much of Cameroon, precautionary saving is particularly strong.
\item \textbf{Financial institutions and habits.} The ease and safety of saving matter enormously. Commercial banks, microfinance institutions (MC$^2$, credit unions such as CAMCCUL), mobile money services, and the traditional \cle{tontines} or \cle{njangis} — rotating savings groups in which members contribute a fixed amount each round and receive the pot in turn — all make saving easier and are deeply rooted in Cameroonian economic life.
\item \textbf{Government policy and income distribution.} Taxation reduces disposable income and hence saving; since richer households have a higher APS, a more unequal distribution of income tends to raise aggregate saving.
\end{itemize}

\begin{savaistu}[Tontines and njangis: African financial engineering]
Long before mobile banking, Cameroonian communities organised saving through \emph{njangis} (in the North-West and South-West) and \emph{tontines}. Members contribute a fixed sum at each meeting and the total is handed to one member in rotation. These institutions perform three functions of a financial system: they mobilise saving, they enforce discipline (social pressure replaces legal contracts), and they provide credit to members who receive the pot early. Economists now study them seriously as a form of \emph{rotating savings and credit association} (ROSCA).
\end{savaistu}

\courssub{The paradox of thrift}

\begin{definition}[The paradox of thrift]
The \cle{paradox of thrift} states that what is virtuous for an individual can be harmful for the economy as a whole: if \emph{all} households simultaneously decide to save a larger fraction of their income, aggregate consumption falls, firms sell less, national income falls through the multiplier process, and total saving may end up \textbf{no higher — or even lower — than before}.
\end{definition}

The intuition runs as follows. One household saving more becomes richer, because its own saving does not reduce its own income. But saving is income not spent: one person's spending is another person's income. If everyone cuts consumption, firms' revenues fall, they lay off workers or cut orders, incomes fall, and since saving depends on income ($S = -a + sY$), the \emph{actual} amount saved falls with income even though the \emph{desire} to save (the propensity to save) has risen.

\begin{exemplebox}[A numerical illustration of the paradox]
Suppose the initial saving function is $S = -100 + 0.2Y$ and equilibrium income is $Y = 1\,000$, so $S = 100$. Now households become more thrifty: the MPS rises from $0.2$ to $0.3$, so planned saving becomes $S' = -100 + 0.3Y$. At the \emph{old} income of $1\,000$, planned saving would be $200$. But consumption has fallen by $100$, which reduces demand, output and income. In a simple closed economy where equilibrium requires $S = I$ with investment fixed at $I = 100$, the new equilibrium satisfies
\[ -100 + 0.3Y = 100 \quad\Longrightarrow\quad Y = \frac{200}{0.3} \approx 666.7. \]
Actual saving is then $S' = -100 + 0.3(666.7) = 100$: exactly the same as before, but with a national income lower by about $333$ billion FCFA. If investment itself falls when sales collapse, total saving can even \emph{decrease}. The attempt to save more has impoverished the nation without increasing saving.
\end{exemplebox}

\begin{aretenir}[Key points]
\begin{itemize}
\item Saving is the part of disposable income not consumed: $S = Y - C$; negative saving is \emph{dissaving}.
\item From $C = a + bY$ we derive $S = -a + (1-b)Y$; the break-even income is $Y^{*} = a/(1-b)$.
\item $APS = S/Y$; $MPS = \Delta S/\Delta Y = 1 - b$; and always $APC + APS = 1$, $MPC + MPS = 1$.
\item Saving depends mainly on income, but also on interest rates, inflation, precautionary motives and financial institutions (banks, microfinance, tontines/njangis).
\item Saving is a \emph{withdrawal} from the circular flow; it is not automatically invested.
\item The paradox of thrift: a general rise in the desire to save reduces national income and may leave total saving unchanged or reduced.
\end{itemize}
\end{aretenir}

\begin{exoresolu}[Saving function, APS, MPS and break-even point]
In the economy of ``Kamerland'', the consumption function is $C = 60 + 0.75Y$ (all figures in billions of FCFA).
\begin{enumerate}
\item Derive the saving function and find the break-even level of income.
\item Compute saving, APS and MPS when $Y = 400$.
\item Verify that $APC + APS = 1$ and $MPC + MPS = 1$ at $Y = 400$.
\item Households decide to save more: the MPS rises to $0.4$. Explain, without calculation, what the paradox of thrift predicts.
\end{enumerate}
\tcblower
\textbf{1. Saving function and break-even point.}
\[ S = Y - C = Y - (60 + 0.75Y) = -60 + 0.25Y. \]
Break-even: $S = 0 \Rightarrow 0.25Y = 60 \Rightarrow Y^{*} = \dfrac{60}{0.25} = 240$ billion FCFA.

\textbf{2. Saving, APS and MPS at $Y = 400$.}
\[ S = -60 + 0.25(400) = 40 \text{ billion FCFA}. \]
\[ APS = \frac{S}{Y} = \frac{40}{400} = 0.10, \qquad MPS = 1 - b = 0.25. \]

\textbf{3. Verification of the identities.}
At $Y = 400$: $C = 60 + 0.75(400) = 360$.
\[ APC = \frac{360}{400} = 0.90 \quad\Rightarrow\quad APC + APS = 0.90 + 0.10 = 1. \checkmark \]
\[ MPC = 0.75 \quad\Rightarrow\quad MPC + MPS = 0.75 + 0.25 = 1. \checkmark \]

\textbf{4. Paradox of thrift.} If the MPS rises to $0.4$, planned saving exceeds planned investment at the old income, consumption demand falls, and through the multiplier national income contracts. Actual saving, which depends on income, falls back towards the (unchanged) level of investment. The economy ends up with lower income and no increase — possibly a decrease — in total saving: the attempt to be more thrifty has made the nation poorer.
\end{exoresolu}

\begin{exercice}[Practice]
The saving function of an economy is $S = -80 + 0.2Y$ (billions of FCFA).
\begin{enumerate}
\item Write down the consumption function.
\item Find the break-even level of income.
\item Compute $S$, $APS$ and $APC$ at $Y = 600$ and at $Y = 300$. Comment on the sign of saving in each case.
\item Income rises from $600$ to $700$. Compute $\Delta S$ and verify the value of the MPS.
\item Give two reasons, using Cameroonian examples, why a low-income household may dissave.
\end{enumerate}
\end{exercice}

\begin{corrige}[Exercise 1]
\textbf{1.} $C = Y - S = Y - (-80 + 0.2Y) = 80 + 0.8Y$.

\textbf{2.} $Y^{*} = \dfrac{80}{0.2} = 400$ billion FCFA.

\textbf{3.} At $Y = 600$: $S = -80 + 0.2(600) = 40$; $APS = 40/600 \approx 0.067$; $APC = 1 - 0.067 = 0.933$. Saving is positive (income above break-even). At $Y = 300$: $S = -80 + 0.2(300) = -20$; $APS = -20/300 \approx -0.067$; $APC = 1 - (-0.067) \approx 1.067$. Saving is negative: the economy dissaves (income below break-even).

\textbf{4.} $\Delta S = 0.2 \times 100 = 20$ billion FCFA; $MPS = \Delta S/\Delta Y = 20/100 = 0.2$, equal to the slope of the saving function, as expected.

\textbf{5.} For example: (i) a household facing an emergency (illness, funeral contribution) spends more than its current income and borrows from a tontine; (ii) a student or unemployed person with little or no income consumes out of family transfers or past savings.
\end{corrige}

\coursec{Practical work: calculating APC, MPC, APS and MPS}

In the previous section we established that saving is the mirror image of consumption: every franc of disposable income is either consumed or saved, so that $Y = C + S$. We also defined four \cle{propensities} — the average and marginal propensities to consume and to save — and we proved the two identities $APC + APS = 1$ and $MPC + MPS = 1$. Knowing the formulas is one thing; the GCE examiner, however, rarely asks you to recite them. Instead, you are given a \emph{data table} (income and consumption figures for a household, a community or a whole economy) and asked to \textbf{compute} the propensities, \textbf{derive} the consumption and saving functions, and \textbf{interpret} your results. This practical work session trains you to do exactly that, quickly, accurately, and in the format expected in the examination hall. It corresponds directly to \textbf{TU 12: Practical work 14} of the official syllabus.

\courssub{Recap of the four propensity formulas}

Before touching any data, let us fix the four formulas firmly in mind. The \emph{average} propensities relate a \textbf{level} of consumption or saving to the \textbf{level} of income at a single point; the \emph{marginal} propensities relate a \textbf{change} in consumption or saving to the \textbf{change} in income between two points.

\begin{propriete}[Formula sheet: the four propensities]
For disposable income $Y$, consumption $C$ and saving $S = Y - C$ (all measured in the same money unit, e.g.\ FCFA):
\begin{align*}
APC &= \frac{C}{Y} & APS &= \frac{S}{Y} = \frac{Y - C}{Y} \\[4pt]
MPC &= \frac{\Delta C}{\Delta Y} & MPS &= \frac{\Delta S}{\Delta Y} = \frac{\Delta Y - \Delta C}{\Delta Y}
\end{align*}
All four are \textbf{pure numbers} (ratios without units), usually expressed as decimals (0.75) or percentages (75\%). The identities $APC + APS = 1$ and $MPC + MPS = 1$ always hold. In the Keynesian consumption function $C = a + bY$, the marginal propensity to consume is exactly the slope: $MPC = b$.
\end{propriete}

\begin{attention}[MPC uses changes, never levels]
The single most frequent error in GCE scripts is computing $MPC = C/Y$ instead of $MPC = \Delta C / \Delta Y$. The average propensity uses \emph{levels}; the marginal propensity uses \emph{changes}. If income rises from $100\,000$ to $120\,000$ FCFA and consumption rises from $90\,000$ to $106\,000$ FCFA, then $MPC = \dfrac{106\,000 - 90\,000}{120\,000 - 100\,000} = \dfrac{16\,000}{20\,000} = 0.8$ — \textbf{not} $106\,000/120\,000 = 0.883$. Write the $\Delta$ symbols explicitly in your working; it protects you from this trap.
\end{attention}

\courssub{Constructing the income–consumption–saving schedule}

Exam questions typically give you two columns of raw data: income $Y$ and consumption $C$ at several levels. Your first task is to \emph{organise} the data before computing anything.

\begin{methode}[Four-step procedure for computing all four propensities]
\begin{enumerate}
\item \textbf{Complete the saving column.} For each row, compute $S = Y - C$. Check the sign: negative saving (dissaving) is possible at low incomes.
\item \textbf{Compute the average propensities row by row.} For each income level, $APC = C/Y$ and $APS = S/Y$. Immediately verify $APC + APS = 1$ on each row.
\item \textbf{Compute the marginal propensities between consecutive rows.} For each pair of successive income levels, $\Delta Y = Y_2 - Y_1$, $\Delta C = C_2 - C_1$, $\Delta S = S_2 - S_1$; then $MPC = \Delta C / \Delta Y$ and $MPS = \Delta S / \Delta Y$. Note that marginal propensities are written \emph{between} rows, so a table with $n$ income levels yields $n - 1$ marginal values.
\item \textbf{Verify and interpret.} Check $MPC + MPS = 1$ for each interval, check that $MPC$ is constant if the data are linear, then state the economic meaning of your results in words.
\end{enumerate}
\end{methode}

\courssub{A complete worked example}

\begin{experience}[Case study: the Mbarga household in Yaoundé]
The Mbarga household's monthly disposable income and consumption were recorded at five different income levels (in thousands of FCFA):

\begin{center}
\begin{tabular}{|c|c|}
\hline
\rowcolor{popblueL} Income $Y$ ('000 FCFA) & Consumption $C$ ('000 FCFA) \\
\hline
100 & 110 \\
\hline
200 & 190 \\
\hline
300 & 270 \\
\hline
400 & 350 \\
\hline
500 & 430 \\
\hline
\end{tabular}
\end{center}

We apply the four-step method. \textbf{Step 1:} saving is $S = Y - C$: at $Y = 100$, $S = -10$ (dissaving); at $Y = 200$, $S = 10$; and so on. \textbf{Steps 2 and 3} give the full schedule below.
\end{experience}

\begin{popfigure}
\begin{center}
\begin{tabular}{|c|c|c|c|c|c|c|}
\hline
\rowcolor{popblueL} $Y$ & $C$ & $S=Y-C$ & $APC=C/Y$ & $APS=S/Y$ & $MPC=\Delta C/\Delta Y$ & $MPS=\Delta S/\Delta Y$ \\
\hline
100 & 110 & $-10$ & 1.10 & $-0.10$ & --- & --- \\
\hline
200 & 190 & 10 & 0.95 & 0.05 & 0.8 & 0.2 \\
\hline
300 & 270 & 30 & 0.90 & 0.10 & 0.8 & 0.2 \\
\hline
400 & 350 & 50 & 0.875 & 0.125 & 0.8 & 0.2 \\
\hline
500 & 430 & 70 & 0.86 & 0.14 & 0.8 & 0.2 \\
\hline
\end{tabular}
\end{center}
\legende{Complete propensity schedule for the Mbarga household (figures in '000 FCFA). Marginal propensities are computed between consecutive rows; the first row has none.}
\end{popfigure}

Let us detail the computations so that you can reproduce them under exam conditions.

\textbf{Step 2 in detail — average propensities.} Take the row $Y = 300$, $C = 270$:
\begin{align*}
APC &= \frac{C}{Y} = \frac{270}{300} = 0.90, &
APS &= \frac{S}{Y} = \frac{30}{300} = 0.10.
\end{align*}
Check: $APC + APS = 0.90 + 0.10 = 1$. \checkmark{} At $Y = 100$, $APC = 110/100 = 1.10$ and $APS = -10/100 = -0.10$: the household consumes \emph{more} than its income, financing the gap by borrowing or drawing down past savings. Again $1.10 + (-0.10) = 1$. \checkmark{}

\textbf{Step 3 in detail — marginal propensities.} Between $Y = 200$ and $Y = 300$:
\begin{align*}
\Delta Y &= 300 - 200 = 100, & \Delta C &= 270 - 190 = 80, & \Delta S &= 30 - 10 = 20, \\
MPC &= \frac{\Delta C}{\Delta Y} = \frac{80}{100} = 0.8, & MPS &= \frac{\Delta S}{\Delta Y} = \frac{20}{100} = 0.2.
\end{align*}
Check: $MPC + MPS = 0.8 + 0.2 = 1$. \checkmark{} The same calculation on every interval gives $MPC = 0.8$ and $MPS = 0.2$: the data are perfectly linear.

\begin{aretenir}[Exam tip: verify before you move on]
Before leaving any propensity calculation, run two instant checks: (i) $APC + APS = 1$ on every row; (ii) $MPC + MPS = 1$ on every interval. If either fails, you have made an arithmetic slip — find it \emph{now}, because all subsequent questions (consumption function, multiplier, equilibrium income) build on these numbers. Examiners award marks for correct method even when a slip is caught and corrected.
\end{aretenir}

\courssub{Deriving the consumption and saving functions from two observations}

When the $MPC$ is constant, the data lie on a straight line $C = a + bY$, and we can recover the two parameters. We already know $b = MPC = 0.8$. To find the autonomous consumption $a$, substitute \emph{any} observation into the equation. Using $Y = 300$, $C = 270$:
\begin{align*}
C &= a + bY \\
270 &= a + 0.8 \times 300 \\
270 &= a + 240 \\
a &= 30.
\end{align*}
Hence the consumption function is $\boxed{C = 30 + 0.8Y}$ (in thousands of FCFA). The saving function follows immediately:
\begin{align*}
S &= Y - C = Y - (30 + 0.8Y) = -30 + 0.2Y,
\end{align*}
that is, $S = -a + (1-b)Y = -30 + 0.2Y$. Notice the symmetry: autonomous saving is $-a = -30$ (dissaving when income is zero), and the slope of the saving function is $MPS = 1 - b = 0.2$.

\begin{methode}[Finding $a$ and $b$ from two observations]
\begin{enumerate}
\item Compute $b = MPC = \dfrac{C_2 - C_1}{Y_2 - Y_1}$ from the two given points.
\item Substitute one point into $C = a + bY$ and solve for $a = C_1 - bY_1$.
\item Write the saving function directly as $S = -a + (1-b)Y$.
\item \textbf{Validate} by checking that the second point satisfies your equation: $a + bY_2$ must equal $C_2$.
\end{enumerate}
\end{methode}

\begin{popfigure}
\begin{center}
\begin{tikzpicture}
\begin{axis}[
    width=11cm, height=8cm,
    xlabel={Income $Y$ ('000 FCFA)}, ylabel={Consumption $C$ ('000 FCFA)},
    xmin=0, xmax=550, ymin=0, ymax=550,
    grid=both, grid style={popline, very thin},
    legend style={at={(0.03,0.97)}, anchor=north west, font=\small},
    axis lines=left,
]
\addplot[domain=0:550, samples=2, thick, popblue] {30 + 0.8*x};
\addlegendentry{$C = 30 + 0.8Y$}
\addplot[domain=0:550, samples=2, dashed, popmuted] {x};
\addlegendentry{$45^{\circ}$ line ($C = Y$)}
\addplot[only marks, mark=*, mark size=2.2pt, poporange] coordinates {(100,110) (200,190) (300,270) (400,350) (500,430)};
\addlegendentry{Observed data}
\draw[dotted, popdark] (axis cs:100,0) -- (axis cs:100,110);
\draw[dotted, popdark] (axis cs:0,110) -- (axis cs:100,110);
\node[font=\small, popdark, anchor=west] at (axis cs:60,150) {break-even};
\node[font=\small, popdark, anchor=west] at (axis cs:60,135) {$Y = 150$};
\end{axis}
\end{tikzpicture}
\end{center}
\legende{The observed data points and the fitted consumption line $C = 30 + 0.8Y$. Below the break-even income of 150 (where the line crosses the $45^{\circ}$ line), the household dissaves; above it, saving is positive.}
\end{popfigure}

The graph confirms the table: all five points lie exactly on the fitted line, the line crosses the $45^{\circ}$ line where $C = Y$, i.e.\ at $30 + 0.8Y = Y$, giving the \cle{break-even income} $Y = 30/0.2 = 150$ thousand FCFA. Below 150, consumption exceeds income ($APC > 1$, dissaving); above 150, saving is positive.

\courssub{Interpreting the results economically}

Numbers without interpretation earn only half the marks. Here is how an examiner expects you to read the results.

An $MPC$ of $0.8$ means that \emph{out of every additional 100 FCFA of disposable income received, the household spends 80 FCFA on consumption and saves 20 FCFA}. For a Cameroonian household like the Mbargas, this is plausible: when a salary earner in Yaoundé receives a pay rise or a trader in Mokolo market enjoys a better month, most of the extra income goes to food, school fees, transport and airtime, while a smaller share is set aside — perhaps in a \emph{tontine} or a savings account. The complementary $MPS$ of $0.2$ tells us the leakage into saving at the margin.

The \emph{average} propensities tell a different story: they describe the \textbf{whole} income, not the extra franc. At $Y = 100$, $APC = 1.10$: the household spends more than it earns — a typical situation for low-income households who must borrow or draw on family support simply to meet basic needs. As income rises, $APC$ falls steadily ($1.10 \to 0.95 \to 0.90 \to 0.875 \to 0.86$) while $APS$ rises: richer households devote a smaller \emph{share} of their total income to consumption, even though their consumption in absolute terms keeps rising. This illustrates Keynes's \cle{fundamental psychological law}: as income increases, consumption increases, but by less than the increase in income.

\begin{attention}[Do not confuse APC with MPC in your interpretation]
Saying ``$APC = 0.9$ means the household spends 90\% of any extra income'' is wrong — that is the meaning of $MPC$. $APC = 0.9$ means the household spends 90\% of its \emph{total} income. Also note that $APC$ can exceed 1 (dissaving) whereas in the simple Keynesian model $MPC$ always lies strictly between 0 and 1.
\end{attention}

\begin{savaistu}[Cameroonian context: tontines and the MPC]
In Cameroon, informal rotating savings and credit associations (\emph{tontines}) are a major savings vehicle, especially for women traders in markets like Mokolo (Yaoundé) or Marché Central (Douala). A high MPC (e.g. 0.8) does not mean households do not save; it means that at the \emph{margin}, 80\% of extra income is consumed. The 20\% marginal saving often flows into tontines, mobile money (MTN MoMo, Orange Money) or informal lending circles. This behaviour influences the size of the national multiplier, which we study later in this chapter.
\end{savaistu}

\courssub{Worked examination-style exercise}

\begin{exoresolu}[GCE-style data response]
The table below shows the national disposable income and aggregate consumption of a hypothetical economy (billions of FCFA).

\begin{center}
\begin{tabular}{|c|c|c|}
\hline
\rowcolor{popblueL} Year & Income $Y$ & Consumption $C$ \\
\hline
1 & 400 & 380 \\
\hline
2 & 600 & 520 \\
\hline
\end{tabular}
\end{center}

\begin{enumerate}
\item Calculate the APC and APS in each year.
\item Calculate the MPC and MPS between the two years.
\item Derive the consumption function and the saving function.
\item Determine the break-even level of income.
\end{enumerate}

\tcblower
\textbf{1. Average propensities.}

Year 1: $S_1 = 400 - 380 = 20$; $APC_1 = \dfrac{380}{400} = 0.95$; $APS_1 = \dfrac{20}{400} = 0.05$. Check: $0.95 + 0.05 = 1$. \checkmark{}

Year 2: $S_2 = 600 - 520 = 80$; $APC_2 = \dfrac{520}{600} \approx 0.867$; $APS_2 = \dfrac{80}{600} \approx 0.133$. Check: $0.867 + 0.133 = 1$. \checkmark{}

\textbf{2. Marginal propensities.} $\Delta Y = 600 - 400 = 200$; $\Delta C = 520 - 380 = 140$; $\Delta S = 80 - 20 = 60$.
\[
MPC = \frac{\Delta C}{\Delta Y} = \frac{140}{200} = 0.7, \qquad MPS = \frac{\Delta S}{\Delta Y} = \frac{60}{200} = 0.3.
\]
Check: $0.7 + 0.3 = 1$. \checkmark{}

\textbf{3. Consumption and saving functions.} With $b = 0.7$, substitute Year 1 into $C = a + bY$:
\[
a = C_1 - bY_1 = 380 - 0.7 \times 400 = 380 - 280 = 100.
\]
Hence $C = 100 + 0.7Y$ and $S = -100 + 0.3Y$. Validation with Year 2: $100 + 0.7 \times 600 = 100 + 420 = 520 = C_2$. \checkmark{}

\textbf{4. Break-even income.} Set $C = Y$: $100 + 0.7Y = Y \Rightarrow 0.3Y = 100 \Rightarrow Y = 333.3$ billion FCFA (approximately). Below this level the economy dissaves; above it, aggregate saving is positive.

\textbf{Interpretation.} Out of every additional billion FCFA of national income, 700 million are consumed and 300 million are saved; autonomous consumption of 100 billion FCFA represents the spending that would occur even at zero income, financed by past savings or borrowing.
\end{exoresolu}

\courssub{Practice exercise}

\begin{exercice}[Propensities for the Ekane family]
The Ekane family in Douala records the following monthly figures (thousands of FCFA):

\begin{center}
\begin{tabular}{|c|c|}
\hline
\rowcolor{popblueL} Income $Y$ & Consumption $C$ \\
\hline
150 & 160 \\
\hline
250 & 230 \\
\hline
350 & 300 \\
\hline
450 & 370 \\
\hline
\end{tabular}
\end{center}

\begin{enumerate}
\item Construct the complete schedule: $S$, $APC$, $APS$ for each row, and $MPC$, $MPS$ for each interval.
\item Verify both identities on your computed values.
\item Derive the consumption function $C = a + bY$ and the saving function.
\item Find the break-even income and interpret the MPC in one sentence for the Ekane family.
\end{enumerate}
\end{exercice}

\begin{corrige}[Exercise: the Ekane family]
\textbf{1. Schedule.} Savings: $S = Y - C$ gives $-10$, $20$, $50$, $80$. Averages: at $Y=150$: $APC = 160/150 \approx 1.067$, $APS \approx -0.067$; at $Y=250$: $APC = 0.92$, $APS = 0.08$; at $Y=350$: $APC \approx 0.857$, $APS \approx 0.143$; at $Y=450$: $APC \approx 0.822$, $APS \approx 0.178$. Marginals: each interval has $\Delta Y = 100$, $\Delta C = 70$, $\Delta S = 30$, so $MPC = 0.7$ and $MPS = 0.3$ throughout.

\textbf{2. Identities.} Each row: e.g.\ $1.067 + (-0.067) = 1$; $0.92 + 0.08 = 1$; etc. Each interval: $0.7 + 0.3 = 1$. \checkmark{}

\textbf{3. Functions.} $b = 0.7$; using $(250, 230)$: $a = 230 - 0.7 \times 250 = 230 - 175 = 55$. So $C = 55 + 0.7Y$ and $S = -55 + 0.3Y$. Check with $(450, 370)$: $55 + 0.7 \times 450 = 55 + 315 = 370$. \checkmark{}

\textbf{4. Break-even.} $55 + 0.7Y = Y \Rightarrow Y = 55/0.3 \approx 183.3$ thousand FCFA. Interpretation: out of every additional 100 FCFA earned, the Ekane family spends 70 FCFA and saves 30 FCFA.
\end{corrige}

\begin{aretenir}[Key points to remember]
\begin{itemize}
\item $APC = C/Y$ and $APS = S/Y$ use \textbf{levels}; $MPC = \Delta C/\Delta Y$ and $MPS = \Delta S/\Delta Y$ use \textbf{changes}.
\item Always complete the saving column first ($S = Y - C$); always verify $APC + APS = 1$ and $MPC + MPS = 1$.
\item With a constant MPC, derive the consumption function by $b = MPC$ then $a = C - bY$ from any observation; the saving function is $S = -a + (1-b)Y$.
\item Break-even income is where $C = Y$ (equivalently $S = 0$): $Y = a/(1-b)$.
\item Interpret marginal propensities in terms of \emph{additional} income and average propensities in terms of \emph{total} income.
\end{itemize}
\end{aretenir}

Mastering these mechanical calculations is essential, because they are the raw material of everything that follows: the consumption and saving functions you have just derived will feed directly into the determination of the equilibrium level of national income (TU 16, TU 17, TU 18), and the MPC you have computed will reappear as the engine of the \emph{multiplier} (TU 19, TU 20) and the \emph{accelerator} (TU 21, TU 22, TU 23) later in this chapter.

\coursec{Investment and investment appraisal}

In the previous section we examined how households decide to consume or save out of their disposable income. We now turn to the other major component of aggregate demand: \textbf{investment}. While saving represents a leakage from the circular flow, investment is an injection that expands the economy's productive capacity. Understanding what drives firms to invest, and how they evaluate potential projects, is essential for explaining fluctuations in national income and for designing policies that promote growth in Cameroon and beyond.

\courssub{Definition of investment in economics}

\begin{definition}[Investment (Capital Formation)]
In macroeconomics, \textbf{investment} (denoted $I$) is the expenditure by firms on \textbf{capital goods} — machinery, equipment, buildings, infrastructure, and changes in inventories — that are used to produce other goods and services in the future. It is the addition to the nation's \textbf{capital stock}.
\end{definition}

\begin{attention}[Financial vs. Real Investment]
In everyday language, "investment" often means buying shares, bonds, or land. In economics, these are \textbf{financial investments} — a transfer of ownership of existing assets. They do not directly create new capital goods. Only the purchase of \textbf{newly produced capital goods} (or additions to inventories) counts as macroeconomic investment. When a Cameroonian entrepreneur buys a new cocoa dryer, that is investment; when she buys shares in SONARA, it is saving/portfolio allocation.
\end{attention}

\begin{aretenir}[Key Distinction]
\begin{itemize}
    \item \textbf{Saving (S)}: Income not spent on consumption (a decision by households).
    \item \textbf{Investment (I)}: Expenditure on capital goods (a decision by firms).
    \item In a closed economy without government, $S = I$ in equilibrium, but they are decisions made by different agents for different reasons.
\end{itemize}
\end{aretenir}

\courssub{Types of investment}

\begin{definition}[Gross and Net Investment]
\begin{itemize}
    \item \textbf{Gross Investment ($I_g$)}: Total expenditure on new capital goods in a period.
    \item \textbf{Depreciation (Capital Consumption Allowance, $D$)}: The value of capital worn out or becoming obsolete during the period.
    \item \textbf{Net Investment ($I_n$)}: $I_n = I_g - D$. It measures the actual increase in the capital stock.
\end{itemize}
\end{definition}

\begin{definition}[Autonomous and Induced Investment]
\begin{itemize}
    \item \textbf{Autonomous Investment ($I_a$)}: Investment that does \emph{not} depend on the current level of national income ($Y$). It is driven by long-term expectations, technological breakthroughs, or government policy. In the simple Keynesian model, it is treated as exogenous (a constant).
    \item \textbf{Induced Investment ($I_{ind}$)}: Investment that varies \emph{positively} with national income (or output). Higher income $\rightarrow$ higher expected sales $\rightarrow$ firms expand capacity. Often modelled as $I_{ind} = bY$ ($b>0$).
\end{itemize}
\end{definition}

\begin{definition}[Replacement Investment]
Investment undertaken solely to replace worn-out or obsolete capital, keeping the capital stock constant. It equals depreciation ($D$). Net investment is zero if only replacement investment occurs.
\end{definition}

\begin{savaistu}[Investment in Cameroon]
Cameroon's gross fixed capital formation has averaged around 20--22\% of GDP in recent years. Major drivers include public infrastructure (roads, dams like Nachtigal, the Kribi deep-sea port) and private investment in agro-industry (CDC, SOCAPALM), telecommunications (MTN, Orange), and the emerging SME sector. The distinction between autonomous (e.g., a new highway) and induced (e.g., a new warehouse because cocoa exports are booming) helps policymakers target stimuli.
\end{savaistu}

\courssub{Determinants of investment}

Firms invest when the expected return exceeds the cost of funds. The main determinants are:

\begin{propriete}[Determinants of Investment]
\begin{enumerate}
    \item \textbf{Expected Rate of Profit (or Marginal Efficiency of Capital, MEC)}: The anticipated annual net return (after operating costs) on an extra unit of capital, expressed as a percentage of its cost. Higher expected profit $\rightarrow$ more investment.
    \item \textbf{Rate of Interest ($r$)}: The cost of borrowing (or opportunity cost of own funds). Higher $r$ $\rightarrow$ fewer projects are profitable $\rightarrow$ less investment.
    \item \textbf{Business Expectations (Animal Spirits)}: Optimism or pessimism about future demand, policy stability, exchange rates. In Cameroon, political stability and the business climate (Doing Business indicators) heavily influence expectations.
    \item \textbf{Technological Change}: New technologies (e.g., digital payment systems, improved cocoa varieties) raise the MEC of new capital, triggering investment.
    \item \textbf{Government Policy}: Corporate tax rates, investment allowances, interest rate subsidies (e.g., via the \emph{Fonds National de l'Emploi} or \emph{Programme d'Appui aux PME}), infrastructure provision, and regulatory environment.
    \item \textbf{Level of Income/Output (Accelerator Effect)}: Current or expected changes in $Y$ induce investment to meet demand (covered in detail in Section 10).
\end{enumerate}
\end{propriete}

\courssub{The Marginal Efficiency of Capital (MEC) and the investment decision}

\begin{definition}[Marginal Efficiency of Capital (MEC)]
The MEC of a capital asset is the \textbf{discount rate} that equates the present value of the expected future net returns (cash flows) from that asset to its supply price (purchase cost). For a single project, the MEC is identical to the \textbf{Internal Rate of Return (IRR)}. The \textbf{MEC schedule} (or investment demand curve) shows the relationship between the total capital stock (or investment) and the MEC of the marginal unit; it slopes downward because of diminishing marginal returns to capital.
\end{definition}

\begin{propriete}[Investment Decision Rule]
A profit-maximising firm will undertake an investment project if and only if:
\[
\boxed{\text{MEC} > r}
\]
where $r$ is the prevailing market rate of interest (cost of capital). The firm invests up to the point where the MEC of the last unit of capital equals the interest rate.
\end{propriete}

\begin{experience}[The MEC Schedule]
Imagine a firm ranks all potential projects by their MEC (highest first). This yields a \textbf{MEC schedule} (or investment demand curve), which slopes \textbf{downward}: as more capital is accumulated, the best remaining projects have lower returns. The interest rate $r$ is a horizontal line (the firm is a price-taker in the capital market). The optimal capital stock is where the MEC curve intersects $r$.
\end{experience}

\begin{popfigure}
\begin{tikzpicture}
\begin{axis}[
    width=12cm, height=8cm,
    axis lines=middle,
    xlabel={Capital Stock ($K$)}, ylabel={Rate (\% / MEC, $r$)},
    xmin=0, xmax=10, ymin=0, ymax=20,
    xtick=\empty, ytick={5,10,15},
    yticklabels={5\%,10\%,15\%},
    clip=false,
    domain=0:10,
    samples=100,
]
\addplot[popblue, thick] {20 - 1.5*x} node[right, pos=0.9, popblue] {MEC schedule};
\addplot[poporange, thick, dashed] coordinates {(0,8) (10,8)} node[right, pos=0.9, poporange] {$r = 8\%$};
\addplot[popgreen, thick, dashed] coordinates {(8,0) (8,8)};
\addplot[popgreen, thick, dashed] coordinates {(0,8) (8,8)};
\node[popgreen, anchor=north] at (axis cs:8,0) {$K^*$};
\node[popgreen, anchor=east] at (axis cs:0,8) {$r$};
\end{axis}
\end{tikzpicture}
\legende{The MEC schedule slopes downward. At interest rate $r$, the firm invests up to $K^*$ where MEC = $r$. A fall in $r$ to $r'$ would increase optimal capital to $K'$.}
\end{popfigure}

\begin{exoresolu}[MEC Calculation]
A firm considers buying a machine costing \SI{10,000,000}{FCFA}. It expects net cash inflows of \SI{3,000,000}{FCFA} per year for 5 years, with no scrap value. Calculate the MEC (IRR) and decide whether to invest if the bank lending rate is 12\%.

\tcblower
\textbf{Solution:} The MEC ($r^*$) solves:
\[
10,000,000 = \frac{3,000,000}{(1+r^*)} + \frac{3,000,000}{(1+r^*)^2} + \frac{3,000,000}{(1+r^*)^3} + \frac{3,000,000}{(1+r^*)^4} + \frac{3,000,000}{(1+r^*)^5}
\]
This is an annuity: $10,000,000 = 3,000,000 \times \text{PVIFA}(r^*,5)$.
$\text{PVIFA}(r^*,5) = 10/3 \approx 3.333$.
From annuity tables (or calculator), PVIFA(15\%,5) = 3.352, PVIFA(16\%,5) = 3.274.
Interpolating: $r^* \approx 15.2\%$.
Since MEC (15.2\%) > interest rate (12\%), the firm \textbf{should invest}.
\end{exoresolu}

\courssub{Investment appraisal methods}

Firms use several techniques to evaluate projects. The syllabus requires mastery of Payback Period and Accounting Rate of Return (ARR), and awareness of discounted cash flow (DCF) methods (NPV, IRR, Profitability Index).

\courssub{Payback Period (PBP)}

\begin{definition}[Payback Period]
The \textbf{payback period} is the time required for the cumulative net cash inflows from a project to equal its initial capital outlay. It measures how quickly the investment is recovered.
\end{definition}

\begin{methode}[Calculating Payback Period]
\begin{enumerate}
    \item List the initial outlay (negative cash flow at time 0) and expected annual net cash inflows.
    \item Compute cumulative cash flow year by year.
    \item Identify the year when cumulative cash flow becomes zero or positive.
    \item If it turns positive during a year, interpolate: 
    \[
    \text{PBP} = \text{Year before recovery} + \frac{\text{Unrecovered cost at start of year}}{\text{Cash flow during year}}
    \]
\end{enumerate}
\end{methode}

\begin{propriete}[Advantages and Limitations of Payback]
\begin{center}
\begin{tabularx}{\linewidth}{|l|X|X|}\hline
\rowcolor{popblueL}
\textbf{Aspect} & \textbf{Advantages} & \textbf{Limitations} \\ \hline
\textbf{Simplicity} & Easy to calculate and understand. & Ignores the \textbf{time value of money} (no discounting). \\ \hline
\textbf{Liquidity focus} & Highlights projects that free up cash quickly — vital for SMEs with tight cash flow. & Ignores cash flows \textbf{after} the payback period (may reject profitable long-term projects). \\ \hline
\textbf{Risk proxy} & Shorter payback $\approx$ lower risk (uncertainty increases with time). & No objective \textbf{accept/reject benchmark} (what is a "good" payback?). \\ \hline
\textbf{Capital rationing} & Useful when capital is severely limited. & Does not measure \textbf{profitability} or shareholder wealth maximisation. \\ \hline
\end{tabularx}
\end{center}
\end{propriete}

\courssub{Average (Accounting) Rate of Return (ARR)}

\begin{definition}[Accounting Rate of Return (ARR)]
\[
\text{ARR} = \frac{\text{Average Annual Accounting Profit}}{\text{Average Investment}} \times 100\%
\]
where \textbf{Average Annual Accounting Profit} = (Total Net Cash Inflows $-$ Total Depreciation) / Project Life, and \textbf{Average Investment} = (Initial Cost + Scrap Value) / 2 (or simply Initial Cost / 2 if straight-line depreciation to zero).
\end{definition}

\begin{methode}[Calculating ARR]
\begin{enumerate}
    \item Determine initial cost, scrap value, project life, and annual net cash inflows.
    \item Compute annual depreciation: $\frac{\text{Initial Cost} - \text{Scrap Value}}{\text{Life}}$.
    \item Compute average annual accounting profit = Average Annual Cash Inflow $-$ Annual Depreciation.
    \item Compute average investment = $\frac{\text{Initial Cost} + \text{Scrap Value}}{2}$.
    \item ARR = (Average Annual Accounting Profit / Average Investment) $\times 100\%$.
    \item Compare ARR with a target/hurdle rate set by management.
\end{enumerate}
\end{methode}

\begin{exoresolu}[Payback and ARR Comparison]
Two mutually exclusive projects, A and B, each require an initial outlay of \SI{50,000,000}{FCFA}. Expected net cash inflows (in millions FCFA):

\begin{center}
\begin{tabular}{|c|ccccc|}\hline
Year & 1 & 2 & 3 & 4 & 5 \\ \hline
Project A & 15 & 15 & 15 & 15 & 15 \\
Project B & 5 & 10 & 20 & 30 & 20 \\ \hline
\end{tabular}
\end{center}
Assume straight-line depreciation to zero over 5 years. The firm's target ARR is 18\% and maximum acceptable payback is 3 years.

\tcblower
\textbf{Payback Period:}
\begin{itemize}
    \item Project A: Cumulative: 15, 30, 45, 60. Payback = 3 + (5/15) = \textbf{3.33 years}.
    \item Project B: Cumulative: 5, 15, 35, 65. Payback = 3 + (15/30) = \textbf{3.5 years}.
    \item Both exceed 3 years $\rightarrow$ both rejected on payback grounds.
\end{itemize}

\textbf{ARR:}
\begin{itemize}
    \item Depreciation per year = 50/5 = 10.
    \item Project A: Avg annual profit = 15 $-$ 10 = 5. Avg investment = (50+0)/2 = 25. ARR = 5/25 = \textbf{20\%}.
    \item Project B: Total inflow = 85. Avg annual inflow = 17. Avg annual profit = 17 $-$ 10 = 7. ARR = 7/25 = \textbf{28\%}.
    \item Both exceed 18\% target $\rightarrow$ both accepted on ARR grounds.
\end{itemize}

\textbf{Conflict:} Payback rejects both; ARR accepts both. This illustrates the different philosophies: liquidity vs. profitability.
\end{exoresolu}

\begin{popfigure}
\begin{center}
\begin{tabular}{|c|ccccc|}\hline
\rowcolor{popblueL}
\textbf{Year} & \textbf{0} & \textbf{1} & \textbf{2} & \textbf{3} & \textbf{4} \\ \hline
\rowcolor{popgreenL}
Project X (FCFA m) & -100 & 40 & 40 & 40 & 40 \\
\rowcolor{poporangeL}
Project Y (FCFA m) & -100 & 10 & 20 & 50 & 80 \\ \hline
\end{tabular}
\end{center}
\legende{Cash flows for two competing projects. Project X has a constant cash flow (payback = 2.5 years). Project Y is back-loaded (payback = 3.25 years) but yields higher total return. Payback favours X; NPV/IRR may favour Y.}
\end{popfigure}

\courssub{Further studies: Discounted Cash Flow (DCF) methods}

Payback and ARR ignore the time value of money. DCF methods discount future cash flows to present value using the cost of capital ($r$).

\begin{definition}[Net Present Value (NPV)]
\[
\text{NPV} = \sum_{t=0}^{n} \frac{C_t}{(1+r)^t}
\]
where $C_t$ is net cash flow at time $t$ ($C_0 < 0$), $r$ is the discount rate (cost of capital). \textbf{Decision rule: Accept if NPV > 0.}
\end{definition}

\begin{definition}[Internal Rate of Return (IRR)]
The discount rate $r^*$ that makes NPV = 0. It is the MEC of the project. \textbf{Decision rule: Accept if IRR > cost of capital.}
\end{definition}

\begin{definition}[Profitability Index (PI)]
\[
\text{PI} = \frac{\text{PV of future cash inflows}}{\text{Initial Outlay}} = 1 + \frac{\text{NPV}}{\text{Initial Outlay}}
\]
\textbf{Decision rule: Accept if PI > 1.} Useful for ranking projects under capital rationing.
\end{definition}

\begin{aretenir}[Why Discounting Matters]
A franc today is worth more than a franc tomorrow because of:
\begin{itemize}
    \item \textbf{Opportunity cost}: Money can earn interest.
    \item \textbf{Inflation}: Purchasing power erodes.
    \item \textbf{Risk}: Future cash flows are uncertain.
\end{itemize}
DCF methods (NPV, IRR, PI) are theoretically superior because they account for all cash flows and the time value of money. NPV is the \textbf{gold standard} for maximising shareholder wealth.
\end{aretenir}

\begin{attention}[Common Pitfalls in Investment Appraisal]
\begin{itemize}
    \item \textbf{Using accounting profit instead of cash flows} for NPV/IRR. Depreciation is a non-cash charge; add it back.
    \item \textbf{Ignoring opportunity costs} (e.g., using land already owned — its market rent is a cost).
    \item \textbf{Including sunk costs} (past expenditures that cannot be recovered).
    \item \textbf{Wrong discount rate}: Use the project's risk-adjusted cost of capital, not the firm's average WACC if risks differ.
    \item \textbf{Mutually exclusive projects with different scales/lives}: NPV is reliable; IRR can give conflicting rankings (use incremental IRR or NPV).
\end{itemize}
\end{attention}

\courssub{Investment in the Cameroonian context}

\begin{savaistu}[Investment Landscape in Cameroon]
\begin{itemize}
    \item \textbf{SMEs (PME)}: Over 90\% of firms. Face high interest rates (often 12--18\% from commercial banks), collateral requirements, and limited access to long-term finance. Payback period is often their primary criterion due to cash-flow constraints.
    \item \textbf{Agricultural Projects}: Cocoa, coffee, cotton, palm oil. Long gestation periods (3--5 years for tree crops) make payback unattractive; NPV/IRR with patient capital (e.g., from \emph{Fonds de Développement Agricole} or development partners) is essential.
    \item \textbf{Public Infrastructure}: Roads (e.g., Yaoundé-Douala highway), energy (Nachtigal, Memve'ele dams), ports (Kribi). Appraised using social cost-benefit analysis (shadow prices, externalities), not just financial returns.
    \item \textbf{Policy Tools}: Investment Charter (Law 2013/004) offers tax holidays, customs duty exemptions, and accelerated depreciation for priority sectors. The \emph{Guichet Unique de Création d'Entreprise} reduces setup time.
\end{itemize}
\end{savaistu}

\courssub{Worked exercise: NPV and IRR}

\begin{exoresolu}[NPV and IRR for a Cameroonian Agro-Project]
A cooperative in the South-West region considers investing \SI{200,000,000}{FCFA} in a cocoa processing unit. Expected net cash inflows (after operating costs, tax) are:
Year 1: 40m, Year 2: 50m, Year 3: 60m, Year 4: 70m, Year 5: 80m. Scrap value at end of Year 5: 20m. Cost of capital = 15\%.

\tcblower
\textbf{NPV Calculation:}
\begin{align*}
\text{NPV} &= -200 + \frac{40}{1.15} + \frac{50}{1.15^2} + \frac{60}{1.15^3} + \frac{70}{1.15^4} + \frac{70+20}{1.15^5} \\
&= -200 + 34.78 + 37.81 + 39.46 + 40.07 + 44.75 \\
&= \mathbf{-3.13\ million\ FCFA}
\end{align*}
NPV < 0 $\rightarrow$ \textbf{Reject} at 15\% cost of capital.

\textbf{IRR Estimation:}
Try 14\%: PV inflows $\approx$ 202.5 $\rightarrow$ NPV $\approx$ +2.5.
Try 15\%: NPV = -3.13.
Interpolation: $\text{IRR} \approx 14\% + \frac{2.5}{2.5+3.13} \times 1\% \approx \mathbf{14.44\%}$.
Since IRR (14.44\%) < Cost of Capital (15\%), \textbf{Reject}. The project does not earn enough to cover the cost of funds.

\textbf{Policy Insight:} If a development bank offers a concessional loan at 10\%, NPV becomes positive and the project becomes viable. This illustrates the role of targeted credit policies.
\end{exoresolu}

\courssub{Summary of key formulas}

\begin{aretenir}[Investment Appraisal Formulas]
\begin{itemize}
    \item \textbf{Payback Period}: Year before recovery $+$ (Unrecovered cost / Cash flow in recovery year)
    \item \textbf{ARR} = $\dfrac{\text{Average Annual Accounting Profit}}{\text{Average Investment}} \times 100\%$
    \item \textbf{NPV} = $\sum \dfrac{C_t}{(1+r)^t}$; Accept if NPV > 0
    \item \textbf{IRR}: Solve $\sum \dfrac{C_t}{(1+\text{IRR})^t} = 0$; Accept if IRR > $r$
    \item \textbf{Profitability Index} = $\dfrac{\text{PV of Inflows}}{\text{Initial Outlay}}$; Accept if PI > 1
    \item \textbf{MEC Decision Rule}: Invest while MEC > $r$
\end{itemize}
\end{aretenir}

\begin{exercice}[Practice Questions]
\begin{enumerate}
    \item Distinguish between gross investment, net investment, and replacement investment. If a firm's capital stock is 500 billion FCFA and depreciation is 5\%, what is replacement investment? If gross investment is 30 billion, what is net investment?
    \item A project costs 80m FCFA, generates 25m/year for 4 years, scrap value 10m. Calculate (a) Payback Period, (b) ARR (straight-line depreciation). The firm's hurdles: max payback 3 years, min ARR 20\%. Should it accept?
    \item Explain why the MEC schedule slopes downward. Show graphically how a fall in the interest rate affects the optimal capital stock.
    \item "Payback period is a measure of risk, not profitability." Discuss.
    \item A Cameroonian SME has two projects with the following cash flows (millions FCFA):
    \begin{center}
    \begin{tabular}{|c|cccc|}\hline
    & Year 0 & Year 1 & Year 2 & Year 3 \\ \hline
    Project Alpha & -100 & 60 & 60 & 0 \\
    Project Beta  & -100 & 20 & 40 & 120 \\ \hline
    \end{tabular}
    \end{center}
    Cost of capital = 10\%. Calculate NPV and IRR for each. Which should be chosen if (a) independent, (b) mutually exclusive? Explain any conflict.
\end{enumerate}
\end{exercice}

\begin{corrige}[Exercise 1]
Replacement Investment = 5\% $\times$ 500bn = 25bn FCFA. Net Investment = Gross $-$ Depreciation = 30 $-$ 25 = 5bn FCFA.
\end{corrige}

\begin{corrige}[Exercise 2]
(a) Cumulative: 25, 50, 75, 100. Payback = 3 + (5/25) = 3.2 years $>$ 3 $\rightarrow$ Reject on payback.
(b) Depreciation = (80-10)/4 = 17.5. Avg annual profit = 25 $-$ 17.5 = 7.5. Avg investment = (80+10)/2 = 45. ARR = 7.5/45 = 16.67\% $<$ 20\% $\rightarrow$ Reject on ARR.
\end{corrige}

\begin{corrige}[Exercise 3]
MEC slopes down because of diminishing marginal returns to capital: the best projects are undertaken first. Graph: MEC curve downward, horizontal $r$ line. Lower $r$ shifts intersection rightward $\rightarrow$ higher $K^*$.
\end{corrige}

\begin{corrige}[Exercise 4]
Payback ignores time value of money and all cash flows after the cutoff. A project with a 2-year payback but zero subsequent returns is less profitable than one with a 4-year payback and 20 years of high returns. Payback favours short-term liquidity, not long-term value creation.
\end{corrige}

\begin{corrige}[Exercise 5]
Alpha NPV @10\%: -100 + 60/1.1 + 60/1.21 = -100 + 54.55 + 49.59 = \textbf{4.14m}. IRR: solve -100 + 60/(1+r) + 60/(1+r)$^{2}$ = 0 $\rightarrow$ r $\approx$ 13.1\%.
Beta NPV @10\%: -100 + 20/1.1 + 40/1.21 + 120/1.331 = -100 + 18.18 + 33.06 + 90.16 = \textbf{41.40m}. IRR $\approx$ 17.4\%.
(a) Independent: Both NPV > 0, both IRR > 10\% $\rightarrow$ Accept both.
(b) Mutually exclusive: NPV says Beta (41.4 > 4.14); IRR says Beta (17.4\% > 13.1\%). No conflict here. If conflict arose (e.g., different scales), NPV is the correct criterion for wealth maximisation.
\end{corrige}

\coursec{Practical work: net present value (NPV)}

In the previous section we examined why firms invest and introduced simple appraisal tools such as the payback period and the accounting rate of return (ARR). Those methods are easy to compute but suffer from a fundamental flaw: they ignore the \cle{time value of money}. A machine that generates 10 million FCFA over five years is not equivalent to receiving 10 million FCFA today. This section equips you with the rigorous \cle{net present value (NPV)} technique, the gold standard for investment appraisal in both the private sector and public project evaluation (e.g. at the Ministry of Economy, Planning and Regional Development — MINEPAT).

\courssub{The time value of money}

\begin{definition}[Time value of money]
Money available today is worth more than the same nominal amount in the future because today's money can be invested to earn interest, because inflation erodes purchasing power, and because future cash flows carry uncertainty. In Cameroon, with the CFA franc pegged to the euro and inflation typically around 2–3\% per year, 1,000 FCFA today can be placed in a Treasury bill (Bons du Trésor) or a term deposit at a commercial bank (e.g. Afriland First Bank, UBA Cameroon) to yield more than 1,000 FCFA next year.
\end{definition}

\begin{savaistu}[The CFA franc and the time value of money]
Because the CFA franc is pegged to the euro (1 EUR = 655.957 FCFA) and the BEAC (Banque des États de l'Afrique Centrale) sets a common policy rate for the CEMAC zone, the risk-free rate used for discounting public and private projects in Cameroon is often anchored on the BEAC's tender rate (taux d'appel d'offres) or the yield on Cameroonian Treasury bonds (OAT Cameroun). This gives a concrete, observable benchmark for the discount rate \(r\).
\end{savaistu}

\courssub{Discounting and the discount factor}

To compare cash flows occurring at different dates we \emph{discount} them to a common date (usually time 0, the present). The \cle{discount factor} for a cash flow received \(n\) years from now at an annual discount rate \(r\) is:
\[
\text{Discount factor} = \frac{1}{(1+r)^n}
\]

\begin{propriete}[Discount factor properties]
\begin{itemize}
\item The discount factor is always between 0 and 1 for \(r>0\).
\item It decreases as \(n\) increases (more distant cash flows are worth less today).
\item It decreases as \(r\) increases (a higher opportunity cost of capital reduces present values).
\end{itemize}
\end{propriete}

\begin{attention}[Common mistake: adding undiscounted cash flows]
Never add a future cash flow (e.g. 5 million FCFA in Year 3) directly to an initial outlay (e.g. –10 million FCFA today) without discounting. That would be like adding metres to kilograms. Always bring every cash flow to present value terms first.
\end{attention}

\courssub{The net present value (NPV) formula}

\begin{propriete}[Net present value formula]
For a project with an initial investment \(C_0\) (a negative cash flow at time 0) and subsequent net cash inflows \(R_t\) at the end of years \(t = 1,2,\dots,n\), the NPV at discount rate \(r\) is:
\[
\boxed{\text{NPV} = \sum_{t=1}^{n} \frac{R_t}{(1+r)^t} - C_0}
\]
If cash flows occur at the beginning of each year, adjust the exponent accordingly. In examinations, assume end-of-year cash flows unless stated otherwise.
\end{propriete}

\begin{aretenir}[Key symbols]
\begin{itemize}
\item \(C_0\) : initial capital outlay (positive number, subtracted in the formula).
\item \(R_t\) : net cash inflow during year \(t\) (after tax, after operating costs, before financing costs).
\item \(r\) : discount rate (cost of capital / required rate of return), expressed as a decimal (e.g. 12\% \(\rightarrow\) 0.12).
\item \(n\) : project life in years.
\end{itemize}
\end{aretenir}

\courssub{Step-by-step NPV computation}

\begin{methode}[5-step NPV procedure]
\begin{enumerate}
\item \textbf{List the cash flows:} Draw a time line. Record \(C_0\) at time 0 and \(R_1, R_2, \dots, R_n\) at the end of each year.
\item \textbf{Choose the discount rate \(r\):} Use the firm's weighted average cost of capital (WACC), the BEAC policy rate plus a risk premium, or the rate specified in the exam question.
\item \textbf{Compute discount factors:} For each year \(t\), calculate \(\frac{1}{(1+r)^t}\). Present factors to at least \textbf{three decimal places} for accuracy.
\item \textbf{Calculate present values:} Multiply each \(R_t\) by its discount factor to get \(PV_t\).
\item \textbf{Sum and decide:} \(\text{NPV} = \sum PV_t - C_0\). Apply the decision rule below.
\end{enumerate}
\end{methode}

\begin{propriete}[NPV decision rule]
\begin{itemize}
\item \textbf{Independent projects:} Accept if \(\text{NPV} > 0\); reject if \(\text{NPV} < 0\); indifferent if \(\text{NPV} = 0\).
\item \textbf{Mutually exclusive projects:} Choose the project with the \textbf{highest positive NPV}. (If all NPVs are negative, reject all.)
\end{itemize}
\end{propriete}

\courssub{Sensitivity of NPV to the discount rate}

Because the discount rate is an estimate (often based on WACC or a hurdle rate), NPV is sensitive to \(r\). A project with large early cash flows is less sensitive than one with large late cash flows. The \cle{NPV profile} plots NPV against \(r\); the point where the curve crosses the horizontal axis is the \cle{internal rate of return (IRR)}.

\begin{popfigure}
\begin{tikzpicture}
\begin{axis}[
    width=12cm, height=7cm,
    axis lines=middle,
    xlabel={Discount rate \(r\)}, ylabel={NPV (million FCFA)},
    xmin=0, xmax=0.30, ymin=-15, ymax=25,
    xtick={0,0.05,0.10,0.12,0.15,0.20,0.25,0.30},
    xticklabels={0\%,5\%,10\%,12\%,15\%,20\%,25\%,30\%},
    ytick={-15,-10,-5,0,5,10,15,20,25},
    grid=major, grid style={dashed, gray!40},
    clip=false
]
\addplot[domain=0:0.30, samples=100, thick, popblue] {20 - 100*x + 200*x^2 - 150*x^3};
\addplot[dashed, popmuted] coordinates {(0,0) (0.30,0)};
\addplot[mark=*, poporange] coordinates {(0.12, 8.5)};
\node[above right, poporange, font=\small] at (axis cs:0.12,8.5) {NPV at 12\% = 8.5m};
\addplot[mark=*, popgreen] coordinates {(0.185, 0)};
\node[below left, popgreen, font=\small] at (axis cs:0.185,0) {IRR \(\approx\) 18.5\%};
\end{axis}
\end{tikzpicture}
\legende{NPV profile of a typical investment project. At the firm's 12\% cost of capital NPV is positive (accept). The IRR (where NPV = 0) is about 18.5\%.}
\end{popfigure}

\courssub{Comparing NPV with payback and ARR}

\begin{center}
\begin{tabularx}{\linewidth}{|l|X|X|X|}
\hline
\rowcolor{popblueL}
\textbf{Criterion} & \textbf{NPV} & \textbf{Payback} & \textbf{ARR} \\
\hline
Time value of money & Yes & No & No \\
Risk consideration & Via discount rate & Implicit (shorter payback preferred) & No \\
Cash vs. accounting profit & Cash flows & Cash flows & Accounting profit \\
Decision rule clarity & Maximize wealth & Arbitrary cut-off & Arbitrary target rate \\
Scale of investment & Handles any scale & Biased against large projects & Biased against large projects \\
\hline
\end{tabularx}
\end{center}

\begin{attention}[Exam trap]
A project can have a short payback period and a high ARR yet a negative NPV if its cash flows occur late and the discount rate is high. Always justify your recommendation with NPV as the primary criterion; mention payback/ARR only as supplementary information.
\end{attention}

\courssub{Worked example: maize-milling machine in the West Region}

\begin{experience}[Case study: KOPA Maize Mill, Bafoussam]
KOPA Cooperative in Bafoussam (West Region) plans to buy a maize-milling machine costing \textbf{25,000,000 FCFA}. The machine has a 5-year useful life, zero scrap value, and is expected to generate the following net cash inflows (after operating costs and taxes):

\begin{center}
\begin{tabularx}{\linewidth}{|c|X|X|X|X|X|}
\hline
\rowcolor{popblueL}
Year & 1 & 2 & 3 & 4 & 5 \\
\hline
Net cash inflow (million FCFA) & 8 & 9 & 10 & 7 & 6 \\
\hline
\end{tabularx}
\end{center}

The cooperative's cost of capital is \textbf{12\%}. Compute the NPV and advise the cooperative.
\end{experience}

\begin{exoresolu}[Solution: KOPA Maize Mill NPV]
\textbf{Step 1: Cash flows}
\(C_0 = 25\) million FCFA.
\(R_1=8,\; R_2=9,\; R_3=10,\; R_4=7,\; R_5=6\) (all in million FCFA).

\textbf{Step 2: Discount rate} \(r = 12\% = 0.12\).

\textbf{Step 3: Discount factors (to 4 d.p. for safety)}
\[
\begin{aligned}
\text{Year 1: } & \frac{1}{1.12} = 0.8929 \\
\text{Year 2: } & \frac{1}{1.12^2} = 0.7972 \\
\text{Year 3: } & \frac{1}{1.12^3} = 0.7118 \\
\text{Year 4: } & \frac{1}{1.12^4} = 0.6355 \\
\text{Year 5: } & \frac{1}{1.12^5} = 0.5674
\end{aligned}
\]

\textbf{Step 4: Present values}
\begin{center}
\begin{tabular}{|c|c|c|c|c|}
\hline
\rowcolor{popblueL}
Year & Cash flow (m FCFA) & Discount factor & Present value (m FCFA) & Cumulative PV (m FCFA) \\
\hline
1 & 8 & 0.8929 & 7.1432 & 7.1432 \\
2 & 9 & 0.7972 & 7.1748 & 14.3180 \\
3 & 10 & 0.7118 & 7.1180 & 21.4360 \\
4 & 7 & 0.6355 & 4.4485 & 25.8845 \\
5 & 6 & 0.5674 & 3.4044 & 29.2889 \\
\hline
\end{tabular}
\end{center}

\textbf{Step 5: NPV}
\[
\text{NPV} = 29.2889 - 25 = \mathbf{4.2889 \text{ million FCFA}}
\]

\textbf{Decision:} NPV > 0, so the project adds value. KOPA Cooperative should \textbf{accept} the investment.

\textbf{Additional insight:} The cumulative PV exceeds the initial outlay during Year 4. The discounted payback period is between 3 and 4 years (approx. 3.5 years), longer than the simple payback (which would be ~2.9 years), illustrating the impact of discounting.
\end{exoresolu}

\begin{popfigure}
\begin{tikzpicture}
\begin{axis}[
    width=12cm, height=7cm,
    axis lines=middle,
    xlabel={Discount rate \(r\)}, ylabel={NPV (million FCFA)},
    xmin=0, xmax=0.30, ymin=-10, ymax=15,
    xtick={0,0.05,0.10,0.12,0.15,0.20,0.25,0.30},
    xticklabels={0\%,5\%,10\%,12\%,15\%,20\%,25\%,30\%},
    ytick={-10,-5,0,5,10,15},
    grid=major, grid style={dashed, gray!40},
    clip=false
]
\addplot[domain=0:0.30, samples=200, thick, popblue] 
    ({x}, {8/(1+x) + 9/(1+x)^2 + 10/(1+x)^3 + 7/(1+x)^4 + 6/(1+x)^5 - 25});
\addplot[dashed, popmuted] coordinates {(0,0) (0.30,0)};
\addplot[mark=*, poporange] coordinates {(0.12, 4.2889)};
\node[above right, poporange, font=\small] at (axis cs:0.12,4.2889) {NPV at 12\% = 4.29m};
\addplot[mark=*, popgreen] coordinates {(0.182, 0)};
\node[below left, popgreen, font=\small] at (axis cs:0.182,0) {IRR \(\approx\) 18.2\%};
\end{axis}
\end{tikzpicture}
\legende{NPV profile for the KOPA maize-milling machine. At the 12\% cost of capital NPV = 4.29 million FCFA (accept). The IRR is approximately 18.2\%.}
\end{popfigure}

\begin{exercice}[Practice: NPV of a cocoa dryer in the South-West Region]
A farmer group in Kumba considers a solar cocoa dryer costing 18,000,000 FCFA. Expected net cash inflows (million FCFA): Year 1: 5, Year 2: 6, Year 3: 6, Year 4: 5, Year 5: 4. The group's cost of capital is 10\%.
\begin{enumerate}
\item Compute the NPV (show discount factors to 4 d.p.).
\item Should they accept the project?
\item Sketch the NPV profile and estimate the IRR.
\end{enumerate}
\end{exercice}

\begin{corrige}[Exercise: cocoa dryer]
\textbf{1. Discount factors at 10\%:}
Year 1: 0.9091, Year 2: 0.8264, Year 3: 0.7513, Year 4: 0.6830, Year 5: 0.6209.
\textbf{Present values:}
Year 1: 5 × 0.9091 = 4.5455
Year 2: 6 × 0.8264 = 4.9584
Year 3: 6 × 0.7513 = 4.5078
Year 4: 5 × 0.6830 = 3.4150
Year 5: 4 × 0.6209 = 2.4836
Sum of PVs = 19.9103 million FCFA.
\textbf{NPV} = 19.9103 – 18 = \textbf{1.9103 million FCFA} (positive).
\textbf{2. Decision:} Accept (NPV > 0).
\textbf{3. NPV profile:} At 0\% NPV = 26 – 18 = 8m. At 10\% NPV = 1.91m. At 15\% NPV ≈ –0.8m. IRR ≈ 13.5\%.
\end{corrige}

\begin{aretenir}[Summary of NPV practical work]
\begin{itemize}
\item NPV discounts all future cash flows to the present using the firm's cost of capital.
\item Formula: \(\text{NPV} = \sum \frac{R_t}{(1+r)^t} - C_0\).
\item Decision rule: NPV > 0 → accept; for mutually exclusive projects, pick the highest NPV.
\item Always present discount factors to at least 3 decimal places in exams.
\item NPV is superior to payback and ARR because it accounts for the time value of money, risk (via \(r\)), and uses cash flows not accounting profits.
\item Sensitivity analysis (NPV profile) shows how the decision changes with the discount rate; the IRR is the discount rate that makes NPV = 0.
\end{itemize}
\end{aretenir}

\coursec{Equilibrium level of national income}

In the previous section we learned how firms evaluate investment projects using Net Present Value. We now shift our focus from the microeconomic decision of a single firm to the macroeconomic outcome for the whole economy. Given the consumption function, the saving function, and the level of planned investment, what determines the total level of national income? Why might the economy settle at a level where resources are unemployed? And what happens when aggregate demand pushes beyond the economy's capacity? This section provides the core Keynesian answers.

\courssub{Aggregate demand and aggregate supply in the Keynesian framework}

In the simple Keynesian model (short run, fixed prices), firms produce exactly what is demanded. There are no involuntary inventory accumulations or shortages at equilibrium. \textbf{Aggregate Supply (AS)} is simply the total output produced, which equals the total income generated: $AS = Y$. \textbf{Aggregate Demand (AD)} is the total planned expenditure on that output.

In a \textbf{two-sector closed economy} (households and firms only), aggregate demand has two components:
\begin{itemize}
    \item \textbf{Consumption (C):} Planned spending by households on goods and services.
    \item \textbf{Investment (I):} Planned spending by firms on capital goods.
\end{itemize}
Thus, $AD = C + I$. In a \textbf{three-sector model} (adding government), $AD = C + I + G$, where $G$ is government expenditure on goods and services. (We treat taxes and transfers later when we introduce the government budget).

\begin{popfigure}
\begin{tikzpicture}
\begin{axis}[
    width=12cm, height=9cm,
    axis lines=middle,
    xlabel={National Income / Output ($Y$)}, ylabel={Expenditure / Income ($E, Y$)},
    xmin=0, xmax=1000, ymin=0, ymax=1000,
    xtick={0,200,400,600,800,1000}, ytick={0,200,400,600,800,1000},
    xticklabels={0,200,400,600,800,1000}, yticklabels={0,200,400,600,800,1000},
    domain=0:1000, samples=2,
    clip=false,
    grid=major, grid style={popline!30},
]
% 45 degree line
\addplot[popdark, thick, domain=0:1000] {x} node[pos=0.85, above right, font=\small] {$Y = E$ (45° line)};
% Consumption function C = a + bY
\addplot[popblue, thick, domain=0:1000] {200 + 0.6*x} node[pos=0.7, above left, font=\small] {$C = a + bY$};
% Aggregate Demand C+I (I autonomous = 150)
\addplot[poporange, very thick, domain=0:1000] {350 + 0.6*x} node[pos=0.7, above right, font=\small] {$AD = C + I$};
% Equilibrium point
\addplot[only marks, mark=*, mark size=3, popgreen] coordinates {(875,875)};
\node[popgreen, anchor=south west, font=\small] at (875,875) {$E_0$ ($Y_e=875$)};
% Dashed lines to axes
\draw[dashed, popmuted] (875,875) -- (875,0) node[below, font=\small] {$Y_e$};
\draw[dashed, popmuted] (875,875) -- (0,875) node[left, font=\small] {$E_e$};
% Autonomous consumption point
\addplot[only marks, mark=*, mark size=2, popblue] coordinates {(0,200)};
\node[popblue, anchor=east, font=\small] at (0,200) {$a$};
% AD intercept
\addplot[only marks, mark=*, mark size=2, poporange] coordinates {(0,350)};
\node[poporange, anchor=east, font=\small] at (0,350) {$a+I$};
\end{axis}
\end{tikzpicture}
\legende{The Keynesian Cross (45° line diagram). Equilibrium occurs where the $AD = C+I$ line crosses the 45° line. At $Y_e$, planned expenditure equals actual output.}
\end{popfigure}

\courssub{The two equilibrium conditions}

\begin{definition}[Equilibrium National Income]
Equilibrium national income ($Y_e$) is the level of income at which \textbf{planned aggregate expenditure equals the actual level of output (income)}. There is no tendency for output to change because firms' sales expectations are met and inventories are stable.
\end{definition}

This definition yields two equivalent equilibrium conditions.

\begin{propriete}[Equivalence of the Income-Expenditure and Injections-Withdrawals Approaches]
In a two-sector economy with no government and no foreign trade, the following two conditions are exactly equivalent:
\begin{enumerate}
    \item \textbf{Income = Expenditure (Keynesian Cross):} $Y = C + I$
    \item \textbf{Withdrawals = Injections (S = I):} $S = I$
\end{enumerate}
\textbf{Proof (examinable):} By definition, $Y = C + S$ (income is either consumed or saved). Substituting into $Y = C + I$ gives $C + S = C + I \Rightarrow S = I$.
\end{propriete}

\begin{aretenir}[Interpretation]
\begin{itemize}
    \item \textbf{$Y = C + I$:} Focuses on the \emph{demand side}. If $Y > C+I$, firms accumulate unsold stocks $\rightarrow$ cut production. If $Y < C+I$, firms run down stocks $\rightarrow$ increase production.
    \item \textbf{$S = I$:} Focuses on the \emph{financial/circular flow}. Saving is a withdrawal from the circular flow; investment is an injection. If $S > I$, withdrawals exceed injections $\rightarrow$ income falls. If $I > S$, injections exceed withdrawals $\rightarrow$ income rises.
\end{itemize}
\end{aretenir}

\courssub{Algebraic determination of equilibrium income}

We assume a linear consumption function $C = a + bY$ (where $a > 0$ is autonomous consumption and $0 < b < 1$ is the MPC) and autonomous investment $I = \bar{I}$.

\begin{methode}[Solving for Equilibrium Income Algebraically (4 Steps)]
\begin{enumerate}
    \item \textbf{Write the equilibrium condition:} $Y = C + I$.
    \item \textbf{Substitute the functions:} $Y = a + bY + \bar{I}$.
    \item \textbf{Collect $Y$ terms on the left:} $Y - bY = a + \bar{I}$.
    \item \textbf{Solve for $Y_e$:} $Y_e = \dfrac{a + \bar{I}}{1 - b}$.
\end{enumerate}
\textbf{Note:} The denominator $(1-b)$ is the MPS. The numerator $(a+\bar{I})$ is total autonomous expenditure.
\end{methode}

\begin{exoresolu}[Finding Equilibrium Income]
\textbf{Statement:} In a hypothetical Cameroonian two-sector economy, the consumption function is $C = 50 + 0.8Y$ (in billions of CFA francs) and planned investment is $\bar{I} = 30$ billion CFA francs.
\begin{enumerate}
    \item Calculate the equilibrium level of national income $Y_e$.
    \item Calculate the equilibrium levels of Consumption ($C_e$), Saving ($S_e$), and verify $S_e = I$.
    \item Suppose investment rises to 50 billion. Find the new $Y_e$.
\end{enumerate}
\tcblower
\textbf{Detailed solution:}
\begin{enumerate}
    \item \textbf{Step 1:} $Y = C + I$.
    \textbf{Step 2:} $Y = 50 + 0.8Y + 30$.
    \textbf{Step 3:} $Y - 0.8Y = 80 \Rightarrow 0.2Y = 80$.
    \textbf{Step 4:} $Y_e = \dfrac{80}{0.2} = \mathbf{400}$ billion CFA francs.
    \item At $Y_e = 400$:
    $C_e = 50 + 0.8(400) = 50 + 320 = \mathbf{370}$ billion.
    $S_e = Y_e - C_e = 400 - 370 = \mathbf{30}$ billion.
    \textbf{Verification:} $S_e = 30 = \bar{I}$. Equilibrium holds.
    \item New $\bar{I} = 50$. Autonomous expenditure $= 50 + 50 = 100$.
    $Y_e' = \dfrac{100}{0.2} = \mathbf{500}$ billion CFA francs.
    The increase in income ($\Delta Y = 100$) is 5 times the increase in investment ($\Delta I = 20$). This is the \textbf{multiplier} at work (see next section).
\end{enumerate}
\end{exoresolu}

\courssub{Graphical determination: The two diagrams}

We have already seen the \textbf{Keynesian Cross (45° line diagram)}. The second diagram plots the \textbf{Saving and Investment functions} against income.

\begin{popfigure}
\begin{tikzpicture}
\begin{axis}[
    width=12cm, height=9cm,
    axis lines=middle,
    xlabel={National Income ($Y$)}, ylabel={Saving / Investment ($S, I$)},
    xmin=0, xmax=1000, ymin=-100, ymax=300,
    xtick={0,200,400,600,800,1000}, ytick={-100,0,100,200,300},
    xticklabels={0,200,400,600,800,1000}, yticklabels={-100,0,100,200,300},
    domain=0:1000, samples=2,
    clip=false,
    grid=major, grid style={popline!30},
]
% Saving function S = -a + (1-b)Y = -50 + 0.2Y
\addplot[popblue, thick, domain=0:1000] {-50 + 0.2*x} node[pos=0.8, above left, font=\small] {$S = -a + (1-b)Y$};
% Investment line (autonomous)
\addplot[poporange, very thick, domain=0:1000] {30} node[pos=0.5, above right, font=\small] {$I = \bar{I}$};
% Equilibrium point
\addplot[only marks, mark=*, mark size=3, popgreen] coordinates {(400,30)};
\node[popgreen, anchor=south west, font=\small] at (400,30) {$E_0$ ($Y_e=400$)};
% Dashed line to x-axis
\draw[dashed, popmuted] (400,30) -- (400,0) node[below, font=\small] {$Y_e$};
% S-intercept
\addplot[only marks, mark=*, mark size=2, popblue] coordinates {(0,-50)};
\node[popblue, anchor=north east, font=\small] at (0,-50) {$-a$};
% Break-even income (where S=0)
\addplot[only marks, mark=*, mark size=2, popmuted] coordinates {(250,0)};
\node[popmuted, anchor=north west, font=\small] at (250,0) {$Y_{break}$};
\end{axis}
\end{tikzpicture}
\legende{The Injections-Withdrawals (S-I) diagram. Equilibrium is at the intersection of the saving curve and the horizontal investment line. Note the same $Y_e = 400$ as in the Keynesian Cross.}
\end{popfigure}

\begin{attention}[Common Mistake: Confusing the Two Diagrams]
Do not mix up the axes!
\begin{itemize}
    \item \textbf{Keynesian Cross:} Vertical axis = \textbf{Expenditure ($E$)}. Lines: $E=Y$ (45°), $C$, $C+I$.
    \item \textbf{S-I Diagram:} Vertical axis = \textbf{Saving / Investment ($S, I$)}. Lines: $S$ (upward sloping), $I$ (horizontal).
\end{itemize}
Both diagrams yield the \textbf{same} $Y_e$. In the exam, you may be asked to draw either one. Label axes clearly!
\end{attention}

\courssub{Equilibrium income vs. full-employment income: The central Keynesian insight}

\begin{propriete}[Equilibrium does not imply Full Employment]
The equilibrium level of income $Y_e$ is determined solely by aggregate demand ($C+I$). The \textbf{full-employment level of income ($Y_f$)} is the maximum output the economy can produce when all available resources (labour, capital, land) are fully employed, given the current technology. \textbf{There is no automatic mechanism in the simple Keynesian model that guarantees $Y_e = Y_f$.}
\end{propriete}

\begin{aretenir}[Three Possible Situations]
\begin{enumerate}
    \item \textbf{$Y_e < Y_f$ (Deflationary / Recessionary Gap):} Aggregate demand is insufficient to buy the full-employment output. Result: \textbf{cyclical unemployment}, idle capacity, recession. \emph{Policy: Expansionary fiscal/monetary policy to raise AD.}
    \item \textbf{$Y_e = Y_f$ (Ideal):} The economy is at full employment with stable prices (in this fixed-price model). \emph{Policy: Maintain.}
    \item \textbf{$Y_e > Y_f$ (Inflationary Gap):} Aggregate demand exceeds the economy's capacity to produce at current prices. Since output cannot exceed $Y_f$ in the short run, the excess demand pulls up \textbf{prices (demand-pull inflation)}. \emph{Policy: Contractionary fiscal/monetary policy to reduce AD.}
\end{enumerate}
\end{aretenir}

\begin{definition}[Deflationary Gap]
The \textbf{deflationary gap} is the vertical distance by which the aggregate demand curve ($C+I$) lies \textbf{below} the 45° line at the full-employment income level $Y_f$. It measures the \textbf{deficiency of aggregate demand} needed to reach full employment.
\end{definition}

\begin{definition}[Inflationary Gap]
The \textbf{inflationary gap} is the vertical distance by which the aggregate demand curve ($C+I$) lies \textbf{above} the 45° line at the full-employment income level $Y_f$. It measures the \textbf{excess of aggregate demand} over the maximum possible output at current prices.
\end{definition}

\begin{attention}[Gaps are measured at $Y_f$, not at $Y_e$]
A very common exam error is to measure the gap at the equilibrium income. By definition, both the deflationary and the inflationary gap are \textbf{vertical distances measured at the full-employment income $Y_f$}, between the $AD$ curve and the 45° line.
\end{attention}

\begin{popfigure}
\begin{tikzpicture}
\begin{axis}[
    width=14cm, height=10cm,
    axis lines=middle,
    xlabel={National Income ($Y$)}, ylabel={Expenditure / Income ($E, Y$)},
    xmin=0, xmax=1000, ymin=0, ymax=1000,
    xtick={0,400,600,800,1000}, ytick={0,400,600,800,1000},
    xticklabels={0,400,600,800,1000}, yticklabels={0,400,600,800,1000},
    domain=0:1000, samples=2,
    clip=false,
    grid=major, grid style={popline!30},
]
% 45 degree line
\addplot[popdark, thick, domain=0:1000] {x};
% Full employment vertical line
\draw[poppurple, very thick, dashed] (600,0) -- (600,600) node[above, font=\small, poppurple] {$Y_f = 600$};
% AD1: Deflationary gap (C+I low)
\addplot[popblue, very thick, domain=0:1000] {200 + 0.5*x};
% AD2: Inflationary gap (C+I high)
\addplot[poporange, very thick, domain=0:1000] {300 + 0.7*x};
% Equilibrium points
\addplot[only marks, mark=*, mark size=3, popblue] coordinates {(400,400)};
\node[popblue, anchor=south east, font=\small] at (400,400) {$E_{def}$};
\addplot[only marks, mark=*, mark size=3, poporange] coordinates {(1000,1000)};
\node[poporange, anchor=south west, font=\small] at (1000,1000) {$E_{inf}$ (notional)};
% Deflationary gap arrow at Yf=600: AD1 = 200+0.5*600 = 500
\draw[popblue, <->, thick] (600,500) -- (600,600) node[midway, right, font=\small, popblue] {Deflationary Gap};
% Inflationary gap arrow at Yf=600: AD2 = 300+0.7*600 = 720
\draw[poporange, <->, thick] (600,720) -- (600,600) node[midway, right, font=\small, poporange] {Inflationary Gap};
% Labels for AD curves
\node[popblue, font=\small] at (200,350) {$AD_1 = C+I$ (low)};
\node[poporange, font=\small] at (200,500) {$AD_2 = C+I$ (high)};
\end{axis}
\end{tikzpicture}
\legende{Deflationary and Inflationary Gaps. At $Y_f=600$, $AD_1$ leaves a deflationary gap (unemployment), while $AD_2$ creates an inflationary gap (excess demand $\rightarrow$ price rises). Both gaps are measured vertically at $Y_f$.}
\end{popfigure}

\courssub{Policy implications: Closing the gaps}

The size of the gap in \emph{expenditure} is not the same as the required change in \emph{income} because of the multiplier. To close a deflationary gap of, say, 50 billion CFA francs at $Y_f$, the government does not need to increase $G$ by 50 billion. It needs to increase autonomous expenditure by $\frac{\text{Gap}}{k}$, where $k = \frac{1}{1-b}$ is the multiplier.

\begin{exemplebox}[Closing a Deflationary Gap]
Suppose $Y_f = 1000$, $C = 100 + 0.75Y$, $I = 100$. Then $AD = 200 + 0.75Y$.
At $Y_f=1000$, $AD = 200 + 750 = 950$. Deflationary gap = $1000 - 950 = 50$.
Multiplier $k = \frac{1}{1-0.75} = 4$.
To close the gap, we need $\Delta Y = 50$. Required $\Delta G$ (or $\Delta I$) $= \frac{50}{4} = 12.5$.
An increase in government spending of only 12.5 billion CFA francs raises equilibrium income by 50 billion, exactly closing the gap.
\end{exemplebox}

\begin{savaistu}[The ``Paradox of Thrift'' in the Keynesian Cross]
If households collectively decide to save more (increase MPS, decrease MPC), the saving curve rotates upward. Paradoxically, at the new equilibrium, \textbf{total saving $S_e$ remains unchanged} (it must still equal the unchanged $I$), but \textbf{national income $Y_e$ falls}. The attempt to save more leads to a recession that prevents the higher saving from materialising. This is a classic GCE essay topic: ``Explain the paradox of thrift using the $S=I$ diagram.''
\end{savaistu}

\courssub{Extension: The three-sector model (with Government)}

When government is added, the equilibrium conditions become:
\begin{align*}
\text{Income-Expenditure:} \quad & Y = C + I + G \\
\text{Injections-Withdrawals:} \quad & S + T = I + G
\end{align*}
where $T$ is net tax revenue (taxes minus transfers). The consumption function becomes $C = a + b(Y - T)$ (disposable income). The algebraic solution is:
\[ Y_e = \frac{a - bT + I + G}{1 - b} \]
The multiplier for government spending ($k_G$) and for taxes ($k_T$) differ:
\[ k_G = \frac{1}{1-b}, \quad k_T = \frac{-b}{1-b} \]
The \textbf{balanced budget multiplier} ($\Delta G = \Delta T$) is exactly 1: $Y$ rises by the amount of the increase in $G$.

\begin{exercice}[Practice: Equilibrium and Gaps]
In a closed economy with government:
$C = 40 + 0.8(Y - T)$, $I = 50$, $G = 60$, $T = 40$. Full employment income $Y_f = 800$.
\begin{enumerate}
    \item Find the equilibrium income $Y_e$.
    \item Determine whether there is an inflationary or deflationary gap, and calculate its size.
    \item By how much must $G$ change to achieve full employment? (Assume $T$ constant).
\end{enumerate}
\end{exercice}

\begin{corrige}[Exercise: Equilibrium and Gaps]
\begin{enumerate}
    \item $Y = C + I + G = 40 + 0.8(Y-40) + 50 + 60 = 40 + 0.8Y - 32 + 110 = 118 + 0.8Y$.
    $Y - 0.8Y = 118 \Rightarrow 0.2Y = 118 \Rightarrow Y_e = \mathbf{590}$.
    \item At $Y_f = 800$, $AD = 118 + 0.8(800) = 118 + 640 = 758$.
    Gap = $Y_f - AD = 800 - 758 = \mathbf{42}$ (Deflationary gap).
    \item Multiplier $k = \frac{1}{1-0.8} = 5$. Need $\Delta Y = 800 - 590 = 210$.
    $\Delta G = \frac{\Delta Y}{k} = \frac{210}{5} = \mathbf{42}$.
    (Check: New $G = 102$. New autonomous spending $= 40 - 32 + 50 + 102 = 160$. New $Y_e = 160/0.2 = 800 = Y_f$.)
\end{enumerate}
\end{corrige}

\begin{aretenir}[Key Takeaways for Section 7]
\begin{itemize}
    \item Equilibrium $Y_e$ is where $AD = Y$ (Keynesian Cross) or $S = I$ (Withdrawals = Injections). Both give the same result.
    \item Algebraically: $Y_e = \frac{\text{Autonomous Expenditure}}{1 - MPC}$.
    \item $Y_e$ can be less than, equal to, or greater than $Y_f$ (full-employment income).
    \item \textbf{Deflationary gap:} $AD < Y$ at $Y_f$ $\rightarrow$ Unemployment. Policy: Increase $G$ or cut $T$.
    \item \textbf{Inflationary gap:} $AD > Y$ at $Y_f$ $\rightarrow$ Demand-pull inflation. Policy: Decrease $G$ or raise $T$.
    \item Gaps are measured \textbf{vertically at $Y_f$}, between the $AD$ curve and the 45° line.
    \item The required policy change is smaller than the gap because of the multiplier effect.
\end{itemize}
\end{aretenir}

\coursec{Practical work: equilibrium and full-employment income}

In the previous section we established the theory of equilibrium national income: an economy is in equilibrium when \cle{aggregate demand} equals national output ($AD = Y$) or, equivalently, when planned \cle{injections} equal planned \cle{withdrawals} ($S = I$ in a two-sector economy). This corresponds to \textbf{TU 16 (Lesson 84)} and \textbf{TU 17 (Practical work 16)} of the official syllabus. Theory, however, is only half of what the GCE examiner demands. This section is a \emph{practical workshop}: you will learn to set up numerical problems from given data, to solve them by \textbf{both} methods, to determine the full-employment level of income, and to measure and interpret the \cle{deflationary gap} and the \cle{inflationary gap} (\textbf{TU 18, Further studies 3}). These calculations appear almost every year in Paper 2 (structured questions), so master the procedure step by step.

\courssub{Setting up a numerical problem}

A typical examination question gives you:
\begin{itemize}
\item a \textbf{consumption function} of the form $C = a + bY$, where $a$ is autonomous consumption and $b$ is the marginal propensity to consume (MPC);
\item a level of \textbf{planned investment} $I$, assumed autonomous (independent of income);
\item sometimes the \textbf{full-employment level of income} $Y_f$.
\end{itemize}

You are asked to find the equilibrium level of income, and then to compare it with full-employment income.

\begin{methode}[The two equivalent approaches to equilibrium]
\begin{enumerate}
\item \textbf{Aggregate demand approach.} Build a schedule with columns $Y$, $C = a + bY$, $I$, and $AD = C + I$. Equilibrium is the row where $AD = Y$. Algebraically, solve $Y = a + bY + I$, giving $Y_e = \dfrac{a + I}{1 - b}$.
\item \textbf{Injections/withdrawals approach.} Derive the saving function $S = Y - C = -a + (1-b)Y$, tabulate $S$ and $I$, and locate the row where $S = I$.
\item \textbf{Always verify} that both methods give the same answer. If they do not, an arithmetic error has been made.
\end{enumerate}
\end{methode}

\begin{attention}[Choose the right saving function]
A frequent error is to write $S = (1-b)Y$ and forget the intercept. Because $C = a + bY$, saving is $S = Y - C = -a + (1-b)Y$. The term $-a$ is \emph{dissaving} at zero income: households consume $a$ even when $Y = 0$ by drawing on past savings or borrowing.
\end{attention}

\courssub{Worked schedule: finding equilibrium by both methods}

\begin{exemplebox}[Setting up the problem]
Suppose in the economy of ``Kamerland'' the consumption function is $C = 100 + 0.8Y$ and planned investment is $I = 200$ (all figures in billions of FCFA). The full-employment level of income is $Y_f = 2000$. We seek the equilibrium income and the state of the economy relative to full employment.
\end{exemplebox}

\textbf{Step 1: build the schedule.} Choose values of $Y$ around the expected equilibrium and compute $C$, $S$, $I$ and $AD$:

\begin{center}
\begin{tabularx}{\linewidth}{|c|X|X|X|X|}
\hline
\rowcolor{popblueL} $Y$ & $C = 100 + 0.8Y$ & $S = -100 + 0.2Y$ & $I$ & $AD = C + I$ \\
\hline
$1000$ & $900$ & $100$ & $200$ & $1100$ \\
\hline
$1250$ & $1100$ & $150$ & $200$ & $1300$ \\
\hline
\rowcolor{popgreenL} $1500$ & $1300$ & $200$ & $200$ & $1500$ \\
\hline
$1750$ & $1500$ & $250$ & $200$ & $1700$ \\
\hline
$2000$ & $1700$ & $300$ & $200$ & $1900$ \\
\hline
\end{tabularx}
\end{center}

\textbf{Step 2: read off the equilibrium.} The highlighted row shows $AD = Y = 1500$ and, simultaneously, $S = I = 200$. Both methods confirm $Y_e = 1500$ billion FCFA.

\textbf{Step 3: verify algebraically.}
\begin{align*}
Y &= C + I = 100 + 0.8Y + 200\\
Y - 0.8Y &= 300\\
0.2Y &= 300 \quad\Rightarrow\quad Y_e = \frac{300}{0.2} = 1500.
\end{align*}
By the injections/withdrawals method: $S = I \Rightarrow -100 + 0.2Y = 200 \Rightarrow 0.2Y = 300 \Rightarrow Y_e = 1500$. \emph{Same answer — the two methods are equivalent.}

\begin{popfigure}
\begin{center}
\begin{tikzpicture}
\begin{axis}[
  width=11cm, height=8.5cm,
  xlabel={National income $Y$ (billions FCFA)},
  ylabel={Expenditure (billions FCFA)},
  xmin=0, xmax=2400, ymin=0, ymax=2400,
  xtick={0,500,1000,1500,2000},
  ytick={0,500,1000,1500,2000},
  grid=both, grid style={popline, very thin},
  legend style={at={(0.03,0.97)}, anchor=north west, font=\small},
  axis lines=left
]
\addplot[domain=0:2400, samples=2, thick, popmuted, dashed] {x};
\addlegendentry{$45^{\circ}$ line ($AD = Y$)}
\addplot[domain=0:2400, samples=2, very thick, popblue] {300 + 0.8*x};
\addlegendentry{$AD = C + I = 300 + 0.8Y$}
\addplot[domain=0:2400, samples=2, very thick, poporange] {400 + 0.8*x};
\addlegendentry{$AD' = AD + 100$ (gap closed)}
\addplot[mark=*, only marks, popdark, mark size=2.5pt] coordinates {(1500,1500) (2000,2000)};
\draw[thick, popgreen, dashed] (axis cs:1500,0) -- (axis cs:1500,1500);
\draw[thick, poppurple, dashed] (axis cs:2000,0) -- (axis cs:2000,2000);
\draw[<->, very thick, poppink] (axis cs:2000,1900) -- (axis cs:2000,2000);
\node[popdark, font=\small, anchor=west] at (axis cs:1520,1350) {$E$: $Y_e = 1500$};
\node[poppurple, font=\small, anchor=west] at (axis cs:2020,1750) {$Y_f = 2000$};
\node[poppink, font=\small, anchor=west] at (axis cs:2020,1950) {gap $=100$};
\end{axis}
\end{tikzpicture}
\legende{Equilibrium at $E$ where the $AD$ line crosses the $45^{\circ}$ line ($Y_e = 1500$). Full-employment income $Y_f = 2000$ requires $AD$ to rise by the vertical deflationary gap of $100$ (shift to $AD'$).}
\end{center}
\end{popfigure}

\courssub{Full-employment income and the deflationary gap}

Equilibrium income is \emph{not necessarily} full-employment income. In our example the economy settles at $Y_e = 1500$ while full employment requires $Y_f = 2000$: there is unemployment of resources. The question is: \textbf{by how much must aggregate demand rise} to push the economy to full employment?

\begin{definition}[Deflationary gap]
The \cle{deflationary gap} is the amount by which aggregate demand at the full-employment level of income falls short of the output the economy could produce at full employment. It is the \textbf{vertical distance} between the $AD$ line and the $45^{\circ}$ line \emph{at} $Y_f$, i.e.\ the extra injection of spending needed to reach full employment.
\end{definition}

\begin{propriete}[Measuring the deflationary gap]
At full-employment income $Y_f$:
\[
\text{Deflationary gap} = Y_f - AD(Y_f).
\]
In our example, $AD(2000) = 300 + 0.8(2000) = 1900$, so the gap is $2000 - 1900 = 100$ billion FCFA. Note the link with the multiplier: since $k = \dfrac{1}{1 - b} = \dfrac{1}{0.2} = 5$, an extra injection of $100$ raises income by $5 \times 100 = 500$, exactly the shortfall $Y_f - Y_e$. Hence equivalently:
\[
\text{Deflationary gap} = \frac{Y_f - Y_e}{k}.
\]
\end{propriete}

\begin{attention}[The gap is measured vertically, not horizontally]
The most common examination mistake is to report the gap as $Y_f - Y_e = 500$. That is the \emph{income shortfall} (horizontal distance). The deflationary gap is a distance \textbf{on the expenditure axis}: here $100$, not $500$. Because of the multiplier, the gap is always \emph{smaller} than the income shortfall it must eliminate. Write the units carefully: the gap is $100$ billion FCFA of \emph{expenditure}.
\end{attention}

\courssub{The inflationary gap}

The mirror-image situation arises when equilibrium income \emph{exceeds} full-employment income.

\begin{definition}[Inflationary gap]
The \cle{inflationary gap} is the amount by which aggregate demand at the full-employment level of income \emph{exceeds} full-employment output:
\[
\text{Inflationary gap} = AD(Y_f) - Y_f.
\]
Since real output cannot expand beyond $Y_f$ in the short run, the excess demand pulls prices up: it is purely inflationary.
\end{definition}

\begin{exemplebox}[An inflationary gap]
Suppose instead $C = 100 + 0.8Y$ and $I = 500$, with $Y_f = 2000$. Then $Y_e = \dfrac{100 + 500}{0.2} = 3000 > Y_f$. At full employment, $AD(2000) = 600 + 0.8(2000) = 2200$, so the inflationary gap is $2200 - 2000 = 200$ billion FCFA. Demand must be \emph{reduced} by $200$ (a fall of $5 \times 200 = 1000$ in income) to restore price stability.
\end{exemplebox}

\courssub{Policy interpretation: closing the gap}

The whole point of measuring the gap is to advise the government. With a deflationary gap of $100$ and a multiplier of $5$:
\begin{itemize}
\item \textbf{Increase government spending} by $\Delta G = 100$ billion FCFA (e.g.\ road construction, school building), which raises income by $5 \times 100 = 500$ and closes the gap exactly; or
\item \textbf{Reduce taxes} to raise disposable income and consumption. Because the tax multiplier is $k_T = -\dfrac{b}{1-b}$ (smaller in absolute value than the spending multiplier $k = \dfrac{1}{1-b}$), the required tax cut is \emph{larger} than $100$: with MPC $= 0.8$, $\Delta T = \dfrac{\text{gap}}{b} = \dfrac{100}{0.8} = 125$ billion FCFA (a reduction of $125$).
\end{itemize}
For an inflationary gap the advice is reversed: cut government spending or raise taxes (contractionary fiscal policy).

\begin{exoresolu}[GCE-style question on equilibrium and gaps]
The economy of Nkamland has the consumption function $C = 200 + 0.75Y$, planned investment $I = 300$, and full-employment income $Y_f = 2400$ (billions of FCFA). (a) Determine the equilibrium level of income using \textbf{both} the aggregate demand and the injections/withdrawals methods. (b) State whether a deflationary or inflationary gap exists and measure it. (c) Advise the government on the change in government spending required to reach full employment.
\tcblower
\textbf{(a) Aggregate demand method:}
\begin{align*}
Y &= C + I = 200 + 0.75Y + 300\\
0.25Y &= 500 \quad\Rightarrow\quad Y_e = 2000.
\end{align*}
\textbf{Injections/withdrawals method:} $S = Y - C = -200 + 0.25Y$. Setting $S = I$:
\[
-200 + 0.25Y = 300 \;\Rightarrow\; 0.25Y = 500 \;\Rightarrow\; Y_e = 2000. \checkmark
\]
Both methods give $Y_e = 2000$ billion FCFA.

\textbf{(b)} Since $Y_e = 2000 < Y_f = 2400$, a \textbf{deflationary gap} exists. At full employment, $AD(2400) = 500 + 0.75(2400) = 2300$, so:
\[
\text{Deflationary gap} = 2400 - 2300 = 100 \text{ billion FCFA}.
\]

\textbf{(c)} The multiplier is $k = \dfrac{1}{1 - 0.75} = 4$. The income shortfall is $2400 - 2000 = 400$, requiring $\Delta Y = 400 = k \times \Delta G$, hence:
\[
\Delta G = \frac{400}{4} = 100 \text{ billion FCFA}.
\]
The government should \emph{increase} its spending by $100$ billion FCFA (or cut taxes by $\frac{100}{0.75} \approx 133.3$ billion FCFA) to close the gap.
\end{exoresolu}

\begin{methode}[Exam technique: when the question says ``using the AD and injections/withdrawals methods'']
\begin{enumerate}
\item Write down $C$, derive $S = Y - C$, and state $I$.
\item Present a small \textbf{schedule table} (4--5 rows) showing $AD = Y$ and $S = I$ at the same income level.
\item Confirm each method \textbf{algebraically} — the table alone may not earn full marks.
\item Compare $Y_e$ with $Y_f$ and name the gap before measuring it.
\item Measure the gap \textbf{vertically}: gap $= Y_f - AD(Y_f)$ (deflationary) or $AD(Y_f) - Y_f$ (inflationary).
\item End with a policy recommendation using $\Delta G = \dfrac{Y_f - Y_e}{k}$.
\end{enumerate}
\end{methode}

\begin{attention}[Other traps to avoid]
\begin{itemize}
\item Do not confuse the \emph{income shortfall} ($Y_f - Y_e$) with the \emph{expenditure gap} ($Y_f - AD(Y_f)$). Examiners penalise this heavily.
\item When computing $AD$ at $Y_f$, substitute $Y_f$ into the $AD$ equation — do not simply reuse $AD$ at equilibrium.
\item A gap exists only relative to $Y_f$: if $Y_e = Y_f$, the economy is at full-employment equilibrium and no gap exists.
\item Keep all figures in the same units (billions of FCFA) throughout the answer.
\end{itemize}
\end{attention}

\begin{exercice}[Practice]
In the economy of Mungo Republic, $C = 150 + 0.8Y$, $I = 250$ (billions of FCFA) and full-employment income is $Y_f = 2500$.
\begin{enumerate}
\item Construct a schedule for $Y = 1500,\ 1750,\ 2000,\ 2250,\ 2500$ showing $C$, $S$, $I$ and $AD$, and identify the equilibrium income by both methods.
\item Verify your answer algebraically.
\item Determine whether a deflationary or inflationary gap exists and measure it.
\item Calculate the change in government spending needed to achieve full employment.
\end{enumerate}
\end{exercice}

\begin{corrige}[Exercise 1]
\textbf{1.} Schedule: at $Y=2000$: $C = 1750$, $S = 250$, $AD = 2000$. Hence $AD = Y$ and $S = I = 250$ at $Y_e = 2000$.
\textbf{2.} $Y = 150 + 0.8Y + 250 \Rightarrow 0.2Y = 400 \Rightarrow Y_e = 2000$; and $S = -150 + 0.2Y = 250 \Rightarrow Y = 2000$. $\checkmark$
\textbf{3.} $Y_e = 2000 < Y_f = 2500$: deflationary gap. $AD(2500) = 400 + 0.8(2500) = 2400$; gap $= 2500 - 2400 = 100$ billion FCFA.
\textbf{4.} $k = \frac{1}{0.2} = 5$; $\Delta Y = 500$, so $\Delta G = \frac{500}{5} = 100$ billion FCFA increase.
\end{corrige}

\begin{aretenir}[Key points]
\begin{itemize}
\item Equilibrium income is found where $AD = Y$ (aggregate demand approach) or where $S = I$ (injections/withdrawals approach); both always give the same $Y_e = \dfrac{a + I}{1 - b}$.
\item Full-employment income $Y_f$ is given by the economy's capacity; $Y_e$ may lie below or above it.
\item Deflationary gap $= Y_f - AD(Y_f)$ (extra injection needed); inflationary gap $= AD(Y_f) - Y_f$ (excess demand to remove).
\item Gaps are measured \textbf{vertically} in expenditure terms, never horizontally along the income axis.
\item Policy: $\Delta G = \dfrac{Y_f - Y_e}{k}$; tax changes must be larger because the tax multiplier is smaller in absolute value ($|k_T| = \frac{b}{1-b} < \frac{1}{1-b} = k$).
\end{itemize}
\end{aretenir}

\coursec{The multiplier}

In the previous section we determined the equilibrium level of national income and saw how an injection of investment, government spending or exports raises income. But a natural question remains: \emph{by how much} does income rise when spending increases? You may have noticed in our worked examples that an injection of 1 billion FCFA can raise national income by \emph{more} than 1 billion FCFA. This surprising result is explained by one of the most famous ideas in Keynesian economics: the \cle{multiplier}.

\courssub{The intuition: one person's spending is another person's income}

Imagine the government awards a contract of 1 billion FCFA to a construction firm in Buea to rehabilitate a road. The firm pays wages to its workers and profits to its owners: 1 billion FCFA of new income has been created. But the story does not stop there. The workers spend most of their new income on food at the market, on transport, on school fees. The traders, drivers and teachers who receive this money in turn spend most of it again. Each round of spending creates new income for someone else, which is partly re-spent, creating yet more income.

At each round, part of the income \emph{leaks out} as saving, so each successive round is smaller than the previous one. The process gradually dies out, but the \emph{total} increase in income is several times larger than the initial injection. This is the multiplier effect.

\begin{definition}[The multiplier]
The \cle{multiplier} ($k$) is the ratio of the total change in national income to the initial change in autonomous spending (investment, government spending or exports) that caused it:
\[ k = \frac{\Delta Y}{\Delta I} \qquad\text{so that}\qquad \Delta Y = k \times \Delta I. \]
\end{definition}

For example, if an increase in investment of 1 billion FCFA leads to a total increase in national income of 5 billion FCFA, then $k = 5/1 = 5$.

\courssub{The multiplier process, round by round}

Let us follow the process precisely. Suppose an investment injection of $\Delta I = 1$ billion FCFA takes place and the \cle{marginal propensity to consume} is $MPC = 0.8$ throughout the economy. This means that every household that receives an extra franc of income spends 0.8 FCFA and saves 0.2 FCFA.

\begin{itemize}
\item \textbf{Round 1:} the investment itself creates 1 billion FCFA of income (wages, profits of the construction firm).
\item \textbf{Round 2:} the recipients spend $0.8 \times 1 = 0.8$ billion FCFA, which becomes income for traders and producers.
\item \textbf{Round 3:} these spend $0.8 \times 0.8 = 0.64$ billion FCFA, and so on.
\end{itemize}

\begin{popfigure}
\begin{center}
\begin{tabular}{|c|c|c|c|}
\hline
\rowcolor{popblueL} \textbf{Round} & \textbf{New income (billion FCFA)} & \textbf{Consumed (0.8)} & \textbf{Saved (0.2)} \\
\hline
1 & 1.000 & 0.800 & 0.200 \\
\hline
2 & 0.800 & 0.640 & 0.160 \\
\hline
3 & 0.640 & 0.512 & 0.128 \\
\hline
4 & 0.512 & 0.410 & 0.102 \\
\hline
5 & 0.410 & 0.328 & 0.082 \\
\hline
$\vdots$ & $\vdots$ & $\vdots$ & $\vdots$ \\
\hline
$n \to \infty$ & $\to 0$ & $\to 0$ & $\to 0$ \\
\hline
\rowcolor{popgreenL} \textbf{Total} & $\Delta Y = 5$ & $4$ & $1$ \\
\hline
\end{tabular}
\end{center}
\legende{Round-by-round multiplier process for $\Delta I = 1$ billion FCFA and $MPC = 0.8$: the total change in income converges to $\Delta Y = k \cdot \Delta I = 5$ billion FCFA.}
\end{popfigure}

The total increase in income is the sum of all the rounds:
\[ \Delta Y = 1 + 0.8 + 0.8^2 + 0.8^3 + \dots \]

\courssub{Derivation of the multiplier formula}

The sum above is a geometric series with first term $\Delta I$ and common ratio $MPC$. Since $0 < MPC < 1$, the series converges:
\begin{align*}
\Delta Y &= \Delta I + MPC \cdot \Delta I + MPC^2 \cdot \Delta I + MPC^3 \cdot \Delta I + \dots \\
&= \Delta I \left(1 + MPC + MPC^2 + MPC^3 + \dots\right) \\
&= \Delta I \times \frac{1}{1 - MPC}
\end{align*}
Dividing both sides by $\Delta I$ gives the multiplier formula.

\begin{propriete}[The multiplier formula (proof examinable)]
\[ k = \frac{\Delta Y}{\Delta I} = \frac{1}{1 - MPC} = \frac{1}{MPS} \]
because $MPC + MPS = 1$, so $1 - MPC = MPS$. The derivation from the geometric series above is instructive and may be required in the examination.
\end{propriete}

A useful check: in equilibrium, total leakages must equal the injection. In our example, total saving is $0.2 + 0.16 + 0.128 + \dots = 1$ billion FCFA, exactly equal to the initial investment of 1 billion FCFA. Equilibrium is restored when $S = I$ again, at a higher level of income.

\begin{propriete}[The size of the multiplier]
The higher the $MPC$ (equivalently, the lower the $MPS$), the larger the multiplier. If $MPC = 0.5$, then $k = 2$; if $MPC = 0.75$, then $k = 4$; if $MPC = 0.9$, then $k = 10$. A larger fraction of income re-spent at each round means the process dies out more slowly, so the total effect is bigger.
\end{propriete}

\courssub{A graphical illustration}

Graphically, an increase in autonomous investment shifts the aggregate demand line $C + I$ upward by $\Delta I$. Equilibrium income rises by a \emph{multiple} of that vertical shift, because the consumption line is flatter than the $45^{\circ}$ line.

\begin{popfigure}
\begin{center}
\begin{tikzpicture}
\begin{axis}[
  width=11cm, height=8.5cm,
  axis lines=left,
  xlabel={National income $Y$ (billion FCFA)},
  ylabel={Aggregate demand (billion FCFA)},
  xmin=0, xmax=12, ymin=0, ymax=12,
  xtick={5,10}, ytick={5,10},
  legend style={at={(0.03,0.97)}, anchor=north west, font=\small},
  clip=false]
\addplot[domain=0:12, dashed, gray] {x};
\addlegendentry{$45^{\circ}$ line}
\addplot[domain=0:11, thick, popblue] {1 + 0.8*x};
\addlegendentry{$C + I_1$}
\addplot[domain=0:10, thick, poporange] {2 + 0.8*x};
\addlegendentry{$C + I_2$}
\addplot[mark=*, popdark] coordinates {(5,5)};
\addplot[mark=*, popdark] coordinates {(10,10)};
\draw[dashed, popmuted] (axis cs:5,0) -- (axis cs:5,5);
\draw[dashed, popmuted] (axis cs:10,0) -- (axis cs:10,10);
\draw[<->, thick, poppink] (axis cs:0.2,1) -- (axis cs:0.2,2) node[midway, left, font=\small] {$\Delta I = 1$};
\draw[<->, thick, popgreen] (axis cs:5,-0.6) -- (axis cs:10,-0.6) node[midway, below, font=\small] {$\Delta Y = 5$};
\end{axis}
\end{tikzpicture}
\end{center}
\legende{With $MPC = 0.8$, an upward shift of the $C+I$ line by $\Delta I = 1$ billion FCFA raises equilibrium income from $Y_1 = 5$ to $Y_2 = 10$: $\Delta Y = k \cdot \Delta I = 5 \times 1 = 5$ billion FCFA.}
\end{popfigure}

Notice that the horizontal distance $\Delta Y$ is much larger than the vertical shift $\Delta I$: this is the multiplier made visible.

\courssub{The multiplier in reverse}

\begin{attention}[The multiplier works in both directions]
The multiplier is not a machine that only creates prosperity. A \emph{fall} in autonomous spending also has a multiplied effect, downward. If investment falls by 1 billion FCFA with $MPC = 0.8$, national income falls by $5 \times 1 = 5$ billion FCFA, as each round of lost income reduces consumption further. This is why a collapse in cocoa export earnings or a cut in public investment can trigger a deep recession: the initial shock is amplified through the whole economy. Never assume the multiplier only applies to increases.
\end{attention}

\courssub{Extensions: government spending, tax and balanced budget multipliers}

\begin{savaistu}[Exam-useful extensions]
The same logic applies to other components of autonomous spending:
\begin{itemize}
\item \textbf{Government spending multiplier:} $k_G = \dfrac{1}{1-MPC}$, identical to the investment multiplier, because government spending enters aggregate demand directly.
\item \textbf{Tax multiplier:} $k_T = \dfrac{-MPC}{1-MPC}$. It is negative (a tax rise reduces income) and smaller in absolute value than $k_G$, because a tax cut is first partly saved before being spent.
\item \textbf{Balanced budget multiplier:} if government spending and taxes rise by the \emph{same} amount, national income still rises by exactly that amount: $k_G + k_T = 1$. Even a fully tax-financed increase in spending is expansionary.
\end{itemize}
These results are a frequent source of structured-question marks at the GCE Advanced Level.
\end{savaistu}

\courssub{Limitations of the multiplier}

In reality, the multiplier in an economy like Cameroon's is considerably smaller than the simple formula suggests. A good examination answer always evaluates:

\begin{itemize}
\item \textbf{Leakages:} besides saving, income leaks out through \emph{taxes} and \emph{imports}. When Cameroonians spend extra income on imported rice, fuel or phones, the spending creates income abroad, not at home. The true multiplier is $k = \dfrac{1}{MPS + MRT + MPM}$, where $MRT$ is the marginal rate of taxation and $MPM$ the marginal propensity to import.
\item \textbf{Inflation:} if the economy is near full employment, extra spending bids up prices rather than raising real output; the \emph{real} multiplier is then much smaller than the money multiplier.
\item \textbf{Spare capacity:} the multiplier assumes unemployed resources (idle factories, jobless workers) exist to respond to extra demand. Without spare capacity, output cannot expand.
\item \textbf{Time lags:} the rounds of re-spending take time — months or years — so the full effect is not immediate.
\item \textbf{Other constraints:} supply bottlenecks (poor roads, electricity shortages), business pessimism and crowding out of private investment by government borrowing can all weaken the process.
\end{itemize}

\courssub{Practical work: calculating and applying the multiplier}

\begin{methode}[Computing the multiplier and the change in income]
\begin{enumerate}
\item Identify the $MPC$ (or $MPS$) from the question. If only one is given, find the other using $MPC + MPS = 1$.
\item Compute the multiplier: $k = \dfrac{1}{1-MPC} = \dfrac{1}{MPS}$.
\item Compute the total change in income: $\Delta Y = k \times \Delta I$ (watch the sign of $\Delta I$).
\item If asked for the new level of income: $Y_{new} = Y_{old} + \Delta Y$.
\end{enumerate}
\end{methode}

\begin{exoresolu}[Worked example: the multiplier in action]
In an economy, the marginal propensity to consume is $0.75$ and equilibrium national income is 400 billion FCFA. Investment then rises by 20 billion FCFA. Calculate (a) the multiplier, (b) the total change in national income, (c) the new equilibrium level of income.
\tcblower
\textbf{Solution.}
\begin{enumerate}
\item $MPC = 0.75$, so $MPS = 1 - 0.75 = 0.25$.
\item Multiplier: $k = \dfrac{1}{MPS} = \dfrac{1}{0.25} = 4$.
\item Total change in income: $\Delta Y = k \times \Delta I = 4 \times 20 = 80$ billion FCFA.
\item New equilibrium income: $Y_{new} = 400 + 80 = 480$ billion FCFA.
\end{enumerate}
\end{exoresolu}

\begin{exoresolu}[Worked example: the reverse multiplier]
The government of a country cuts public investment by 10 billion FCFA. The marginal propensity to save is $0.2$. Calculate the total effect on national income.
\tcblower
\textbf{Solution.} $MPS = 0.2$, so $k = \dfrac{1}{0.2} = 5$. Then $\Delta Y = k \times \Delta I = 5 \times (-10) = -50$ billion FCFA. National income \emph{falls} by 50 billion FCFA: the cut is multiplied fivefold through reduced consumption at each round.
\end{exoresolu}

\begin{exoresolu}[Worked example: tax multiplier]
In an economy with $MPC = 0.8$, the government reduces lump-sum taxes by 5 billion FCFA. Calculate the total change in national income.
\tcblower
\textbf{Solution.}
$MPC = 0.8$, so $k_T = \dfrac{-MPC}{1-MPC} = \dfrac{-0.8}{0.2} = -4$.
$\Delta T = -5$ billion FCFA (a tax cut).
$\Delta Y = k_T \times \Delta T = (-4) \times (-5) = +20$ billion FCFA.
National income rises by 20 billion FCFA. Note that $|k_T| = 4 < k_G = 5$.
\end{exoresolu}

\begin{exoresolu}[Worked example: balanced budget multiplier]
The government increases spending by 10 billion FCFA and simultaneously raises lump-sum taxes by 10 billion FCFA. If $MPC = 0.75$, find the net change in national income.
\tcblower
\textbf{Solution.}
$k_G = \dfrac{1}{1-0.75} = 4$; $k_T = \dfrac{-0.75}{0.25} = -3$.
$\Delta Y = k_G \Delta G + k_T \Delta T = 4 \times 10 + (-3) \times 10 = 40 - 30 = 10$ billion FCFA.
Alternatively, balanced budget multiplier $= k_G + k_T = 4 - 3 = 1$, so $\Delta Y = 1 \times 10 = 10$ billion FCFA.
\end{exoresolu}

\begin{attention}[Exam tip: state MPC or MPS explicitly]
Before computing $k$, always write down the value of $MPC$ or $MPS$ given in the question, and show the conversion $MPS = 1 - MPC$ if needed. Examiners award marks for the formula $k = 1/(1-MPC)$, for correct substitution, and for the final answer with units (billion FCFA). A common mistake is to confuse $k = 1/MPS$ with $k = 1/MPC$, or to report $\Delta Y$ without adding it to the initial income when the \emph{new level} is requested.
\end{attention}

\begin{exercice}[Practice]
\begin{enumerate}
\item Calculate the multiplier when (a) $MPC = 0.5$; (b) $MPC = 0.9$; (c) $MPS = 0.4$.
\item In an economy with $MPS = 0.25$, exports of cocoa rise by 15 billion FCFA. Calculate the total change in national income.
\item Equilibrium income is 600 billion FCFA and $MPC = 0.8$. By how much must investment increase to raise income to 700 billion FCFA?
\item If $MPC = 0.6$, calculate (a) the government spending multiplier, (b) the tax multiplier, (c) the balanced budget multiplier. Verify that $k_G + k_T = 1$.
\end{enumerate}
\end{exercice}

\begin{corrige}[Exercise 1]
\begin{enumerate}
\item (a) $k = 1/(1-0.5) = 2$; (b) $k = 1/(1-0.9) = 10$; (c) $k = 1/0.4 = 2.5$.
\item $k = 1/0.25 = 4$; $\Delta Y = 4 \times 15 = 60$ billion FCFA.
\item Required $\Delta Y = 100$ billion FCFA; $k = 1/(1-0.8) = 5$; $\Delta I = \Delta Y / k = 100/5 = 20$ billion FCFA.
\item (a) $k_G = 1/(1-0.6) = 2.5$; (b) $k_T = -0.6/0.4 = -1.5$; (c) $k_G + k_T = 2.5 - 1.5 = 1$.
\end{enumerate}
\end{corrige}

\begin{aretenir}[Key points]
\begin{itemize}
\item The multiplier measures the amplified effect of a change in autonomous spending on national income: $k = \Delta Y / \Delta I$.
\item $k = \dfrac{1}{1-MPC} = \dfrac{1}{MPS}$; the higher the $MPC$, the larger the multiplier.
\item The process works round by round: spending becomes income, which is partly re-spent, until leakages absorb the whole injection.
\item The multiplier works in reverse: cuts in spending multiply downward.
\item Extensions: $k_G = \frac{1}{1-MPC}$, $k_T = \frac{-MPC}{1-MPC}$, balanced budget multiplier $= 1$.
\item In practice the multiplier is reduced by taxes, imports, inflation, capacity limits and time lags: $k = \frac{1}{MPS + MRT + MPM}$.
\end{itemize}
\end{aretenir}

In the next section we shall see that investment itself responds to changes in income: this is the \emph{accelerator}, whose interaction with the multiplier can generate cumulative booms and slumps.

\coursec{The accelerator and the multiplier--accelerator interaction}

In the previous section we studied the \cle{multiplier}: an initial injection of spending (for example autonomous investment) raises national income by a multiple of itself, because one person's spending becomes another person's income, part of which is spent again. The multiplier therefore explains how \emph{investment affects income}. But the relationship between investment and income runs in both directions: the \emph{growth of income} also affects investment. This reverse link is captured by the \cle{accelerator principle}. When the multiplier and the accelerator operate together, they can generate the recurring booms and slumps known as \emph{trade cycles} (business cycles). This section explains the accelerator, shows how to perform accelerator calculations, and then combines it with the multiplier.

\courssub{The accelerator principle}

\textbf{Intuition.} Imagine a bakery in Douala whose ovens can produce at most $10\,000$ loaves per week. As long as demand is $8\,000$ loaves and stable, the bakery has no reason to buy new ovens; it only replaces worn-out machines (\emph{replacement investment}). But suppose demand rises from $8\,000$ to $9\,000$ loaves, and then to $10\,000$. To meet the \emph{growing} demand, the bakery must buy additional ovens: this is \emph{net investment}. Now notice something crucial: if demand stops growing and stays at $10\,000$ loaves, the bakery no longer needs extra ovens --- net investment falls back to \emph{zero}, even though demand is at its highest level ever. What matters for investment is not how \emph{high} demand is, but how \emph{fast it is growing}. This is the heart of the accelerator.

\begin{definition}[The accelerator principle]
The \cle{accelerator principle} states that the level of \cle{net investment} depends on the \emph{rate of change} of national income (or output), not on its level. When income is rising, firms need additional capital goods to produce the extra output, so net investment is positive; when income grows more slowly, net investment falls; when income stops growing, net investment collapses to zero (only replacement investment remains).
\end{definition}

\begin{definition}[Net investment and gross investment]
\cle{Net investment} is the change in the economy's capital stock: $I^{\text{net}} = \Delta K$. \cle{Gross investment} is total spending on capital goods, which equals net investment plus \cle{replacement investment} (depreciation): $I^{\text{gross}} = I^{\text{net}} + D$, where $D$ is the capital consumption allowance.
\end{definition}

To make the accelerator precise, we assume firms wish to maintain a fixed relationship between their capital stock and the output they produce.

\begin{definition}[Capital--output ratio]
The \cle{capital--output ratio}, denoted $v$, is the amount of capital required to produce one unit of output per period:
\[
v=\frac{K}{Y}.
\]
For example, if $v=3$, then \SI{3}{}~FCFA of capital equipment (machines, buildings) is needed to produce \SI{1}{}~FCFA of output per year. It is also called the \cle{accelerator coefficient}.
\end{definition}

If the desired capital stock is $K_t=vY_t$, then the \emph{change} in the desired capital stock between two periods is $K_t-K_{t-1}=v(Y_t-Y_{t-1})$. Since net investment is precisely the change in the capital stock, we obtain the accelerator formula.

\begin{propriete}[The accelerator formula]
\[
\boxed{\,I_t=v\,\Delta Y_t=v\left(Y_t-Y_{t-1}\right)\,}
\]
where $I_t$ is \emph{net} investment in period $t$, $v$ is the capital--output ratio (accelerator coefficient) and $\Delta Y_t$ is the change in income (output) between period $t-1$ and period $t$. Gross investment equals net investment plus replacement (depreciation) investment. The derivation follows directly from $K_t=vY_t$ and is examinable.
\end{propriete}

\begin{attention}[Change, not level!]
The most common examination mistake is to compute investment as $I=vY$ instead of $I=v\Delta Y$. The accelerator responds to the \textbf{change} in income, not to its \textbf{level}. If income is high but constant, $\Delta Y=0$ and net investment is zero. If income growth merely \emph{slows down} (income still rising, but by less than before), net investment \emph{falls}. Always compute $\Delta Y_t=Y_t-Y_{t-1}$ first.
\end{attention}

\courssub{Autonomous and induced investment}

The accelerator forces us to distinguish two types of investment:

\begin{definition}[Autonomous and induced investment]
\begin{itemize}
\item \cle{Autonomous investment} is investment that is \emph{independent of the current level or change of national income}. It is determined by factors such as new inventions, population growth, government development projects or expectations about the distant future. Example: the State building the Memve'ele hydroelectric dam as part of a long-term energy strategy.
\item \cle{Induced investment} is investment \emph{caused by changes in income or demand}. It is the investment described by the accelerator: $I_t=v\Delta Y_t$. Example: a cocoa-processing firm in Ebolowa buying extra grinders because demand for its cocoa butter is rising.
\end{itemize}
\end{definition}

In the simple multiplier model of the previous sections, all investment was treated as autonomous (a constant, $\bar{I}$). The accelerator adds induced investment, making investment a \emph{function of the change in income}.

\courssub{Practical work: calculations on the acceleration principle}

\begin{methode}[Four-step procedure for accelerator calculations]
\begin{enumerate}
\item \textbf{Compute the change in output} for each period: $\Delta Y_t=Y_t-Y_{t-1}$.
\item \textbf{Compute the required capital stock}: $K_t=vY_t$.
\item \textbf{Compute net investment}: $I_t=v\Delta Y_t$ (equivalently $K_t-K_{t-1}$). If $\Delta Y_t<0$, net investment is negative (firms let capital wear out).
\item \textbf{Add replacement investment} (if given) to obtain gross investment: $I^{\text{gross}}_t=I_t+D$, where $D$ is depreciation per period.
\end{enumerate}
\end{methode}

\begin{exoresolu}[Accelerator calculation with a capital--output ratio of 3]
A firm requires \SI{3}{}~FCFA of capital to produce \SI{1}{}~FCFA of output per year ($v=3$). Replacement investment is $20$ billion FCFA per year. Output (in billions of FCFA) over six years is: $Y_0=100$, $Y_1=110$, $Y_2=125$, $Y_3=135$, $Y_4=140$, $Y_5=140$. Calculate net investment and gross investment for each year.
\tcblower
\textbf{Step 1: changes in output.} $\Delta Y_1=10$; $\Delta Y_2=15$; $\Delta Y_3=10$; $\Delta Y_4=5$; $\Delta Y_5=0$.

\textbf{Steps 2 and 3: required capital and net investment} ($I_t=3\Delta Y_t$):

\begin{center}
\begin{tabular}{|c|c|c|c|c|c|}
\hline
\rowcolor{popblueL} Year & $Y_t$ & $\Delta Y_t$ & $K_t=3Y_t$ & Net $I_t=3\Delta Y_t$ & Gross $I_t=I_t+20$ \\
\hline
0 & 100 & -- & 300 & -- & -- \\
\hline
1 & 110 & 10 & 330 & 30 & 50 \\
\hline
2 & 125 & 15 & 375 & 45 & 65 \\
\hline
3 & 135 & 10 & 405 & 30 & 50 \\
\hline
4 & 140 & 5 & 420 & 15 & 35 \\
\hline
5 & 140 & 0 & 420 & 0 & 20 \\
\hline
\end{tabular}
\end{center}

\textbf{Interpretation.} Output rises \emph{throughout} years 1 to 4, yet net investment already \emph{falls} from year 3 onwards because output is growing more slowly ($15$, then $10$, then $5$). In year 5 output is at its highest level ($140$) but is no longer growing, so net investment collapses to zero and gross investment falls to replacement only ($20$). This illustrates perfectly why the accelerator magnifies fluctuations: a mere \emph{slowdown in the growth of demand} --- not even a fall in demand --- causes investment to collapse.
\end{exoresolu}

\begin{attention}[Negative net investment]
If output falls, $\Delta Y<0$ and the formula gives negative net investment: firms do not even replace all their worn-out machines, so the capital stock shrinks. In practice, gross investment cannot fall below zero (you cannot ``un-buy'' a machine), so the accelerator works asymmetrically on the way down: net disinvestment is limited by the amount of depreciation.
\end{attention}

\courssub{Assumptions and limitations of the accelerator}

The simple accelerator is a powerful but mechanical model. It rests on strong assumptions:

\begin{itemize}
\item \textbf{No spare capacity.} Firms are assumed to be producing at full capacity, so any extra demand requires new capital. If firms have idle machines, they can raise output \emph{without} investing, and the accelerator does not operate.
\item \textbf{Constant capital--output ratio.} Technology is assumed fixed. In reality, technical progress and changes in relative factor prices alter $v$ over time.
\item \textbf{No role for expectations.} Firms are assumed to react mechanically to \emph{past} changes in output. In reality, investment is forward-looking: a firm invests only if it \emph{expects} the rise in demand to last. A temporary blip in sales will not trigger the purchase of a new factory.
\item \textbf{Instantaneous adjustment.} The model assumes capital can be installed immediately; in reality there are long gestation lags (a dam or a cement plant takes years to build). A more realistic \emph{lagged accelerator} would write $I_t=v(Y_{t-1}-Y_{t-2})$.
\item \textbf{No financial constraints.} Firms are assumed able to finance any desired investment, ignoring interest rates, credit rationing and retained profits.
\end{itemize}

Despite these limitations, the accelerator captures a real and important mechanism: investment is the most \emph{volatile} component of aggregate demand precisely because it is tied to the \emph{growth rate} of demand rather than its level.

\courssub{Further studies: the multiplier--accelerator interaction and the trade cycle}

We now combine the two mechanisms studied in this chapter:

\begin{itemize}
\item \textbf{Multiplier:} investment $\rightarrow$ income. An injection of investment raises income by a multiple: $\Delta Y=k\Delta I$ with $k=\dfrac{1}{1-MPC}$ (in a simple closed economy with no government).
\item \textbf{Accelerator:} income $\rightarrow$ investment. A rise in income induces further investment: $I_t=v\Delta Y_t$.
\end{itemize}

\begin{aretenir}[Multiplier--accelerator interaction as an explanation of business cycles]
An initial autonomous injection (say, government spending on road construction) raises income through the \emph{multiplier}. The \emph{rising income} then induces new investment through the \emph{accelerator}. This induced investment is itself multiplied into still higher income, which induces still more investment: the economy enters a cumulative \textbf{boom}. But the boom contains the seeds of its own end: as income approaches full-employment capacity, its \emph{growth} slows down. The moment income grows more slowly, the accelerator makes induced investment \emph{fall} --- even though income is still rising. Falling investment is then multiplied into falling income, which makes investment fall further: a cumulative \textbf{recession}. When income has fallen enough, replacement needs and new autonomous injections restart the process. The interaction of the multiplier and the accelerator therefore generates self-sustaining \cle{trade cycles}: recurring fluctuations of income around its long-run trend, with phases of expansion, peak, contraction and trough.
\end{aretenir}

\begin{popfigure}
\begin{center}
\begin{tikzpicture}
\begin{axis}[
  width=0.85\linewidth, height=6.5cm,
  axis lines=middle,
  xlabel={Time (years)}, ylabel={National income $Y$},
  xlabel style={anchor=west}, ylabel style={anchor=south},
  xmin=0, xmax=13, ymin=-1.2, ymax=2.6,
  xtick=\empty, ytick=\empty,
  domain=0:12.5, samples=200, smooth]
\addplot[popmuted, dashed, thick] {0.15*x};
\addplot[popblue, very thick] {0.15*x + 0.9*sin(deg(1.05*x))*exp(-0.02*x)};
\node[popmuted, font=\small, align=center] at (axis cs:11.2,2.15) {Long-run\\trend};
\node[popblue, font=\small] at (axis cs:1.6,1.35) {Boom};
\node[popblue, font=\small] at (axis cs:4.6,0.35) {Recession};
\node[popblue, font=\small] at (axis cs:7.8,1.7) {Recovery};
\node[popblue, font=\small] at (axis cs:10.9,0.75) {Slowdown};
\end{axis}
\end{tikzpicture}
\end{center}
\legende{Income oscillating around its trend under the combined multiplier--accelerator interaction: the trade cycle.}
\end{popfigure}

The logic of the cycle can be summarised in a chain:
\[
\text{Injection} \xrightarrow{\ \text{multiplier}\ } \Delta Y \uparrow \xrightarrow{\ \text{accelerator}\ } I \uparrow \xrightarrow{\ \text{multiplier}\ } \Delta Y \uparrow\uparrow \ \cdots
\]
and, symmetrically, on the way down: slower growth of $Y$ $\Rightarrow$ collapse of induced $I$ $\Rightarrow$ multiplied fall of $Y$ $\Rightarrow$ further disinvestment.

\begin{savaistu}[Cocoa, oil and the amplified cycle in Cameroon]
Cameroon's economy is highly exposed to world prices of a few export commodities, notably cocoa, coffee, crude oil, cotton and timber. A rise in the world price of cocoa raises export earnings: this injection is multiplied through the economy as farmers, transporters and traders spend their higher incomes. Rising incomes then induce investment --- farmers plant more trees, firms buy lorries and processing equipment (the accelerator). When world prices fall, export earnings drop, the multiplier works in reverse, and because income growth slows, induced investment collapses: the downturn is \emph{amplified} far beyond the initial price shock. This is why economists recommend that commodity-dependent economies like those of the CEMAC zone \emph{diversify} their production and build stabilisation funds during booms, so as to dampen the multiplier--accelerator swings.
\end{savaistu}

\begin{exercice}[Applying the accelerator]
A country has a capital--output ratio of $v=2.5$ and annual replacement investment of $10$ billion FCFA. National income (billions of FCFA) is: $Y_0=200$, $Y_1=220$, $Y_2=236$, $Y_3=248$, $Y_4=252$, $Y_5=252$.
\begin{enumerate}
\item Calculate net investment and gross investment for years 1 to 5.
\item In which year does net investment reach its maximum? Explain why it is not the year with the highest income.
\item What happens to gross investment in year 5? Explain using the accelerator principle.
\end{enumerate}
\end{exercice}

\begin{corrige}[Exercise 1]
\textbf{1.} Using $I_t=2.5\,\Delta Y_t$ and gross investment $=I_t+10$:
\begin{center}
\begin{tabular}{|c|c|c|c|c|}
\hline
\rowcolor{popblueL} Year & $Y_t$ & $\Delta Y_t$ & Net $I_t$ & Gross $I_t$ \\
\hline
1 & 220 & 20 & 50 & 60 \\
\hline
2 & 236 & 16 & 40 & 50 \\
\hline
3 & 248 & 12 & 30 & 40 \\
\hline
4 & 252 & 4 & 10 & 20 \\
\hline
5 & 252 & 0 & 0 & 10 \\
\hline
\end{tabular}
\end{center}
\textbf{2.} Net investment is maximal in year 1 ($50$ billion), because that is when income grows \emph{fastest} ($\Delta Y=20$). The accelerator responds to the rate of change of income, not its level: income is highest in years 4--5, but investment is then at its lowest.
\textbf{3.} In year 5 income stops growing ($\Delta Y=0$), so net investment collapses to zero and gross investment falls to replacement investment only ($10$ billion). A mere stop in the \emph{growth} of income --- not a fall in income --- is enough to cause the collapse of investment, which the multiplier will then transmit to the whole economy.
\end{corrige}

\begin{aretenir}[Key points of the section]
\begin{itemize}
\item The \cle{accelerator} states that net investment depends on the \emph{change} in income: $I_t=v\Delta Y_t$, where $v$ is the capital--output ratio.
\item \emph{Induced} investment is explained by the accelerator; \emph{autonomous} investment is independent of income changes.
\item The accelerator magnifies fluctuations: a slowdown in the \emph{growth} of demand can make investment collapse, even while demand is still rising.
\item Its limitations: assumes no spare capacity, a constant $v$, no expectations, instant adjustment, no financial constraints.
\item Combined with the multiplier ($\Delta Y=k\Delta I$), the accelerator explains \cle{trade cycles}: injections are multiplied, rising income induces investment, and the mutual reinforcement generates booms and recessions --- a mechanism clearly visible in commodity-dependent economies such as Cameroon.
\end{itemize}
\end{aretenir}

\newpage

\coursec{Integration activity}

\courssub{Integration Activity: The Buea Agro-Processing Project}

This integration activity closes the chapter by bringing together, in \textbf{one single decision-making situation}, the four core skills you have acquired:

\begin{itemize}
\item \textbf{Investment appraisal}: computing the \cle{Net Present Value} (NPV) and the \cle{payback period} of a project and formulating an invest / do-not-invest recommendation;
\item \textbf{Propensity calculations}: deriving the \cle{MPC}, \cle{MPS}, \cle{APC} and \cle{APS} from household data and building the consumption function;
\item \textbf{Income determination}: finding the \cle{equilibrium level of income} by both the aggregate demand ($Y = C + I$) method and the injections--withdrawals ($S = I$) method;
\item \textbf{Multiplier analysis}: measuring the total effect of an investment injection on national (or local) income and identifying a possible \cle{deflationary gap} or \cle{inflationary gap} relative to full-employment income.
\end{itemize}

Beyond the calculations, the activity trains you in the \emph{professional} dimension of economics: writing a short, argued \textbf{advisory note} and defending it orally, exactly as an economist working for a cooperative, a local council or the Ministry of Agriculture would do.

\begin{methode}[Organising the group work]
\begin{enumerate}
\item Form groups of 4 to 5 students and appoint a coordinator, a calculator-checker and a presenter.
\item Read the situation carefully and identify the four blocks of data (cash flows, discount rate, household survey, macroeconomic data).
\item Answer the tasks in order: each answer feeds the next one.
\item Prepare a one-page advisory note (calculations + recommendation) and a 5-minute presentation.
\end{enumerate}
\end{methode}

\courssub{Situation}

\textbf{The Buea Agro-Processing Project.} The \emph{Fako Cassava Farmers' Cooperative} (FCFC), based in Buea, South West Region, groups 240 smallholder farmers. Every year the cooperative loses a large part of its cassava harvest because fresh cassava rots within a few days and must be sold cheaply and quickly. The management committee is considering the construction of a \textbf{cassava-processing plant} (production of garri, cassava flour and starch) costing \textbf{20 million FCFA}, financed by a loan and the members' contributions.

Before signing anything, the committee hires your group of young economists to answer two questions: \emph{Will the project pay?} and \emph{What will it do to the income of the Buea local economy?}

\textbf{Document 1 -- Projected net cash flows of the plant} (in million FCFA). The appropriate discount rate, given the cost of the loan and the risk, is $10\%$.

\begin{center}
\begin{tabularx}{\linewidth}{|l|X|X|X|X|X|X|}
\hline
\rowcolor{popblueL} Year & 0 & 1 & 2 & 3 & 4 & 5 \\
\hline
Net cash flow (million FCFA) & $-20$ & $4$ & $6$ & $7$ & $7$ & $6$ \\
\hline
\end{tabularx}
\end{center}

\textbf{Document 2 -- Survey of a representative household in Buea.} When the household's annual disposable income rose from \textbf{500 million FCFA} to \textbf{600 million FCFA}, its annual consumption rose from \textbf{420 million FCFA} to \textbf{486 million FCFA}. The committee assumes the whole local economy behaves like this representative agent, so that the local consumption function is linear: $C = a + bY$.

\begin{attention}[Units consistency]
In macroeconomic models we work in \textbf{million FCFA}. The household survey figures are therefore expressed in million FCFA (500, 600, 420, 486) to be directly compatible with the investment and income data in Documents 1 and 3. This is a standard \emph{representative agent} simplification.
\end{attention}

\textbf{Document 3 -- The local economy (figures in million FCFA).} The local economy is simplified to two sectors (households and firms): no government, no foreign trade. The planned investment injection is precisely the plant: $I = 20$. The full-employment level of income of the local economy is estimated at $Y_f = 350$.

\courssub{Tasks}

\begin{enumerate}
\item \textbf{Investment appraisal.} Using a discount rate of $10\%$ (discount factors: year 1: $0.909$; year 2: $0.826$; year 3: $0.751$; year 4: $0.683$; year 5: $0.621$), compute the present value of each cash flow and the NPV of the project. Compute the payback period. Should the cooperative invest? Justify.
\item \textbf{Propensities.} From Document 2, calculate the MPC, MPS, and the APC and APS at the initial income level. Deduce the consumption function $C = a + bY$ of the local economy (work in million FCFA).
\item \textbf{Multiplier.} Calculate the investment multiplier $k$. Estimate the \emph{total} change in local income generated by the 20 million FCFA injection, and explain in two or three sentences the mechanism behind it.
\item \textbf{Equilibrium income.} Determine the equilibrium level of income $Y_e$ (a) by the aggregate demand method $Y = C + I$; (b) by the withdrawals--injections method $S = I$. Verify that both methods give the same result and check the consistency with your answer to Task 3.
\item \textbf{Full employment.} Compare $Y_e$ with the full-employment income $Y_f = 350$. Is there an inflationary or a deflationary gap? Calculate the size of the gap (the additional injection needed to reach full employment) and propose one realistic way for the cooperative or the council to close it.
\item \textbf{Advisory note.} Write your one-page note: recommendation on the project, expected impact on local income, and remaining macroeconomic problem. Defend it in 5 minutes.
\end{enumerate}

\courssub{Model Answer}

\textbf{Task 1 -- Investment appraisal.}

\begin{center}
\begin{tabularx}{\linewidth}{|l|X|X|X|}
\hline
\rowcolor{popblueL} Year & Cash flow (million FCFA) & Discount factor ($10\%$) & Present value (million FCFA) \\
\hline
0 & $-20$ & 1.000 & $-20.000$ \\
1 & 4 & 0.909 & 3.636 \\
2 & 6 & 0.826 & 4.956 \\
3 & 7 & 0.751 & 5.257 \\
4 & 7 & 0.683 & 4.781 \\
5 & 6 & 0.621 & 3.726 \\
\hline
\multicolumn{3}{|l|}{Total present value of inflows (Years 1--5)} & 22.356 \\
\hline
\end{tabularx}
\end{center}

\[
\text{NPV} = \sum_{t=0}^{5} \frac{CF_t}{(1+r)^t} = -20 + 22.356 = \mathbf{+2.356 \text{ million FCFA}} > 0.
\]

Since the NPV is \textbf{positive}, the project creates value in present-value terms: the cooperative should invest.

\textbf{Payback period} (non-discounted): cumulative cash flows are $4$; $10$; $17$; $24$ (million FCFA). The initial 20 million are recovered during Year 4:
\[
\text{Payback} = 3 + \frac{20 - 17}{7} = 3 + \frac{3}{7} \approx \mathbf{3.43 \text{ years}} \quad (\text{about 3 years and 5 months}).
\]
This is acceptable for an agro-industrial plant.

\begin{attention}[Common mistake]
Do not compare the undiscounted sum of inflows ($4+6+7+7+6 = 30$) directly with the cost. Money received in Year 5 is worth less than money today: discounting is the whole point of the NPV criterion. The payback period ignores the time value of money; NPV does not.
\end{attention}

\textbf{Task 2 -- Propensities and consumption function.}

From Document 2:
\[
\Delta Y = 600 - 500 = 100 \text{ million FCFA}, \qquad
\Delta C = 486 - 420 = 66 \text{ million FCFA}.
\]

\[
\text{MPC} = \frac{\Delta C}{\Delta Y} = \frac{66}{100} = \mathbf{0.66}, \qquad
\text{MPS} = 1 - \text{MPC} = \mathbf{0.34}.
\]

At the initial income level ($Y_0 = 500$):
\[
\text{APC} = \frac{C_0}{Y_0} = \frac{420}{500} = \mathbf{0.84}, \qquad
\text{APS} = 1 - \text{APC} = \mathbf{0.16}.
\]

\begin{aretenir}[Interpretation]
$\text{APC} > \text{MPC}$ indicates the presence of \textbf{autonomous consumption} ($a > 0$): even if income were zero, households would still consume by dissaving or borrowing.
\end{aretenir}

The consumption function is $C = a + bY$ with $b = \text{MPC} = 0.66$. Using the initial point $(Y_0, C_0) = (500, 420)$:
\[
420 = a + 0.66 \times 500 \implies a = 420 - 330 = \mathbf{90}.
\]
Hence, for the whole local economy (in million FCFA):
\[
\boxed{C = 90 + 0.66Y}
\]

\textbf{Task 3 -- The investment multiplier.}

\[
k = \frac{1}{1 - \text{MPC}} = \frac{1}{\text{MPS}} = \frac{1}{0.34} \approx \mathbf{2.94}.
\]

Total impact of the 20 million FCFA injection:
\[
\Delta Y = k \times \Delta I = 2.94 \times 20 \approx \mathbf{58.8 \text{ million FCFA}}.
\]

\begin{experience}[The multiplier mechanism in words]
The 20 million spent on the plant become income for builders, transporters and input suppliers in Buea. They spend $66\%$ (their MPC) of this new income — about 13.2 million — on garri, rent, school fees, etc. The recipients of that 13.2 million in turn spend $66\%$ of it (8.7 million), and so on. The infinite geometric series $20 + 20(0.66) + 20(0.66)^2 + \cdots$ sums to $20 / (1-0.66) = 58.8$ million FCFA.
\end{experience}

\textbf{Task 4 -- Equilibrium level of income.}

\paragraph{(a) Aggregate demand method ($Y = C + I$):}
\[
Y = 90 + 0.66Y + 20 \implies Y - 0.66Y = 110 \implies 0.34Y = 110
\]
\[
Y_e = \frac{110}{0.34} \approx \mathbf{323.5 \text{ million FCFA}}.
\]

\paragraph{(b) Withdrawals--injections method ($S = I$):}
Saving function: $S = Y - C = Y - (90 + 0.66Y) = -90 + 0.34Y$.
Set $S = I = 20$:
\[
-90 + 0.34Y = 20 \implies 0.34Y = 110 \implies Y_e = \frac{110}{0.34} \approx \mathbf{323.5 \text{ million FCFA}}.
\]

\begin{aretenir}[Verification]
Both methods yield the same $Y_e = 323.5$ million FCFA. \checkmark
\end{aretenir}

\paragraph{Consistency check with Task 3:}
Without the plant ($I = 0$), equilibrium would be $Y_0 = \frac{90}{0.34} \approx 264.7$ million FCFA.
The rise due to the project is $323.5 - 264.7 = 58.8$ million FCFA, exactly the multiplier result $\Delta Y = k \Delta I$.

\textbf{Task 5 -- Full employment and the deflationary gap.}

\[
Y_e \approx 323.5 < Y_f = 350.
\]
The economy settles \emph{below} full employment: there is a \textbf{deflationary gap} (recessionary gap).

Income shortfall: $\Delta Y_{\text{gap}} = Y_f - Y_e = 350 - 323.5 = 26.5$ million FCFA.
Additional injection required to close the gap:
\[
\Delta I_{\text{gap}} = \frac{\Delta Y_{\text{gap}}}{k} = \Delta Y_{\text{gap}} \times \text{MPS} = 26.5 \times 0.34 \approx \mathbf{9.0 \text{ million FCFA}}.
\]

\begin{savaistu}[Policy option for Buea]
A realistic complementary injection: the Buea Council could finance a \textbf{rural feeder road} (approx. 9 million FCFA) linking the cooperative's farms to the plant, reducing transport costs and post-harvest losses. This public investment would, via the multiplier, lift local income to the full-employment level of 350 million FCFA.
\end{savaistu}

\textbf{Task 6 -- Advisory note (summary for the FCFC Management Committee).}

\begin{center}
\begin{tabularx}{\linewidth}{|X|}
\hline
\rowcolor{popblueL} \textbf{ADVISORY NOTE -- Buea Agro-Processing Project} \\
\hline
\textbf{1. Project viability:} \textbf{INVEST.} NPV = +2.36 million FCFA at 10\% discount rate. Payback period $\approx$ 3.4 years. The project is financially sound. \\
\textbf{2. Macroeconomic impact:} The 20 million FCFA injection will raise local equilibrium income by 58.8 million FCFA (multiplier $k \approx 2.94$), from 264.7 to 323.5 million FCFA. \\
\textbf{3. Remaining gap:} A \textbf{deflationary gap} of 26.5 million FCFA persists (full employment = 350 million). An additional 9 million FCFA of investment (e.g., a council-funded feeder road or a second processing line) would close the gap. \\
\textbf{4. Recommendation:} Approve the plant construction \emph{and} lobby the Buea Council for the complementary infrastructure investment to achieve full employment. \\
\hline
\end{tabularx}
\end{center}

\begin{aretenir}[What this activity demonstrates for the GCE]
One single real decision mobilises the entire chapter:
\begin{itemize}
\item NPV and payback judge the \emph{project} (micro/finance);
\item MPC/MPS describe \emph{behaviour} (consumption/saving functions);
\item The multiplier translates an injection into \emph{income} (income determination);
\item The comparison $Y_e$ vs $Y_f$ reveals the \emph{macroeconomic gap} (policy).
\end{itemize}
In the GCE examination, structured questions often chain exactly these steps: master the chain, not just each link.
\end{aretenir}

\newpage

\coursec{Methods and worked examples}

\courssub{Skill 1: Calculating APC, MPC, APS and MPS}

\begin{methode}[Calculating the average and marginal propensities]
To compute the propensities to consume and to save from data on income ($Y$) and consumption ($C$):
\begin{enumerate}
\item Check the units (usually billions of FCFA) and identify income $Y$ and consumption $C$ for each period.
\item Compute saving if needed: $S = Y - C$ (in a simple economy with no government and no foreign trade).
\item Average propensities are \emph{ratios at one point}:
\[ APC = \frac{C}{Y}, \qquad APS = \frac{S}{Y} = 1 - APC. \]
\item Marginal propensities are \emph{ratios of changes between two periods}:
\[ MPC = \frac{\Delta C}{\Delta Y}, \qquad MPS = \frac{\Delta S}{\Delta Y} = 1 - MPC. \]
\item Always verify the two identities: $APC + APS = 1$ and $MPC + MPS = 1$.
\item Present answers as decimals (e.g. $0.8$) or percentages ($80\%$), as required by the question.
\end{enumerate}
\textbf{Reflex:} average propensities use \emph{levels}; marginal propensities use \emph{changes}. Never mix the two.
\end{methode}

\begin{exoresolu}[Propensities for a household in Bamenda]
The table below shows the disposable income and consumption of a household in Bamenda over three years (figures in thousands of FCFA).

\begin{center}
\begin{tabular}{|c|c|c|}
\hline
\rowcolor{popblueL} Year & Income $Y$ & Consumption $C$ \\
\hline
2021 & 2000 & 1900 \\
\hline
2022 & 2500 & 2300 \\
\hline
2023 & 3000 & 2700 \\
\hline
\end{tabular}
\end{center}

\begin{enumerate}
\item Calculate APC and APS in each year.
\item Calculate MPC and MPS between 2021 and 2022, and between 2022 and 2023.
\item Comment on the results.
\end{enumerate}
\tcblower
\textbf{Step 1: compute saving.} $S = Y - C$:
2021: $S = 2000 - 1900 = 100$; 2022: $S = 2500 - 2300 = 200$; 2023: $S = 3000 - 2700 = 300$ (thousand FCFA).

\textbf{Step 2: average propensities.}
\begin{align*}
APC_{2021} &= \frac{1900}{2000} = 0.95, & APS_{2021} &= \frac{100}{2000} = 0.05;\\
APC_{2022} &= \frac{2300}{2500} = 0.92, & APS_{2022} &= \frac{200}{2500} = 0.08;\\
APC_{2023} &= \frac{2700}{3000} = 0.90, & APS_{2023} &= \frac{300}{3000} = 0.10.
\end{align*}
Check: $0.95 + 0.05 = 0.92 + 0.08 = 0.90 + 0.10 = 1$. \checkmark

\textbf{Step 3: marginal propensities.}
Between 2021 and 2022: $\Delta Y = 500$, $\Delta C = 400$, $\Delta S = 100$:
\[ MPC = \frac{400}{500} = 0.8, \qquad MPS = \frac{100}{500} = 0.2. \]
Between 2022 and 2023: $\Delta Y = 500$, $\Delta C = 400$, $\Delta S = 100$:
\[ MPC = \frac{400}{500} = 0.8, \qquad MPS = \frac{100}{500} = 0.2. \]
Check: $0.8 + 0.2 = 1$. \checkmark

\textbf{Step 4: comment.} The APC falls as income rises (from $0.95$ to $0.90$), illustrating Keynes's fundamental psychological law: as income increases, consumption increases but by less than the increase in income. The MPC is constant at $0.8$, which means the consumption function here is linear: each extra 1000 FCFA of income generates 800 FCFA of extra consumption and 200 FCFA of extra saving.
\end{exoresolu}

\courssub{Skill 2: Deriving and using the consumption and saving functions}

\begin{methode}[Working with $C = a + bY$ and $S = -a + (1-b)Y$]
\begin{enumerate}
\item Identify \cle{autonomous consumption} $a$ (consumption when $Y = 0$, financed by past savings or borrowing) and the \cle{MPC} $b$.
\item The saving function follows immediately: $S = Y - C = -a + (1-b)Y$, where $1 - b = MPS$.
\item To find the \emph{break-even income} (where $C = Y$, so $S = 0$), solve $a + bY = Y$, giving $Y = \dfrac{a}{1-b}$.
\item To find consumption or saving at a given income, substitute the value of $Y$.
\item If given two observations $(Y_1, C_1)$ and $(Y_2, C_2)$, first compute $b = \dfrac{\Delta C}{\Delta Y}$, then $a = C_1 - bY_1$.
\end{enumerate}
\textbf{Reflex:} below break-even income, saving is negative (\emph{dissaving}); above it, saving is positive.
\end{methode}

\begin{exoresolu}[Consumption function of the Cameroonian economy]
Suppose the aggregate consumption function of an economy is $C = 200 + 0.75Y$, where $C$ and $Y$ are in billions of FCFA.
\begin{enumerate}
\item State the autonomous consumption and the MPC.
\item Derive the saving function.
\item Calculate the break-even level of income.
\item Calculate consumption and saving when national income is 2000 billion FCFA.
\end{enumerate}
\tcblower
\textbf{1.} Autonomous consumption $a = 200$ billion FCFA (consumption even when income is zero); $MPC = b = 0.75$.

\textbf{2.} Saving function:
\[ S = Y - C = Y - (200 + 0.75Y) = -200 + 0.25Y. \]
The MPS is $0.25 = 1 - 0.75$.

\textbf{3.} Break-even income: set $S = 0$ (or $C = Y$):
\[ -200 + 0.25Y = 0 \implies Y = \frac{200}{0.25} = 800 \text{ billion FCFA.} \]
At $Y = 800$, households consume exactly their income.

\textbf{4.} At $Y = 2000$:
\[ C = 200 + 0.75(2000) = 200 + 1500 = 1700 \text{ billion FCFA,} \]
\[ S = -200 + 0.25(2000) = 300 \text{ billion FCFA.} \]
Check: $C + S = 1700 + 300 = 2000 = Y$. \checkmark\ Also $APC = 1700/2000 = 0.85$, which is greater than the MPC ($0.75$): the APC always exceeds the MPC when autonomous consumption is positive, and the gap narrows as income rises.
\end{exoresolu}

\courssub{Skill 3: Investment appraisal — payback period and average rate of return}

\begin{methode}[Simple appraisal methods]
\begin{enumerate}
\item \textbf{Payback period:} cumulate the annual net cash inflows until the initial cost is recovered. If recovery occurs during a year, interpolate: extra months $= \dfrac{\text{amount still needed}}{\text{cash flow of that year}} \times 12$.
\item \textbf{Average rate of return (ARR):}
\[ ARR = \frac{\text{average annual profit (or net cash flow)}}{\text{initial cost of investment}} \times 100. \]
\item \textbf{Decision rules:} choose the project with the \emph{shortest} payback period or the \emph{highest} ARR; reject projects whose payback exceeds the firm's maximum acceptable period.
\item \textbf{Limits to mention in evaluation:} both methods ignore the time value of money; payback also ignores cash flows arising after the payback date.
\end{enumerate}
\end{methode}

\begin{exoresolu}[Choosing between two agro-processing machines]
A firm in Douala hesitates between two machines for processing cassava. Each costs 12 000 000 FCFA. Expected annual net cash inflows (FCFA) are:

\begin{center}
\begin{tabular}{|c|c|c|}
\hline
\rowcolor{popblueL} Year & Machine A & Machine B \\
\hline
1 & 4 000 000 & 6 000 000 \\
\hline
2 & 4 000 000 & 4 000 000 \\
\hline
3 & 4 000 000 & 3 000 000 \\
\hline
4 & 4 000 000 & 2 000 000 \\
\hline
5 & 4 000 000 & 1 000 000 \\
\hline
\end{tabular}
\end{center}

\begin{enumerate}
\item Compute the payback period of each machine.
\item Compute the ARR of each machine.
\item Advise the firm, stating one limit of each method.
\end{enumerate}
\tcblower
\textbf{1. Payback periods.}
Machine A: cumulative inflows: 4M; 8M; 12M. The cost (12M) is recovered exactly at the end of year 3: payback $= 3$ years.

Machine B: cumulative inflows: 6M; 10M; 13M. After 2 years, 10M is recovered; 2M is still needed out of year 3's 3M:
\[ \text{extra time} = \frac{2}{3} \times 12 = 8 \text{ months}. \]
Payback of B $= 2$ years 8 months. \textbf{B pays back faster.}

\textbf{2. ARR.}
Machine A: total inflow $= 5 \times 4 = 20$M; average $= 4$M per year:
\[ ARR_A = \frac{4}{12} \times 100 = 33.3\%. \]
Machine B: total inflow $= 6+4+3+2+1 = 16$M; average $= 16/5 = 3.2$M:
\[ ARR_B = \frac{3.2}{12} \times 100 = 26.7\%. \]
\textbf{A has the higher ARR.}

\textbf{3. Advice.} The two methods conflict: B is better on payback (liquidity, less risk), A is better on ARR (overall profitability). If the firm faces cash constraints, it may prefer B; if it seeks total return, A. \emph{Limits:} payback ignores the 8M that A earns after year 3; ARR ignores the timing of receipts — a franc received in year 1 is worth more than one received in year 5. This is why discounted methods such as NPV are superior.
\end{exoresolu}

\courssub{Skill 4: Calculating the Net Present Value (NPV)}

\begin{methode}[Computing and interpreting NPV]
\begin{enumerate}
\item Write the discount factor for each year $t$: $\dfrac{1}{(1+r)^t}$, where $r$ is the rate of interest (cost of capital) as a decimal.
\item Discount each net cash flow: $PV_t = \dfrac{R_t}{(1+r)^t}$.
\item Sum the discounted inflows and subtract the initial cost $C_0$:
\[ NPV = -C_0 + \sum_{t=1}^{n} \frac{R_t}{(1+r)^t}. \]
\item \textbf{Decision rule:} accept if $NPV > 0$ (the project earns more than the cost of capital); reject if $NPV < 0$; between mutually exclusive projects, choose the highest NPV.
\item Keep at least 2--3 decimal places in discount factors to avoid rounding errors.
\end{enumerate}
\textbf{Reflex:} the higher the discount rate, the lower the NPV — distant receipts shrink most.
\end{methode}

\begin{exoresolu}[NPV of a cocoa-drying investment]
A cooperative in the Centre Region considers buying a cocoa dryer costing 10 000 000 FCFA. It expects net cash inflows of 3 000 000 FCFA per year for 4 years. The cost of capital is $10\%$.
\begin{enumerate}
\item Calculate the NPV of the project.
\item Should the cooperative invest? Justify.
\item State one advantage of NPV over the payback method.
\end{enumerate}
\tcblower
\textbf{1.} Discount factors at $r = 0.10$: $\frac{1}{1.1} = 0.9091$; $\frac{1}{1.1^2} = 0.8264$; $\frac{1}{1.1^3} = 0.7513$; $\frac{1}{1.1^4} = 0.6830$.

\begin{center}
\begin{tabular}{|c|c|c|c|}
\hline
\rowcolor{popblueL} Year & Cash flow (FCFA) & Discount factor & Present value (FCFA) \\
\hline
1 & 3 000 000 & 0.9091 & 2 727 300 \\
\hline
2 & 3 000 000 & 0.8264 & 2 479 200 \\
\hline
3 & 3 000 000 & 0.7513 & 2 253 900 \\
\hline
4 & 3 000 000 & 0.6830 & 2 049 000 \\
\hline
\multicolumn{3}{|r|}{Total PV of inflows} & 9 509 400 \\
\hline
\end{tabular}
\end{center}

\[ NPV = 9 509 400 - 10 000 000 = -490 600 \text{ FCFA}. \]

\textbf{2.} $NPV < 0$: the project should be \textbf{rejected}. Discounted at $10\%$, the inflows are worth only about 9.51 million FCFA, less than the 10 million cost. Equivalently, the project yields less than the $10\%$ the cooperative could earn elsewhere. (Note: undiscounted inflows total 12 million — this shows why discounting matters.)

\textbf{3.} Advantage: NPV takes account of the \emph{time value of money} and uses \emph{all} cash flows over the project's life, whereas payback ignores both.
\end{exoresolu}

\courssub{Skill 5: Determining the equilibrium level of national income}

\begin{methode}[Equilibrium: aggregate demand and injections--withdrawals approaches]
\begin{enumerate}
\item In a closed economy without government, aggregate demand is $AD = C + I$, with $C = a + bY$.
\item \textbf{AD approach:} set $Y = AD$: solve $Y = a + bY + I$, giving
\[ Y_e = \frac{a + I}{1 - b}. \]
\item \textbf{Injections--withdrawals approach:} equilibrium requires planned injections $=$ planned withdrawals: $I = S$, with $S = -a + (1-b)Y$. Solve $I = -a + (1-b)Y$ — the same $Y_e$ results.
\item To find equilibrium with a government and foreign trade, add $G$ and $(X - M)$ to injections and $T$ and $M$ to withdrawals: equilibrium where $I + G + X = S + T + M$.
\item Always \emph{verify} by computing $C + I$ at $Y_e$ and checking it equals $Y_e$.
\end{enumerate}
\textbf{Reflex:} equilibrium income is where planned spending equals output — not necessarily full-employment income.
\end{methode}

\begin{exoresolu}[Equilibrium income in a simple economy]
In a hypothetical economy, $C = 100 + 0.8Y$ and planned investment $I = 150$ (billions of FCFA).
\begin{enumerate}
\item Find the equilibrium level of income using the aggregate demand approach.
\item Confirm the result using the injections--withdrawals approach.
\item If full-employment income is 1500 billion FCFA, comment on the situation.
\end{enumerate}
\tcblower
\textbf{1. Aggregate demand approach.}
\[ Y = C + I = 100 + 0.8Y + 150 \implies Y - 0.8Y = 250 \implies 0.2Y = 250 \]
\[ Y_e = \frac{250}{0.2} = 1250 \text{ billion FCFA.} \]
Verification: $C = 100 + 0.8(1250) = 1100$; $C + I = 1100 + 150 = 1250 = Y$. \checkmark

\textbf{2. Injections--withdrawals approach.} Saving function: $S = -100 + 0.2Y$. Equilibrium: $S = I$:
\[ -100 + 0.2Y = 150 \implies 0.2Y = 250 \implies Y_e = 1250 \text{ billion FCFA.} \]
Same result, as expected: at equilibrium, the withdrawal of saving (250) exactly matches the injection of investment... precisely, $S = 150 = I$. \checkmark

\textbf{3.} Full-employment income ($1500$) exceeds equilibrium income ($1250$) by $250$ billion FCFA. The economy is in \emph{under-employment equilibrium}: resources are idle. To reach full employment, aggregate demand must rise. The required increase in autonomous spending is not 250 but $\dfrac{250}{k}$ where $k = \frac{1}{1-0.8} = 5$, i.e. only $50$ billion FCFA of extra investment or government spending — this is the \emph{deflationary gap}. The gap between $AD$ and full-employment output is smaller than the output gap because of the multiplier.
\end{exoresolu}

\courssub{Skill 6: Calculating and applying the multiplier}

\begin{methode}[The investment multiplier]
\begin{enumerate}
\item Identify the MPC ($b$) or the MPS.
\item Compute the multiplier:
\[ k = \frac{1}{1 - MPC} = \frac{1}{MPS}. \]
In an economy with taxation ($t$) and imports ($m$), use $k = \dfrac{1}{MPS + MRT + MPM} = \dfrac{1}{1 - MPC_{\text{domestic}}}$.
\item Apply it: $\Delta Y = k \times \Delta I$ (or $\Delta G$, or any change in autonomous spending).
\item Reverse problems: given a target $\Delta Y$ (e.g. to close a deflationary gap), required spending $= \dfrac{\Delta Y}{k}$.
\item Remember the logic: one person's spending is another's income; spending is re-spent in successive rounds, with a fraction $b$ passed on each time, so $k = 1 + b + b^2 + \dots = \frac{1}{1-b}$.
\end{enumerate}
\textbf{Reflex:} the larger the leakages (saving, taxes, imports), the smaller the multiplier. In an open economy like Cameroon's, high import propensity reduces $k$.
\end{methode}

\begin{exoresolu}[Effect of public works spending]
In an economy, $MPC = 0.75$. The government increases its spending on road construction by 40 billion FCFA.
\begin{enumerate}
\item Calculate the value of the multiplier.
\item Calculate the total change in national income.
\item Show the first three rounds of spending that produce this result.
\item If the full-employment ceiling is only 120 billion FCFA above the current income, what problem arises?
\end{enumerate}
\tcblower
\textbf{1.} $k = \dfrac{1}{1 - 0.75} = \dfrac{1}{0.25} = 4$.

\textbf{2.} $\Delta Y = k \times \Delta G = 4 \times 40 = 160$ billion FCFA.

\textbf{3.} Round-by-round (billions of FCFA):
\begin{itemize}
\item Round 1: government spends $40$ → income of contractors and workers rises by $40$.
\item Round 2: recipients spend $0.75 \times 40 = 30$ → new income of $30$.
\item Round 3: $0.75 \times 30 = 22.5$ → new income of $22.5$.
\end{itemize}
Continuing: $40 + 30 + 22.5 + 16.875 + \dots$ is a geometric series with ratio $0.75$:
\[ \Delta Y = \frac{40}{1 - 0.75} = 160 \text{ billion FCFA.} \checkmark \]

\textbf{4.} If income can only rise by 120 billion before full employment, the extra demand of 160 billion exceeds available output by 40 billion FCFA: beyond full employment, additional spending bids up \emph{prices} rather than output. This excess is the \emph{inflationary gap}: demand-pull inflation results. The government should inject only $\dfrac{120}{4} = 30$ billion FCFA to reach full employment without inflation.
\end{exoresolu}

\courssub{Skill 7: The acceleration principle}

\begin{methode}[Accelerator calculations]
\begin{enumerate}
\item The \cle{accelerator} coefficient is $v = \dfrac{K}{Y}$ (capital--output ratio) or is given directly: $I_t = v\,\Delta Y = v\,(Y_t - Y_{t-1})$.
\item Compute the change in output (or demand) between the two periods.
\item Multiply by $v$ to obtain \emph{induced net investment}.
\item If replacement investment is given, add it to get \emph{total} investment.
\item Key reflexes: investment depends on the \emph{change} in income, not its level; if income grows at a constant rate, induced investment is constant; if income merely grows more slowly, induced investment \emph{falls} — this explains investment's volatility.
\end{enumerate}
\end{methode}

\begin{exoresolu}[Accelerator in a bottling firm]
A bottling firm in Yaoundé needs 2 FCFA of capital equipment to produce 1 FCFA of output per year ($v = 2$). Demand for its drinks evolves as follows (millions of FCFA): Year 1: 500; Year 2: 600; Year 3: 700; Year 4: 750; Year 5: 750. Replacement investment is 50 million FCFA each year.
\begin{enumerate}
\item Calculate induced (net) investment each year.
\item Calculate total investment each year.
\item Comment on the pattern observed.
\end{enumerate}
\tcblower
\textbf{1. and 2.} Apply $I_{net} = v\,\Delta Y$ with $v = 2$, then add replacement investment of 50.

\begin{center}
\begin{tabular}{|c|c|c|c|c|}
\hline
\rowcolor{popblueL} Year & Output $Y$ & $\Delta Y$ & Net investment $2\Delta Y$ & Total investment \\
\hline
1 & 500 & -- & -- & 50 \\
\hline
2 & 600 & 100 & 200 & 250 \\
\hline
3 & 700 & 100 & 200 & 250 \\
\hline
4 & 750 & 50 & 100 & 150 \\
\hline
5 & 750 & 0 & 0 & 50 \\
\hline
\end{tabular}
\end{center}

(figures in millions of FCFA)

\textbf{3. Comment.} Output rises throughout years 1--4 and never falls, yet total investment collapses from 250 to 50. Two lessons follow. First, investment depends on the \emph{rate of change} of demand: when demand grows at a constant rate (years 2--3), net investment is constant; when demand grows more slowly (year 4, $+50$ instead of $+100$), net investment falls by half even though demand is still rising. Second, when demand stops growing (year 5), net investment drops to zero and only replacement investment remains. This magnified response of investment to changes in output is the acceleration principle; combined with the multiplier (investment → income → consumption → investment), it generates cumulative booms and slumps, i.e. the trade cycle.
\end{exoresolu}

\begin{attention}[Frequent errors in the examination]
\begin{itemize}
\item Confusing $APC = C/Y$ with $MPC = \Delta C / \Delta Y$ — read the question carefully.
\item Forgetting that $k = 1/MPS$, not $1/MPC$.
\item Computing NPV without subtracting the initial cost, or discounting the initial cost itself (it occurs at year 0).
\item In accelerator questions, using the \emph{level} of output instead of its \emph{change}, and forgetting to add replacement investment when total investment is requested.
\item Claiming the deflationary gap equals the output gap: the gap in \emph{spending} is $\Delta Y / k$, much smaller than the gap in output.
\end{itemize}
\end{attention}

\newpage

\coursec{Exercises}

\courssub{I check my knowledge}

\begin{exercice}[Multiple choice questions]
For each question, choose the single correct answer.
\begin{enumerate}
\item The marginal propensity to consume (MPC) is defined as:
\begin{enumerate}
\item total consumption divided by total income;
\item the change in consumption divided by the change in income;
\item the change in income divided by the change in consumption;
\item total saving divided by total income.
\end{enumerate}
\item In a closed economy without government, equilibrium national income occurs when:
\begin{enumerate}
\item consumption equals investment;
\item saving equals investment;
\item saving equals consumption;
\item exports equal imports.
\end{enumerate}
\item If MPC $= 0.75$, the value of the simple multiplier is:
\begin{enumerate}
\item $1.33$;
\item $3$;
\item $4$;
\item $0.25$.
\end{enumerate}
\item The accelerator principle states that net investment depends on:
\begin{enumerate}
\item the rate of interest only;
\item the level of national income;
\item the \emph{rate of change} of output (income);
\item the level of government spending.
\end{enumerate}
\end{enumerate}
\end{exercice}

\begin{exercice}[True or false — justify your answer]
State whether each of the following statements is \textbf{true} or \textbf{false}, and justify your answer in two or three sentences.
\begin{enumerate}
\item When income is zero, consumption must also be zero.
\item APC $+$ APS $= 1$ and MPC $+$ MPS $= 1$ at every level of income.
\item In the circular flow of an open economy with government, taxes, saving and imports are injections, while government spending, investment and exports are leakages (withdrawals).
\item A positive net present value (NPV) means an investment project should be rejected.
\end{enumerate}
\end{exercice}

\begin{exercice}[Definitions]
Define each of the following terms precisely, as required at the GCE Advanced Level:
\begin{enumerate}
\item the circular flow of income;
\item autonomous consumption;
\item the marginal propensity to save (MPS);
\item induced investment;
\item the multiplier;
\item the deflationary gap.
\end{enumerate}
\end{exercice}

\begin{exercice}[Direct calculations]
\begin{enumerate}
\item A household's income rises from $200$ million FCFA to $260$ million FCFA, and its consumption rises from $170$ million FCFA to $218$ million FCFA. Calculate the MPC and the MPS.
\item Given the consumption function $C = 40 + 0.8Y$, write down the corresponding saving function.
\item Given MPC $= 0.6$, calculate the value of the simple multiplier.
\item A machine costs $5{,}000{,}000$ FCFA and yields a net cash flow of $2{,}000{,}000$ FCFA per year. Calculate its payback period.
\end{enumerate}
\end{exercice}

\courssub{I practise}

\begin{exercice}[Consumption and saving schedule]
The following schedule shows national income $Y$ and consumption $C$ (in million FCFA):
\begin{center}
\begin{tabular}{|c|c|c|c|c|c|}
\hline
$Y$ & $0$ & $100$ & $200$ & $300$ & $400$ \\
\hline
$C$ & $60$ & $140$ & $220$ & $300$ & $380$ \\
\hline
\end{tabular}
\end{center}
\begin{enumerate}
\item Calculate APC, MPC, APS and MPS at each level of income.
\item Derive the consumption function and the saving function.
\item Verify the two identities APC $+$ APS $= 1$ and MPC $+$ MPS $= 1$.
\item At what level of income is saving equal to zero? Comment.
\end{enumerate}
\end{exercice}

\begin{exercice}[The circular flow in an open economy]
\begin{enumerate}
\item Draw a clearly labelled diagram of the circular flow of income for an \textbf{open economy with a government sector}, showing households, firms, the government, the financial market and the foreign sector.
\item Identify on your diagram three leakages (withdrawals) and three injections.
\item State the equilibrium condition of national income in this economy in terms of leakages and injections.
\item Explain briefly what happens to the level of national income if injections exceed leakages.
\end{enumerate}
\end{exercice}

\begin{exercice}[Investment appraisal: NPV and payback]
A firm in Douala considers buying a machine costing $10{,}000{,}000$ FCFA, which is expected to yield net cash flows of $3{,}000{,}000$; $4{,}000{,}000$; $4{,}000{,}000$ and $2{,}000{,}000$ FCFA over four years. The discount rate is $12\%$.
\begin{enumerate}
\item Compute the net present value (NPV) of the project and advise the firm.
\item Calculate the payback period of the project.
\item State one advantage and one limitation of the payback method compared with the NPV method.
\end{enumerate}
\end{exercice}

\begin{exercice}[Multiplier and accelerator]
\begin{enumerate}
\item Given MPC $= 0.8$, calculate the multiplier and the total change in national income resulting from a new government road investment of $5$ billion FCFA in the Far North Region. Show the process round by round for the first three rounds.
\item The capital-output ratio of an economy is $3$ and output rises from $200$ to $260$ billion FCFA. Calculate induced net investment. What happens to net investment if output then stays constant at $260$ billion FCFA? Comment on the implication for economic stability.
\end{enumerate}
\end{exercice}

\courssub{I prepare for the GCE}

\begin{exercice}[Equilibrium income and the deflationary gap (structured question)]
In an economy, $C = 50 + 0.75Y$ and $I = 150$ (all values in million FCFA).
\begin{enumerate}
\item Define the equilibrium level of national income. \hfill (2 marks)
\item Determine the equilibrium level of income using the aggregate demand approach. \hfill (4 marks)
\item Determine the equilibrium level of income using the saving--investment (leakages--injections) approach, and show that both approaches give the same result. \hfill (4 marks)
\item If the full-employment level of income is $1{,}000$ million FCFA, calculate the deflationary gap. \hfill (3 marks)
\item Illustrate the situation with a clearly labelled Keynesian cross diagram. \hfill (4 marks)
\item Suggest one fiscal policy measure the government could use to close this gap, and calculate the change in government spending required. \hfill (3 marks)
\end{enumerate}
\end{exercice}

\begin{exercice}[Inflationary and deflationary gaps (structured question)]
\begin{enumerate}
\item Explain, with the aid of diagrams, the difference between an inflationary gap and a deflationary gap. \hfill (8 marks)
\item Using the aggregate demand and aggregate supply (or Keynesian cross) framework, explain why a deflationary gap is associated with unemployment. \hfill (4 marks)
\item Suggest and explain two fiscal policy measures the Cameroonian government could use to close a deflationary gap, indicating one possible limitation of each measure in the Cameroonian context. \hfill (8 marks)
\end{enumerate}
\end{exercice}

\begin{exercice}[Essay — GCE Paper 2 style]
\begin{enumerate}
\item ``An increase in the desire to save can make a nation poorer.'' Discuss this statement with reference to the circular flow of income and the paradox of thrift. \hfill (20 marks)
\end{enumerate}
\emph{Guidance:} your essay should define saving and the circular flow; explain the paradox of thrift using the saving--investment equilibrium and the multiplier; consider the conditions under which higher saving may instead promote growth (e.g.\ financing of investment); and reach a balanced, well-argued conclusion illustrated, where relevant, with a diagram.
\end{exercice}

\newpage

\coursec{Solutions to the exercises}

\begin{corrige}[Exercise 1 — Multiple choice questions]
\begin{enumerate}
\item \textbf{(b)}: by definition, MPC $= \dfrac{\Delta C}{\Delta Y}$, the fraction of an additional unit of income that is consumed. Answer (a) defines the \emph{average} propensity to consume, and (c) inverts the ratio.
\item \textbf{(b)}: in a closed economy without government, the only leakage is saving and the only injection is investment; equilibrium requires $S = I$. Answer (d) concerns an open economy, and (a) and (c) have no equilibrium meaning.
\item \textbf{(c)}: $k = \dfrac{1}{1 - \text{MPC}} = \dfrac{1}{1 - 0.75} = \dfrac{1}{0.25} = 4$.
\item \textbf{(c)}: the accelerator states that induced net investment is proportional to the \emph{change} in output: $I_n = v\,\Delta Y$, where $v$ is the capital--output ratio. A constant level of output implies zero induced net investment.
\end{enumerate}
\end{corrige}

\begin{corrige}[Exercise 2 — True or false]
\begin{enumerate}
\item \textbf{False.} Even when income is zero, households must consume a minimum amount to survive. This \cle{autonomous consumption} $a$ is financed by drawing down past savings or by borrowing (dissaving). Hence $C = a > 0$ when $Y = 0$, and the consumption curve cuts the vertical axis above the origin.
\item \textbf{True.} Disposable income is either consumed or saved: $Y = C + S$. Dividing through by $Y$ gives $\dfrac{C}{Y} + \dfrac{S}{Y} = 1$, i.e.\ APC $+$ APS $= 1$; taking changes gives $\Delta Y = \Delta C + \Delta S$, so dividing by $\Delta Y$ gives MPC $+$ MPS $= 1$.
\item \textbf{False.} The statement reverses the two concepts. Taxes, saving and imports are \cle{leakages} (withdrawals) because they remove spending power from the domestic circular flow; government spending, investment and exports are \cle{injections} because they add spending to the flow.
\item \textbf{False.} A positive NPV means the present value of the expected net cash flows exceeds the initial cost of the project; the project adds value and should be \emph{accepted}. A negative NPV signals rejection.
\end{enumerate}
\end{corrige}

\begin{corrige}[Exercise 3 — Definitions]
\begin{enumerate}
\item \textbf{Circular flow of income:} the continuous movement of income and expenditure between the main sectors of the economy (households, firms, government, financial market, foreign sector), in which households supply factor services to firms in exchange for factor incomes, and use those incomes to buy the goods and services produced by firms, so that one sector's expenditure becomes another sector's income.
\item \textbf{Autonomous consumption:} the part of consumption that does not depend on the current level of income; it is the minimum consumption undertaken even when income is zero (financed by dissaving or borrowing), denoted $a$ in the consumption function $C = a + bY$.
\item \textbf{Marginal propensity to save (MPS):} the fraction of an additional unit of income that is saved: $\text{MPS} = \dfrac{\Delta S}{\Delta Y} = 1 - \text{MPC}$.
\item \textbf{Induced investment:} investment that is caused by changes in the level (or rate of growth) of national income or output, as opposed to autonomous investment, which is independent of income. According to the accelerator, $I_n = v\,\Delta Y$.
\item \textbf{Multiplier:} the ratio of the total change in national income to the initial change in autonomous expenditure that caused it: $k = \dfrac{\Delta Y}{\Delta I} = \dfrac{1}{1 - \text{MPC}} = \dfrac{1}{\text{MPS}}$ in the simple closed-economy model.
\item \textbf{Deflationary gap:} the amount by which aggregate demand at the full-employment level of income falls short of the level needed to purchase the full-employment output; equivalently, the increase in autonomous expenditure required to raise equilibrium income to the full-employment level. It is associated with unemployment of resources.
\end{enumerate}
\end{corrige}

\begin{corrige}[Exercise 4 — Direct calculations]
\begin{enumerate}
\item MPC $= \dfrac{\Delta C}{\Delta Y} = \dfrac{218 - 170}{260 - 200} = \dfrac{48}{60} = 0.8$; hence MPS $= 1 - 0.8 = 0.2$.
\item $S = Y - C = Y - (40 + 0.8Y) = -40 + 0.2Y$. Autonomous saving is $-40$ (dissaving at zero income) and MPS $= 0.2$.
\item $k = \dfrac{1}{1 - \text{MPC}} = \dfrac{1}{1 - 0.6} = \dfrac{1}{0.4} = 2.5$.
\item Payback period $= \dfrac{\text{initial cost}}{\text{annual net cash flow}} = \dfrac{5{,}000{,}000}{2{,}000{,}000} = 2.5$ years, i.e.\ 2 years and 6 months.
\end{enumerate}
\end{corrige}

\begin{corrige}[Exercise 5 — Consumption and saving schedule]
\begin{enumerate}
\item Between any two successive rows, $\Delta C = 80$ and $\Delta Y = 100$, so MPC $= \dfrac{80}{100} = 0.8$ (constant) and MPS $= 1 - 0.8 = 0.2$ at all levels. The average propensities are:
\begin{center}
\begin{tabular}{|c|c|c|c|c|}
\hline
$Y$ & $S = Y - C$ & APC $= C/Y$ & APS $= S/Y$ & APC $+$ APS \\
\hline
$100$ & $-40$ & $1.4$ & $-0.4$ & $1$ \\
$200$ & $-20$ & $1.1$ & $-0.1$ & $1$ \\
$300$ & $0$ & $1.0$ & $0$ & $1$ \\
$400$ & $20$ & $0.95$ & $0.05$ & $1$ \\
\hline
\end{tabular}
\end{center}
APC and APS are not defined at $Y = 0$ (division by zero). Note that APC falls as income rises, approaching the MPC of $0.8$.
\item Autonomous consumption is the value of $C$ when $Y = 0$, so $a = 60$; with MPC $= 0.8$:
\[ C = 60 + 0.8Y \qquad\text{and}\qquad S = Y - C = -60 + 0.2Y. \]
\item From the last column of the table, APC $+$ APS $= 1$ at every level of income; and MPC $+$ MPS $= 0.8 + 0.2 = 1$. Both identities are verified.
\item $S = 0$ when $-60 + 0.2Y = 0$, i.e.\ $Y = 300$ million FCFA. This is the \emph{break-even} level of income, at which households consume exactly all their income. Below $Y = 300$ households dissave ($S < 0$, financed by past savings or borrowing); above it they make positive savings.
\end{enumerate}
\end{corrige}

\begin{corrige}[Exercise 6 — The circular flow in an open economy]
\begin{enumerate}
\item The diagram must contain five sectors and the following flows:
\begin{itemize}
\item \textbf{Households $\to$ firms:} factor services (labour, land, capital, enterprise); \textbf{firms $\to$ households:} factor incomes (wages, rent, interest, profit).
\item \textbf{Households $\to$ firms:} consumption expenditure $C$; \textbf{firms $\to$ households:} goods and services.
\item \textbf{Households $\to$ financial market:} saving $S$; \textbf{financial market $\to$ firms:} investment $I$.
\item \textbf{Households and firms $\to$ government:} taxes $T$; \textbf{government $\to$ households and firms:} government spending $G$ (public goods, transfers, salaries).
\item \textbf{Domestic economy $\to$ foreign sector:} payments for imports $M$; \textbf{foreign sector $\to$ domestic firms:} receipts from exports $X$.
\end{itemize}
\item Leakages (withdrawals): saving $S$, taxes $T$, imports $M$. Injections: investment $I$, government spending $G$, exports $X$.
\item Equilibrium condition: total leakages equal total injections:
\[ S + T + M = I + G + X. \]
\item If injections exceed leakages ($I + G + X > S + T + M$), more spending enters the circular flow than leaves it. Aggregate demand rises, firms observe falling stocks and rising sales, so they expand output and hire more factors of production. National income therefore \emph{rises}; as income grows, saving, tax revenue and imports also grow, until leakages have risen enough to restore equality with injections at a higher equilibrium level of income.
\end{enumerate}
\end{corrige}

\begin{corrige}[Exercise 7 — Investment appraisal: NPV and payback]
\begin{enumerate}
\item Discounting each net cash flow at $12\%$:
\begin{align*}
\text{PV} &= \frac{3{,}000{,}000}{1.12} + \frac{4{,}000{,}000}{1.12^{2}} + \frac{4{,}000{,}000}{1.12^{3}} + \frac{2{,}000{,}000}{1.12^{4}} \\
&\approx 2{,}678{,}571 + 3{,}188{,}776 + 2{,}847{,}121 + 1{,}271{,}036 \\
&\approx 9{,}985{,}504 \text{ FCFA}.
\end{align*}
\[ \text{NPV} = 9{,}985{,}504 - 10{,}000{,}000 \approx -14{,}500 \text{ FCFA}. \]
The NPV is (slightly) negative: at a $12\%$ discount rate the project does not quite recover its cost in present-value terms, so on purely financial grounds the firm should \emph{reject} it (or renegotiate the cost, or accept it only if the appropriate discount rate is slightly below $12\%$).
\item Cumulative cash flows: $3$ million after year 1; $7$ million after year 2; $11$ million after year 3. The cost of $10$ million is recovered during year 3:
\[ \text{Payback} = 2 + \frac{10 - 7}{4} = 2.75 \text{ years} \quad (\text{2 years and 9 months}). \]
\item \textbf{Advantage of payback:} it is simple to compute and to understand, and it is useful where liquidity and risk matter, because it favours projects that recover cash quickly. \textbf{Limitation:} it ignores the time value of money (no discounting) and disregards all cash flows arising after the payback date, so it can reject highly profitable long-lived projects. The NPV method, by contrast, discounts every cash flow over the whole life of the project and measures the net wealth created.
\end{enumerate}
\end{corrige}

\begin{corrige}[Exercise 8 — Multiplier and accelerator]
\begin{enumerate}
\item Multiplier: $k = \dfrac{1}{1 - \text{MPC}} = \dfrac{1}{1 - 0.8} = 5$. Total change in income:
\[ \Delta Y = k \times \Delta G = 5 \times 5 = 25 \text{ billion FCFA}. \]
Round-by-round process:
\begin{itemize}
\item \textbf{Round 1:} the government spends $5$ billion on the road; this becomes the income of engineers, workers and suppliers in the Far North Region ($+5$ billion).
\item \textbf{Round 2:} they consume $0.8 \times 5 = 4$ billion, which becomes income for traders, transporters, etc.\ ($+4$ billion).
\item \textbf{Round 3:} the recipients consume $0.8 \times 4 = 3.2$ billion ($+3.2$ billion).
\end{itemize}
The process continues with each round equal to $0.8$ times the previous one. The total is the sum of a geometric series:
\[ \Delta Y = 5 + 4 + 3.2 + \dots = \frac{5}{1 - 0.8} = 25 \text{ billion FCFA}. \]
\item Induced net investment:
\[ I_n = v\,\Delta Y = 3 \times (260 - 200) = 3 \times 60 = 180 \text{ billion FCFA}. \]
If output then stays constant at $260$ billion, $\Delta Y = 0$ and induced net investment falls to \emph{zero}, even though output remains at its highest level. \textbf{Comment:} the accelerator makes investment extremely volatile: investment depends not on the \emph{level} of demand but on its \emph{growth rate}. A mere slowdown in the growth of demand (from $+60$ to $0$) causes investment to collapse, which through the multiplier reduces income, which further reduces investment. This interaction is a major source of economic instability and helps explain trade cycles.
\end{enumerate}
\end{corrige}

\begin{corrige}[Exercise 9 — Equilibrium income and the deflationary gap]
\begin{enumerate}
\item \textbf{(2 marks)} The equilibrium level of national income is the level at which aggregate demand (planned expenditure) equals aggregate output, $Y = C + I$ — equivalently, at which planned leakages equal planned injections, $S = I$ — so that there is no tendency for income to change. \emph{Mark scheme: 1 mark for $AD = $ output or $S = I$; 1 mark for "no tendency to change".}
\item \textbf{(4 marks)} Aggregate demand approach:
\begin{align*}
Y &= C + I = 50 + 0.75Y + 150 \\
Y - 0.75Y &= 200 \\
0.25Y &= 200 \quad\Rightarrow\quad Y_e = 800 \text{ million FCFA}.
\end{align*}
\emph{Mark scheme: 1 mark for setting $Y = C + I$; 2 marks for correct algebra; 1 mark for the result with units.}
\item \textbf{(4 marks)} Saving function: $S = Y - C = Y - (50 + 0.75Y) = -50 + 0.25Y$. Setting leakages equal to injections:
\[ S = I \;\Rightarrow\; -50 + 0.25Y = 150 \;\Rightarrow\; 0.25Y = 200 \;\Rightarrow\; Y_e = 800 \text{ million FCFA}. \]
Both approaches give the same equilibrium, $Y_e = 800$ million FCFA, because the two conditions $Y = C + I$ and $S = I$ are algebraically equivalent (subtract $C$ from both sides of the first). \emph{Mark scheme: 1 mark for the saving function; 2 marks for solving; 1 mark for the equivalence comment.}
\item \textbf{(3 marks)} The required increase in income is $\Delta Y = Y_f - Y_e = 1{,}000 - 800 = 200$ million. The deflationary gap is the vertical shortfall of aggregate demand at $Y_f$:
\[ \text{Gap} = \Delta Y \times \text{MPS} = 200 \times 0.25 = 50 \text{ million FCFA}. \]
(Equivalently, $\Delta G = \Delta Y / k = 200/4 = 50$ million.) \emph{Mark scheme: 1 mark for $\Delta Y = 200$; 1 mark for the method; 1 mark for the answer of 50 million FCFA.}
\item \textbf{(4 marks)} The Keynesian cross diagram shows: income $Y$ on the horizontal axis, expenditure on the vertical axis; the $45^{\circ}$ line (points where expenditure equals income); the aggregate demand line $AD = 200 + 0.75Y$ with intercept $200$ and slope $0.75$; its intersection with the $45^{\circ}$ line at $Y_e = 800$; the full-employment income $Y_f = 1{,}000$ marked to the right of $Y_e$; and the vertical distance of $50$ million between the $AD$ line and the $45^{\circ}$ line at $Y_f$, labelled ``deflationary gap''. \emph{Mark scheme: 1 mark per correctly drawn and labelled element: axes and $45^{\circ}$ line; $AD$ line; $Y_e$ and $Y_f$; the gap.}
\item \textbf{(3 marks)} The government can use an \textbf{expansionary fiscal policy}, for example increasing its own expenditure on roads, schools or hospitals. With $k = \dfrac{1}{\text{MPS}} = \dfrac{1}{0.25} = 4$:
\[ \Delta G = \frac{\Delta Y}{k} = \frac{200}{4} = 50 \text{ million FCFA}. \]
An increase in government spending of $50$ million FCFA closes the gap. \emph{Mark scheme: 1 mark for the policy; 1 mark for $k = 4$; 1 mark for $\Delta G = 50$ million.}
\end{enumerate}
\end{corrige}

\begin{corrige}[Exercise 10 — Inflationary and deflationary gaps]
\begin{enumerate}
\item \textbf{(8 marks) Definitions and diagrams.} A \cle{deflationary gap} is the amount by which aggregate demand at the full-employment level of income falls \emph{short} of the output the economy can produce at full employment: equilibrium income $Y_e$ lies \emph{below} full-employment income $Y_f$, and resources (especially labour) are unemployed. On the Keynesian cross diagram, the $AD$ line cuts the $45^{\circ}$ line to the \emph{left} of $Y_f$; the gap is the vertical distance between the $45^{\circ}$ line and the $AD$ line measured at $Y_f$. An \cle{inflationary gap} is the opposite situation: aggregate demand at full employment \emph{exceeds} the full-employment output, so $Y_e$ (in money terms) lies to the \emph{right} of $Y_f$. Since real output cannot rise beyond $Y_f$, the excess demand bids up prices: the gap is purely inflationary. On the diagram, the $AD$ line cuts the $45^{\circ}$ line to the right of $Y_f$, and the gap is again the vertical distance at $Y_f$, this time with $AD$ above the $45^{\circ}$ line. \emph{Mark scheme: 2 marks per correct definition; 2 marks per correct labelled diagram (or one diagram showing both cases clearly).}
\item \textbf{(4 marks)} With a deflationary gap, aggregate demand is insufficient to buy the full-employment output. Firms cannot sell all they could produce, so they reduce output to the equilibrium level $Y_e < Y_f$ and lay off workers; capital equipment also lies idle. Because the demand for labour is derived from the demand for goods, deficient aggregate demand translates directly into \emph{demand-deficient (cyclical) unemployment}: workers are willing to work at the going wage but no jobs are offered. Only an increase in aggregate demand can restore full employment.
\item \textbf{(8 marks)} Two fiscal measures for Cameroon:
\begin{itemize}
\item \textbf{Increase government expenditure} (e.g.\ on roads, energy, schools, hospitals, or agricultural support in the regions). This injects spending directly into the circular flow and raises income by a multiple of the initial injection through the multiplier. \emph{Limitation:} the Cameroonian government's budget is constrained by limited tax revenue and by debt-servicing obligations; financing extra spending by borrowing may raise the public debt or crowd out private borrowing, and part of the spending leaks into imports, reducing the multiplier.
\item \textbf{Reduce taxation} (e.g.\ lower income tax or VAT), which raises households' disposable income and therefore consumption. \emph{Limitation:} with a high marginal propensity to import and a high marginal propensity to save among some income groups, much of the tax cut leaks out of the circular flow, weakening its effect; moreover tax cuts reduce government revenue, which is already narrow in Cameroon, risking larger budget deficits.
\end{itemize}
\emph{Mark scheme: 2 marks per measure correctly explained, 2 marks per relevant limitation, applied to the Cameroonian context.}
\end{enumerate}
\end{corrige}

\begin{corrige}[Exercise 11 — Essay: the paradox of thrift]
\textbf{Indicative mark scheme (20 marks):} definitions and framework (4); explanation of the paradox using the circular flow, the $S = I$ equilibrium and the multiplier (8); counter-arguments and conditions under which saving promotes growth (5); balanced conclusion (2); quality of argument, diagram and examples (1).

\textbf{Introduction.} Saving is the part of disposable income that is not consumed; in the circular flow of income it is a leakage, while investment is the corresponding injection. Intuitively, thrift appears virtuous: what is true for an individual, however, need not be true for the whole nation. The statement refers to Keynes's \cle{paradox of thrift}.

\textbf{Development 1 — the paradox.} In the simple Keynesian model, equilibrium income is determined by $S = I$, where saving rises with income ($S = -a + (1-b)Y$) while investment is autonomous. Suppose all households decide to save more at every level of income: the saving curve shifts upward. Consumption, a component of aggregate demand, falls; firms face unsold goods and cut output and employment; national income falls through the multiplier process ($\Delta Y = k \times \Delta C$, with $k = 1/\text{MPS}$). Income keeps falling until saving has fallen back to equality with the unchanged level of investment. The result: the nation ends up with a \emph{lower} income and \emph{no more} (possibly even less) total saving than before. An attempt by everyone to save more makes the nation poorer — this is the paradox of thrift, an example of the fallacy of composition. A diagram showing the upward shift of the saving curve against a horizontal investment line, with equilibrium income falling from $Y_1$ to $Y_2$, earns full credit.

\textbf{Development 2 — qualifications.} The paradox holds under restrictive assumptions: investment is autonomous, resources are unemployed, and the interest rate does not adjust. In the longer run, higher saving can \emph{finance} investment: through banks and financial markets (in Cameroon, for instance, via the CEMAC banking system and the Douala Stock Exchange), household savings can fund firms' capital goods, raising productive capacity and future growth. If investment is interest-elastic, a larger supply of loanable funds lowers interest rates and stimulates $I$. Moreover, in an economy already at full employment, extra saving releases resources for investment rather than reducing output. High saving rates have historically financed rapid capital accumulation in several Asian economies.

\textbf{Conclusion.} The statement is valid in the short run of a demand-deficient economy, where investment does not automatically absorb extra saving: a general rise in thrift then reduces aggregate demand and, through the multiplier, makes the nation poorer. But it is not a universal law: where financial intermediation channels saving into productive investment, and where the economy operates near full employment, greater thrift can raise investment, capacity and future income. The wise policy is therefore not to discourage saving but to ensure that savings are effectively mobilised and transformed into investment — a major challenge for financial development in Cameroon and the CEMAC region.
\end{corrige}

\newpage

\coursec{Summary sheet}

\begin{popfigure}
\begin{tikzpicture}[scale=0.9, every node/.style={font=\small}]
\tikzset{
  central/.style={circle, draw=popdark, fill=popblue, text=white, minimum width=3cm, text width=2.5cm, align=center, font=\bfseries\small},
  branch/.style={rectangle, rounded corners, draw=popdark, fill=poporange, text=white, minimum width=2.5cm, text width=2.2cm, align=center, font=\small},
  branchB/.style={branch, fill=popgreen},
  branchC/.style={branch, fill=poppurple},
  branchD/.style={branch, fill=poppink},
  branchE/.style={branch, fill=popgold},
  branchF/.style={branch, fill=popblueL},
}
\node[central, align=center] (center){Circular Flow\\ of Income};
\node[branch, align=center] (b1) at (90:4){Macroeconomic\\ Variables};
\node[branchB, align=center] (b2) at (30:4){Consumption\\ Function};
\node[branchC, align=center] (b3) at (-30:4){Saving\\ Function};
\node[branchD, align=center] (b4) at (-90:4){Investment\\ and Appraisal};
\node[branchE, align=center] (b5) at (-150:4){Equilibrium\\ and Multiplier};
\node[branchF, align=center] (b6) at (150:4){Accelerator\\ and Interaction};
\foreach \n in {b1,b2,b3,b4,b5,b6} \draw[popline, thick] (center) -- (\n);
\end{tikzpicture}
\legende{Mind map of the Circular Flow of Income chapter}
\end{popfigure}

\begin{bilan}
\textbf{Key Definitions}
\begin{itemize}
\item \cle{Circular flow of income}: Continuous flow of money, goods and services between households and firms; includes leakages (saving, taxes, imports) and injections (investment, government spending, exports).
\item \cle{Macroeconomic variables}: Aggregate output (Y), consumption (C), saving (S), investment (I), government expenditure (G), net exports (X-M), aggregate demand (AD), aggregate supply (AS).
\item \cle{Consumption function}: Relationship between consumption and disposable income, typically $C = a + bY_d$ where $a$ is autonomous consumption and $b$ is MPC.
\item \cle{Saving function}: Derived from consumption function: $S = -a + (1-b)Y_d$; $MPS = 1-MPC$.
\item \cle{Investment}: Expenditure on capital goods; autonomous or induced (accelerator). Appraisal methods: Payback period, ARR, NPV, IRR.
\item \cle{Equilibrium national income}: $Y = AD = C+I+G+(X-M)$ or Injections = Withdrawals ($I+G+X = S+T+M$).
\item \cle{Multiplier}: $k = \frac{1}{1-MPC} = \frac{1}{MPS}$; shows change in equilibrium income from a change in autonomous spending.
\item \cle{Accelerator}: $I_t = v\Delta Y_{t-1}$; induced investment depends on rate of change of output.
\item \cle{Inflationary/Deflationary gaps}: Difference between equilibrium income and full-employment income.
\end{itemize}

\textbf{Formulas and Laws to Know}
\begin{itemize}
\item $APC = \frac{C}{Y}$, $MPC = \frac{\Delta C}{\Delta Y}$; $APS = \frac{S}{Y}$, $MPS = \frac{\Delta S}{\Delta Y}$.
\item $MPC + MPS = 1$ (for a closed economy without taxes).
\item Multiplier $k = \frac{1}{1-MPC} = \frac{1}{MPS}$; with taxes and imports: $k = \frac{1}{1-MPC(1-t)+m}$.
\item Equilibrium condition: $Y = C + I + G + X - M$.
\item Injection-Withdrawal equality: $I + G + X = S + T + M$.
\item Net Present Value: $NPV = \sum_{t=1}^{n} \frac{CF_t}{(1+r)^t} - I_0$; accept if $NPV > 0$.
\item Accelerator coefficient $v = \frac{\Delta I}{\Delta Y}$.
\item Interaction: Multiplier-accelerator model can generate cyclical fluctuations (Samuelson-Hicks).
\end{itemize}

\textbf{Methods}
\begin{itemize}
\item \textbf{Calculating APC, MPC, APS, MPS}: Use a table of income and consumption/saving; compute averages and marginal changes; verify $MPC+MPS=1$.
\item \textbf{Investment appraisal}: Compute payback period, ARR, NPV (discount cash flows at cost of capital), IRR (rate where NPV=0). Compare mutually exclusive projects using NPV.
\item \textbf{Finding equilibrium income}: 
  \begin{enumerate}
  \item Write consumption function $C = a + bY_d$ (with $Y_d = Y - T$).
  \item Aggregate demand $AD = C + I + G + X - M$.
  \item Set $Y = AD$ and solve for $Y^*$.
  \item Alternatively, use $I+G+X = S+T+M$ and solve.
  \end{enumerate}
\item \textbf{Multiplier calculation}: Identify autonomous spending change $\Delta A$; compute $\Delta Y = k \times \Delta A$.
\item \textbf{Accelerator calculation}: Given $v$ and $\Delta Y$, find induced investment $\Delta I = v \Delta Y$.
\item \textbf{Gap analysis}: Compute full-employment income $Y_f$; compare with $Y^*$; inflationary gap if $Y^* > Y_f$, deflationary if $Y^* < Y_f$.
\end{itemize}

\textbf{Common Mistakes}
\begin{itemize}
\item Confusing average propensities (APC, APS) with marginal propensities (MPC, MPS).
\item Forgetting that $MPC+MPS=1$ only holds when there are no taxes or imports; in open economy with taxes, $MPC+MPS < 1$.
\item Using nominal values instead of real values when calculating propensities or multiplier.
\item Misapplying NPV: discounting at wrong rate, omitting initial outlay, or ignoring timing of cash flows.
\item In equilibrium, mixing up injections and withdrawals; e.g., treating imports as injection.
\item Assuming multiplier works instantly; ignoring time lags and capacity constraints.
\item In accelerator, using current $\Delta Y$ instead of lagged $\Delta Y$.
\item Not drawing labelled diagrams (circular flow, AD/AS, multiplier process) in essay questions.
\end{itemize}

\textbf{GCE Examination Tips}
\begin{itemize}
\item Always state the formula before substituting numbers; show all steps for full marks.
\item Practice numerical questions on APC/MPC/APS/MPS, multiplier, NPV, and accelerator; they are frequent in Paper 2.
\item For essay questions, structure: definition, diagram, explanation, evaluation (advantages/limitations), Cameroonian context (e.g., cocoa exports as injection, CFA franc stability, SONARA investment).
\item Know the difference between inflationary and deflationary gaps and appropriate fiscal/monetary policies.
\item Use the injection-withdrawal approach for open economy equilibrium; it is often quicker than AD=AS.
\item Remember that the multiplier is larger when MPC is high and leakages (taxes, imports) are low.
\item In accelerator questions, clarify whether investment is autonomous or induced; the accelerator only explains induced part.
\item Revise the interaction: multiplier-accelerator can produce cycles; understand the roles of $k$ and $v$ in Samuelson's model.
\item Time management: allocate about 1 minute per mark; for 20-mark essays, spend 20 minutes with clear paragraphs.
\end{itemize}
\end{bilan}

\findecours

\end{document}
